%%%%%%%%%%%%%%%%%%%%%%%%%%%%%%%%%%%%%%%%
% datoteka diploma_predloga.tex
%
% ta verzija sproducira pdf datoteko v pdf/A formatu.
%
% vzorčna datoteka za pisanje diplomskega dela v formatu LaTeX
% na UL Fakulteti za računalništvo in informatiko
%
% Leta 2026 je predlogo posodobil Žiga Emeršič. Predlogo je razdelil na več datotek, uvedel spremenljivke, dodal PDF-A in posodobil navodila (tudi združil določena poglavja in pripravil nova).
% Leta 2021 je predlogo posodobil Franc Solina.
% Leta 2010 je predlogo pripravil Gašper Fijavž.
%
% za upravljanje z literaturo ta vezija uporablja BibLaTeX
%
% svetujemo uporabo Overleaf.com - na tej spletni implementaciji LaTeXa ta vzorec zagotovo pravilno deluje
%

% Združljivost s pdfa (sicer nekoliko starejšo različico 1.7)
\DocumentMetadata{pdfstandard = a-2b, pdfversion  = 1.7}

\documentclass[a4paper,12pt,openright]{book}
%\documentclass[a4paper, 12pt, openright, draft]{book}  Nalogo preverite tudi z opcijo draft, ki pokaže, katere vrstice so predolge! Pozor, v draft opciji, se slike ne pokažejo!

% \flushbottom (privzeto za book) razteguje presledke, da se besedilo poravna do
% spodnjega roba strani - to na straneh s slikami povzroči velike umetne praznine.
% \raggedbottom to prepreči, strani se končajo naravno.
\raggedbottom


%%%%%%%%%%%%%%%%%%%%%%%%%%%%%%%%%%%%%%%%
%  PODATKI O DIPLOMI
%%%%%%%%%%%%%%%%%%%%%%%%%%%%%%%%%%%%%%%%

\newcommand{\tauthor}{Nik Katanović Leban}
\newcommand{\temail}{nk17425@student.uni-lj.si}
\newcommand{\tmentor}{viš. pred. dr. Igor Rožanc}

\newcommand{\tstopnja}%
{VISOKOŠOLSKI STROKOVNI ŠTUDIJSKI PROGRAM\\ PRVE STOPNJE}

\newcommand{\tinst}%
{Fakulteta za računalništvo in informatiko}

\newcommand{\tvrsta}{Diplomska naloga na visokošolskem programu prve stopnje Računalništvo in informatika}

% TODO: Ask Igor Rožanc what should i do about this.
% Opis iz STUDISa:
\newcommand{\tdescription}
{V diplomskem delu zasnujte prototip spletne aplikacije za prodajo vozil z vključeno samodejno oceno tržne vrednosti vozil na podlagi primerljivih oglasov na spletu. V ta namen preučite obstoječe rešitve, opredelite zahteve in zasnujte aplikacijo, pri čemer posebno pozornost namenite učinkovitemu vrednotenju vozil. Z uporabo ustreznih tehnologij implementirajte delujoč prototip ter ga ustrezno analizirajte in ovrednotite.}
\newcommand{\tdescriptionEn}
{In the thesis, design a prototype of a web application for vehicle sales with integrated automated market value estimation based on comparable online listings. For this purpose, examine existing solutions, specify the requirements, and design the application, with a special focus on effective vehicle valuation. Implement a working prototype using appropriate technologies, then analyze and evaluate it accordingly.}

% Naslov:
\newcommand{\ttitle}{Prototip spletne aplikacije za prodajo vozil z ocenjevanjem njihove vrednosti}
\newcommand{\ttitleEn}{Prototype of a web application for selling vehicles with vehicle valuation}

% Ključne besede:
\newcommand{\tkeywords}{spletna aplikacija, vozila, vrednotenje, React, Express.js, Node.js, PostgreSQL, Docker}
\newcommand{\tkeywordsEn}{web application, vehicles, valuation, React, Express.js, Node.js, PostgreSQL, Docker}

% Zahvala (opcijsko):
\newcommand{\tzahvala}
{Rad bi se zahvalil svojemu mentorju Igorju Rožancu, za vso aktivno pomoč in svetovanje pri razvoju te diplomske naloge. Zahvalil bi se tudi vsem sodelavcem iz podjetja CargoX, ki so z mano delili svoje znanje in izkušnje, ter mi omogočili, da sem se lahko lotil tako zahtevnega projekta. Posebna zahvala gre moji družini, ki me je ves čas podpirala in mi stala ob strani.}

%%%%%%%%%%%%%%%%%%%%%%%%%%%%%%%%%%%%%%%%
%  VKLJUČITEV NUJNIH ELEMENTOV 1/2
%%%%%%%%%%%%%%%%%%%%%%%%%%%%%%%%%%%%%%%%

\input{podporno/preambula.tex} 
\begin{document}

%%%%%%%%%%%%%%%%%%%%%%%%%%%%%%%%%%%%%%%%
% POVZETEK
%%%%%%%%%%%%%%%%%%%%%%%%%%%%%%%%%%%%%%%%

\newcommand{\tpovzetek}
{Diplomsko delo obravnava problem določanja ustrezne tržne vrednosti rabljenega vozila, pri prodaji ali nakupu vozila, zaradi nezadostnega pregleda stanja na trgu, stranke težko ocenijo primernost ponujene cene in vozilo kupijo po previsoki ceni ali ga prodajo pod njegovo tržno vrednostjo. Zato smo razvili prototip spletne aplikacije za objavo in pregled oglasov vozil, ki uporabnikom omogoča enostavno in varno prodajo ter nakup vozil. Ključen del predlagane rešitve je sistem za avtomatsko ocenjevanje vrednosti vozil, ki temelji na zbiranju podatkov iz zunanjih virov, njihovi normalizaciji in primerjavi vozila s podobnimi vozili na trgu. Pri oceni vrednosti se upoštevajo ključne lastnosti vozila, kot so znamka, model, letnik in prevoženi kilometri. Za pridobivanje podatkov so bili uporabljeni javno dostopni podatkovni viri ter postopki samodejnega zajema podatkov s spletnih strani za prodajo vozil. Zbrane podatke sistem obdela in uporabi za izračun ocenjene tržne vrednosti vozila. Rezultat diplomskega dela je delujoč prototip spletne aplikacije za prodajo vozil s podporo za samodejno vrednotenje vozil, razvit z uporabo tehnologij React~\cite{react-docs}, TypeScript~\cite{typescript}, Node.js~\cite{nodejs}, Express.js~\cite{expressjs-docs}, Python~\cite{python3.11}, PostgreSQL~\cite{postgresql-docs}, Redis~\cite{redis-data-store} in Docker~\cite{docker}.}

\newcommand{\tabstract}
{This thesis addresses the problem of determining the appropriate market value of a used vehicle. When selling or buying a vehicle, customers have an insufficient overview of the situation on the market, so they find it difficult to assess whether the offered price is reasonable and end up buying a vehicle at too high a price or selling it below its market value. We therefore developed a prototype of a web application for publishing and browsing vehicle listings, which enables users to sell and buy vehicles simply and safely. A key part of the proposed solution is a system for automatic vehicle valuation, based on collecting data from external sources, normalising it, and comparing a vehicle with similar vehicles on the market. The valuation takes into account key vehicle characteristics such as make, model, year, and mileage. Publicly available data sources and automated data collection from vehicle sales websites were used to obtain the data. The system processes the collected data and uses it to compute the estimated market value of a vehicle. The result of the thesis is a working prototype of a web application for selling vehicles with support for automatic vehicle valuation, developed using the technologies React~\cite{react-docs}, TypeScript~\cite{typescript}, Node.js~\cite{nodejs}, Express.js~\cite{expressjs-docs}, Python~\cite{python3.11}, PostgreSQL~\cite{postgresql-docs}, Redis~\cite{redis-data-store}, and Docker~\cite{docker}.}

%%%%%%%%%%%%%%%%%%%%%%%%%%%%%%%%%%%%%%%%
%  VKLJUČITEV NUJNIH ELEMENTOV 2/2
%%%%%%%%%%%%%%%%%%%%%%%%%%%%%%%%%%%%%%%%

\input{podporno/naslovnestrani.tex}
\mainmatter
\setcounter{page}{1}
\pagestyle{fancy}

%%%%%%%%%%%%%%%%%%%%%%%%%%%%%%%%%%%%%%%%
% GLAVNA VSEBINA
%%%%%%%%%%%%%%%%%%%%%%%%%%%%%%%%%%%%%%%%

\chapter{Uvod}
\label{ch:uvod}

Nakup in prodaja rabljenih vozil predstavljata pomemben del trga rabljenih vozil.
Pri tem se kupci in prodajalci pogosto srečujejo s težavo določanja primerne tržne
vrednosti vozila, saj nimajo zadostnega pregleda nad trenutno ponudbo in gibanjem cen na
trgu. Raznolikost informacij je še posebej izrazita,
ker je vsako vozilo unikatno glede na svoje stanje, opremo in zgodovino
uporabe~\cite{gegic2019car,eurotax}.\\ Najpogostejši pristop k oceni vrednosti je zato primerjava
 primerljivih prodaj podobnih vozil iz obstoječih oglasov, vendar je samostojno izvajanje takšne
primerjave časovno potratno in nezanesljivo. Oglasi lahko vsebujejo nerealno visoke ali nizke cene,
uporabniki pa pogosto nimajo dovolj znanja in izkušenj, zato lahko vozilo prodajo pod njegovo tržno
vrednostjo ali pa ga kupijo po previsoki ceni.\\
Na trgu že obstajajo rešitve za ocenjevanje vrednosti vozil, vendar so mnoge plačljive, 
prilagojene tujim trgom ali pa temeljijo na zastarelih podatkih, zaradi česar njihove ocene
pogosto ne odražajo dejanskega stanja na slovenskem trgu rabljenih vozil.\\
Cilj diplomskega dela je razvoj prototipa spletne aplikacije za objavo in pregled oglasov
vozil, ki vsebuje tudi sistem za samodejno ocenjevanje vrednosti vozil. Sistem temelji na
zbiranju podatkov iz zunanjih virov, njihovi normalizaciji in analizi primerljivih vozil na
trgu~\cite{laender2002}. Pri določanju ocenjene vrednosti se upoštevajo ključne lastnosti
vozila, kot so znamka, model, letnik prve registracije in število prevoženih kilometrov.
Dolgoročni cilj sistema je privabiti zadostno število uporabnikov, da na platformi
objavijo lastne oglase vozil, s čemer bo potreba po podatkih iz zunanjih virov
postopno upadla, tržno vrednost vozil pa bo mogoče ocenjevati na podlagi preverjenih oglasov, ki so zbrani
znotraj platforme same.\\
Na začetku diplomskega dela si zastavljamo izhodišče, da je mogoče s pomočjo zbranih tržnih podatkov in
primerjalne analize podobnih vozil uporabniku ponuditi uporabno in dovolj natančno oceno
tržne vrednosti vozila. Pri tem želimo preveriti, ali vključitev takšnega sistema v
prototip izboljša uporabniško izkušnjo pri nakupu in prodaji vozil.\\
Preostanek diplomskega dela je organiziran na naslednji način. Poglavje~\ref{ch:sorodne-resitve}
predstavi podobne rešitve za prodajo in vrednotenje vozil. Poglavje~\ref{ch:struktura} podrobneje
predstavi oglaševanje vozil, na katerem temelji razvita rešitev, in
postopke zajema podatkov s spletnih virov. Poglavje~\ref{ch:razvoj} je jedro naše rešitve in opiše
opredelitev zahtev, načrtovanje podatkovne baze in arhitekture sistema, implementacijo zalednega in
odjemalskega dela aplikacije, delovanje algoritma za vrednotenje vozil ter infrastrukturo in
avtomatizacijo namestitve. V poglavju~\ref{ch:analiza} je rešitev analizirana, diplomsko delo pa se zaključi s povzetkom doseženih rezultatov,
 kritično presojo omejitev rešitve in smeri za nadaljnje izboljšave.

\chapter{Sorodne rešitve}
\label{ch:sorodne-resitve}
Na področju prodaje rabljenih vozil obstaja več spletnih aplikacij, ki uporabnikom omogočajo objavo in pregled oglasov vozil.
 Pokaže pa se, da imajo le nekatere možnost vrednotenja vozil, od teh je večina je plačljivih. 
 Med najbolj znanimi tujimi ponudniki sta \textbf{mobile.de}~\cite{mobile-de} in \textbf{AutoScout24}~\cite{autoscout24},
 ki združujeta veliko število oglasov iz različnih držav ter uporabnikom omogočata filtriranje vozil po številnih kriterijih.
 Poleg njiju obstaja še vrsta drugih platform, ki uporabnikom bodisi zgolj omogočajo prodajo vozil bodisi ponujajo oceno
njihove vrednosti, med njimi so \textbf{AutoTrader}~\cite{autotrader}, \textbf{La Centrale}~\cite{lacentrale}, \textbf{Otomoto
}~\cite{otomoto}, \textbf{Cars.com}~\cite{cars-com} in \textbf{Kelley Blue Book}~\cite{kbb}. \\
Za slovenski trg je najbolj razširjena platforma \textbf{AvtoNet}~\cite{avto-net}, ki
predstavlja naše osrednje mesto za objavo in iskanje rabljenih vozil. Poleg aplikcij obstajajo tudi
specializirane storitve za vrednotenje vozil, najbolj znani sta \textbf{Eurotax}~\cite{eurotax} in \textbf{CARFAX}~\cite{carfax}.\\
 Platforma \textbf{AvtoNet}~\cite{avto-net} je
izšla leta 1997 in danes velja za prevladujočo aplikacijo za prodajo vozil na
slovenskem trgu. Uporabniku omogoča prodajo raznovrstnih vozil (kot so
avtomobili, motorna kolesa in plovila). Kljub dolgi prisotnosti na trgu ni
doživela večjih posodobitev, zato je njen uporabniški vmesnik zastarel in ponuja
precej prostora za izboljšave. Za objavo oglasa se mora uporabnik prijaviti,
medtem ko za brskanje po oglasih prijava ni potrebna. Prikaz posameznega oglasa na \textbf{AvtoNet-u}~\cite{avto-net}
je predstavljen na sliki~\ref{fig:avto-net-example}.
\begin{figure}[H]
  \centering
  \includegraphics[width=0.9\textwidth]{primer-avto-net.png}
  \caption{Primer oglasa na Avto.net-u}
  \label{fig:avto-net-example}
\end{figure}
Uporabniški vmesnik je pregleden in omogoča enostavno orientacijo, zato
  se uporabnik na njem hitro znajde. Glavna pomanjkljivost so predvsem manjkajoče
  dodatne funkcionalnosti, ki bi izboljšale uporabniško izkušnjo. Dober primer je obrazec za objavo oglasa.
  Vneseni podatki se recimo ob prehodu na
  naslednji korak ne shranijo, zato mora uporabnik ob vrnitvi na prejšnji korak, ponovno vpisati celotno vsebino.
   Obrazec je tudi razmeroma dolg in uporabniku ne prikazuje,
  koliko korakov ima še do zaključka objave. Platforma sicer dela zanesljivo
  in zadovoljuje osnovne potrebe uporabnikov. Uporabniški vmesnik za objavo oglasa je prikazan na
  sliki~\ref{fig:avto-net-objava}.
\begin{figure}[H]
  \centering
  \includegraphics[width=0.9\textwidth]{primer-objave-avto-net.png}
  \caption{Objava oglasa na Avto.net-u}
  \label{fig:avto-net-objava}
\end{figure}
 \textbf{Mobile.de}~\cite{mobile-de} ima novejši uporabniški vmesnik, ki je bolj pregleden in omogoča enostavno navigacijo po spletnem mestu.
 Uporabniku omogoča da objavi oglas, brska po oglasih in filtrira vozila po različnih kriterijih. Za razliko od \textbf{AvtoNet-a}~\cite{avto-net}, \textbf{Mobile.de}~\cite{mobile-de} omogoča tudi oceno vrednosti vozila.
  Pri posameznih oglasih prikazuje tudi mnenje, ali je cena ugodna, primerna ali previsoka. Mnenje je
  predstavljeno z petstopenjsko barvno lestvico: pri ugodni ceni je lestvica
  obarvana z zeleno, pri previsoki rdeče, ostale stopnje pa so označene z
  vmesnimi barvami. Takšen pristop uporabniku omogoča hitrejše odločanje, 
  saj mu ni treba ročno primerjati oglasov med seboj. Ocena temelji na
  samodejni primerjavi s podobnimi vozili v ponudbi. Pri tem pa platforma ne ponuja
  možnosti samostojnega vrednotenja vozila na podlagi uporabniško vnesenih
  podatkov, kadar oglas še ne obstaja. Prikaz osnovnih podatkov
  vozila z mnenjem je predstavljen na sliki~\ref{fig:mobile-de-example}.
\begin{figure}[H]
    \centering
    
\includegraphics[width=0.5\textwidth]{primer-mobile-de.png}
    \caption{Primer ocene vozila na Mobile.de}
    \label{fig:mobile-de-example}
\end{figure}
Specializirana storitev za vrednotenje vozil je \textbf{Eurotax}~\cite{eurotax}. Ki se pogosto uporablja v industriji vozil, zavarovalništvu in pri trgovcih z vozili.
 Njegove ocene temeljijo na obsežnih zbirkah podatkov o vozilih in tržnih cenah,
  vendar je dostop do storitve praviloma plačljiv. Zaradi poslovnega modela in
   omejenega dostopa podrobna analiza delovanja sistema ni mogoča.\\
Na severnoameriškem trgu je najbolj razširjena tovrstna storitev \textbf{CARFAX}~\cite{carfax}, ki
  poleg vrednotenja vozila ponuja tudi vpogled v njegovo celotno zgodovino.
  Vsako vozilo je določeno z enolično identifikacijsko številko (VIN)~\cite{iso3779}, na podlagi
  katere storitev pripravi poročilo o zgodovini vozila. To poročilo vključuje
  podatke o številu prejšnih lastnikov, servisni zgodovini, zgodovini registracij, prijavljenih prometnih nesrečah in škodi.
\textbf{CARFAX}~\cite{carfax} ponuja funkcijo vrednotenja vozila, ki za razliko od preprostih ocen, ki upoštevajo zgolj znamko, model
  in letnik, vključuje zgodovino vozila, prevožene kilometre, škodo zaradi nesreč,
   servisne posege ter uporabo vozila (osebna,
  najemniška ali komercialna raba). Vrednost se prilagaja tudi glede na regijo,
  saj se ponudba in povpraševanje razlikujeta po posameznih območjih.
Kljub obsežni funkcionalnosti je storitev žal prilagojena za severnoameriški
  trg, zato ni neposredno uporabna za nas.
% TODO: Add picture of carfax
Pregled obstoječih rešitev~\ref{tab:sorodne-resitve} je pokazal, da večina platform
 bodisi omogoča prodajo vozil brez naprednega sistema za vrednotenje 
 bodisi ponuja ločene plačljive storitve za ocenjevanje vrednosti vozil.
 Namen naše rešitve je združiti obe funkcionalnosti v enotnemu prototipu aplikacije,
  ki uporabnikom omogoča objavo vozil, pregled trga in samodejno oceno vrednosti vozila 
  na podlagi obstoječih primerljivih vozil.
\begin{table}[H]
\centering
\caption{Primerjava izbranih obstoječih rešitev za prodajo in vrednotenje vozil.}
\label{tab:sorodne-resitve}
\begin{tabularx}{\textwidth}{@{}l X c X@{}}
\toprule
\textbf{Rešitev} & \textbf{Trg} & \textbf{Prodaja oglasov} & \textbf{Ocena vrednosti vozila} \\
\midrule
mobile.de & Evropa, najbolj razširjen na nemškem trgu & Da & Osnoven prikaz cene (zeleno -- ugodna, rdeče -- precenjena) \\
AutoScout24 & Evropa--Nemčija, Belgija, Italija & Da & Priporočilo cene ob objavi oglasa \\
Avto.net & Slovenija & Da & Nima \\
Eurotax & Evropa--Nemčija, Slovenija, Švica, Italija & Ne & Natančno poročilo o ceni \\
CARFAX & ZDA & Ne & Podrobne ocene glede na VIN \\
\bottomrule
\end{tabularx}
\end{table}

\chapter{Ozadje}
\label{ch:struktura}
V tem poglavju predstavimo temeljne pojme in postopke, na katerih temelji naše delo. Najprej natančneje opišemo ključne
tehnologije s tega področja,
 kaj je oglas rabljenega vozila in kateri atributi vozila so za nas relevantni. 
 V nadaljevanju predstavimo še pristope za zajem podatkov,
  s katerimi smo pridobili gradivo za sistem vrednotenje vozil.

\section{Oglas in atributi vozila}
\textbf{Oglas} je vsebina, ki je namenjena predstavitvi izdelka ali storitve, ki se trži, prodaja ali kupuje~\cite{kotler2016marketing}.
Oglaševanje ima dolgo zgodovino; tiskani oglasi so se v evropskih časopisih redno pojavljali že od nekdaj, množično pa se začelo oglaševati z časopisi in revijami v 19.~in 20.~stoletju~\cite{tungate2007adland}.
Z razširjenostjo interneta se je oglaševanje večinoma preneslo v elektronsko obliko, kjer danes poteka pretežni del prodaje in nakupa rabljenih vozil.
V našem primeru obravnavamo spletne oglase za prodajo vozila.
 Vsebovati morajo vse ključne podatke. Katero vozilo se prodaja in njegove tehnične značilnosti,
  ki kupcu pomagajo pri odločitvi o nakupu. 
 Poleg tega mora oglas vsebovati še tržno ceno in kontaktne podatke prodajalca. 
 Oglas ostane objavljen na izbranem mestu, do odločitve umika ali nakupa vozila.

\subsection{Atributi vozila}
\textbf{Atributi} vozila ki vplivajo na ceno so njegove tehnične ali fizične lastnosti. 
To pa so znamka, model, izvedba (angl.\ \emph{trim}), barva, opis in prostornina motorja, vrsta goriva, vrsta menjalnika,
 pogon (sprednji, zadnji ali štirikolesni), letnik izdelave, datum prve registracije, prevoženi kilometri,
 moč motorja (izražena v kilovatih ali konjskih močeh), emisijski razred, poraba goriva, tip karoserije,
 število vrat, število sedežev, število prestav, število lastnikov, servisna zgodovina, stanje vozila (rabljeno, novo, poškodovano),
 država izvora, oprema (npr.\ klimatska naprava, navigacija, usnjeni sedeži) in cena.
  V naši rešitvi se moramo omejiti, zato se osredotočimo predvsem na nekaj 
  ključnih atributov: znamko, model, prevožene kilometre, letnik izdelave in ceno.
\textbf{Znamka} in \textbf{model} služita kot identifikatorja vozila, saj na njuni podlagi vemo, za katero vozilo gre.
 Naslednje dva atributa sta tista, ki najbolj vplivata na razliko vrednosti med kontkretnimi primeri istega modela vozila. 
 \textbf{Prevoženi kilometri} so pomembni, ker izražajo uporabo vozila; večje kot je število prevoženih kilometrov, večja je verjetnost okvar in obrabe. 
 \textbf{Letnik izdelave} pove starost vozila in s tem nakazuje pričakovane prihodnje stroške vzdrževanja. 
 \textbf{Cena} je glavni podatek, saj zagotavlja primerjalno vrednost med vozili. Na sliki~\ref{fig:primer-oglas} prikazujemo primer oglasa v našem prototipu.
\begin{figure}[H]
    \centering
    \includegraphics[width=0.8\textwidth]{primer-oglas.png}
    \caption{Primer oglasa na prototipu}
    \label{fig:primer-oglas}
\end{figure}

\section{Zajem podatkov iz spletnih virov}
\label{sec:zajem-podatkov}
\textbf{Spletni vir} je poljubna spletna storitev ali stran, ki ponuja podatke uporabne za obdelavo,
 na primer oglasni portal, javna podatkovna zbirka ali programski vmesnik.
Zajem podatkov iz spletnih virov predstavlja temeljni korak pri našem pristopu za vrednotenje rabljenih vozil na trgu.
Obseg starosti in kakovosti zbranih podatkov neposredno
vpliva na zanesljivost modelov za ocenjevanje cen~\cite{laender2002}. V splošnem ločimo tri
pristope za pridobivanje podatkov iz spleta: izkoriščanje javno dostopnih
podatkovnih zbirk, dostop prek vmesnikov API-jev ter samodejni zajem podatkov s spletnih strani~\cite{laender2002}

\subsection{Javno dostopne podatkovne zbirke}

\textbf{Javno dostopne podatkovne zbirke} (angl.\ \emph{open datasets}) predstavljajo enega najprimernejših virov za zbiranje strukturiranih podatkov,
 saj so podatki v njih že organizirani, kategorizirani in pripravljeni za neposredno
  analitično obdelavo~\cite{janssen2012open}. Primeri takšnih zbirk so odprti vladni portali, (kot sta Statistični urad Republike 
 Slovenije (SURS)~\cite{surs} in Eurostat~\cite{eurostat}), akademski repozitoriji ter portali prometnih agencij.
  Slabost tega pristopa je omejena razpoložljivost aktualnih tržnih podatkov.
Javne zbirke običajno ne vsebujejo oglasov z aktualnimi prodajnimi cenami vozil,
ki so nujno potrebni za ocenjevanje tržne vrednosti~\cite{gegic2019car}.
Za tuje trge sicer obstajajo javno dostopne zbirke oglasov, pridobljeni iz
oglasnih portalov, (npr. nemški trg (eBay Kleinanzeigen)~\cite{ebay_de_dataset}), ki so vendar za slovenski trg takšna javna zbirka
ne obstaja.

\subsection{Dostop prek vmesnikov API}
\textbf{Dostop do podatkov prek vmesnika API} (angl.\ \emph{application programming interface}) običajno poteka tako, da se uporabnik najprej registrira pri ponudniku~\cite{oauth2-rfc6749}.
Ob registraciji uporabnik prejme nosilni žeton (angl.\ \emph{bearer token}), ki ga mora priložiti v vsakemu zahtevku
 za namen avtentikacije~\cite{oauth2-rfc6749}. Večina vmesnikov API je plačljivih, saj ponudniki ustvarjajo prihodek s prodajo podatkov. 
 Brezplačni vmesniki so zato pogosto pomanjkljivi in vsebujejo zastarele podatke. V našem primeru smo vmesnike API uporabili
  predvsem za začetno polnitev podatkovne baze z vsemi znamkami in modeli vozil, za kar zadostujejo tudi brezplačni ponudniki.

\subsection{Samodejni zajem podatkov s spletnih strani}

\textbf{Samodejni zajem podatkov s spletnih strani} (angl.\ \emph{web scraping}) je postopek avtomatiziranega
pridobivanja strukturiranih informacij iz spletnih strani z razčlenjevanjem
njihove vsebine~\cite{laender2002}. Ta pristop se pogosto uporablja pri zbiranju
oglasov za prodajo rabljenih vozil. Na voljo je veliko spletnih strani, ki
vsebujejo obsežne zbirke takšnih oglasov, iz vsakega oglasa pa je mogoče
pridobiti ključne strukturirane atribute vozila, (npr. znamka, model, letnik
izdelave, prevoženi kilometri, vrsto goriva, tip menjalnika in navedena cena~\cite{gegic2019car}).
 Zajem poteka z orodji za avtomatizirano
brskanje po spletnih straneh, (kot so Playwright~\cite{playwright-docs} ali Selenium~\cite{selenium-docs}), ki simulirajo obisk uporabnika, ter z
razčlenjevanjem strukture dokumenta HTML za izluščitev relevantnih
podatkovnih polj~\cite{laender2002}. Problem zajema podatkov na tak način je ta,
da so podatki na spletni strani namenjeni za prikaz in ne za posredovanje podatkov uporabnika. Zato
 se morajo lastniki spletnih strani strinjati z tem, saj so podatki njihova last.

\subsection{Razčlenjevanje vsebine HTML}

Spletna stran je v osnovi dokument v jeziku HTML
(angl.\ \emph{HyperText Markup Language}), ki je organiziran v drevesno strukturo
DOM (angl.\ \emph{Document Object Model})~\cite{olston2010}. Vsak element
dokumenta (naslov, odstavek in podatkovne celice), je v tej strukturi
predstavljen kot vozlišče, ki ima lahko starševska in podrejena vozlišča.
Orodja za razčlenjevanje (kot sta knjižnici
\texttt{BeautifulSoup}~\cite{beautifulsoup-docs} in \texttt{lxml}~\cite{lxml-docs} v jeziku Python), to strukturo pretvorijo
v obliko, ki so primerna za programsko obdelavo.\\
Za izluščevanje podatkovnih elementov iz razčlenjene strukture se v praksi
najpogosteje uporabljata dve vrsti izbirnikov (angl.\ \emph{selectors}): izbirniki CSS in izbirniki XPath~\cite{darmawan2022}.
 Izbirniki CSS so bili prvotno zasnovani za oblikovanje
spletnih strani, a jih pri zajemu podatkov preusmerimo v poizvedovanje po
strukturi DOM~\cite{olston2010}. Izbirniki XPath (angl.\ \emph{XML Path
Language}) pa predstavljajo namenski poizvedovalni jezik za navigacijo po drevesnih
strukturah XML in HTML, ki poleg preprostega iskanja po oznakah in razredih
vsebuje tudi premikanje (navzgor po drevesu, iskanje po vsebini besedilnih
vozlišč ter vgrajene funkcije, (kot sta \texttt{contains()} in
\texttt{position()}~\cite{darmawan2022})).\\
V kontekstu zajemanja podatkov iz
oglasnih portalov za vozila to pomeni, da lahko izbirniki iz razčlenjenega
dokumenta izluščimo vsa ključna podatkovna polja oglasa.

\chapter{Razvoj rešitve}
\label{ch:razvoj}
Cilj diplomske naloge je bila izdelava prototipa spletne aplikacije, v katerem bodo uporabniki lahko
objavljali oglase za prodajo svojih vozil. Pred začetkom razvoja smo morali najprej jasno opredeliti zahteve.
Sledila ji je zasnova načrta naše rešitve, nato smo morali izbrati ustrezne
tehnologije, kar smo storili na podlagi lastnih izkušenj in predznanja ter pregleda
različnih forumov in strokovnih člankov. Ko smo zaključili razvoj prototipa spletne aplikacije, smo
se lotili še ločenega razvoja sistema za vrednotenje vozil, ki je napisan v programskem jeziku Python.
Python smo izbrali, ker je zelo primeren za obdelavo podatkov in razpolaga z bogatim
ekosistemom knjižnic, namenjenih analizi podatkov.\\

\section{Zahteve}
\label{sec:zahteve}
Pred začetkom razvoja smo opredelili zahteve, ki jih delimo na funkcionalne in nefunkcionalne. 
Pri opredelitvi funkcionalnih zahtev ločimo dve vrsti uporabnikov: \textbf{anonimnega} in \textbf{registriranega}. 
Anonimni uporabnik lahko dostopa le do pregleda oglasov, medtem ko je registriranemu uporabniku omogočen dostop
 do celotne funkcionalnosti sistema. Funkcionalne zahteve so naslednje:
\begin{itemize}
    \item Uporabnik se lahko registrira.
    \item Registriran uporabnik se lahko prijavi in odjavi.
    \item Oglaševanje vozil je omogočeno le registriranim uporabnikom.
    \item Registriran uporabnik lahko za vozilo z vnešenimi atributi (Znamka, model, letnik, prevoženi kilometri) oceni njegovo vrednost.
    \item Registriran uporabnik lahko všečka oglase drugih uporabnikov.
    \item Registriran uporabnik lahko ažurira svoje objavljene in všečkane oglase.
    \item Anonimni uporabnik lahko pregleduje objavljene oglase, ne more pa jih objavljati, všečkati ali ovrednotiti vozil.
\end{itemize}
Na sliki~\ref{fig:use-case} prikazujemo diagram primerov uporabe, ki povzema interakcije registriranega in anonimnega uporabnika s sistemom.
\begin{figure}[H]
    \centering
    \includegraphics[width=0.8\textwidth]{use-case.png}
    \caption{Diagram primerov uporabe za registriranega in anonimnega uporabnika.}
    \label{fig:use-case}
\end{figure}
Nefunkcionalne zahteve opredeljujejo kakovostne lastnosti sistema in so naslednje:
\begin{itemize}
    \item Spletna aplikacija je hitra in odzivna. Uporabniški vmesnik se na dejanja
    uporabnika odzove praktično takoj (v manj kot $100$~ms), zahtevki HTTP proti
    zalednemu sistemu (z izjemo vrednotenja vozila) pa se v večini primerov
    obdelajo v manj kot $300$~ms.
    \item Sistem za vrednotenje vozil vrne oceno v kratkem času (manj kot 2 sekundi).
    \item Uporabniški vmesnik je intuitiven in enostaven za uporabo. Uporabnik
    lahko glavne naloge (ocena vrednosti vozila, objava oglasa, prijava) opravi
    brez predhodnih navodil, saj so poti do njih vidne z glavne strani, obrazci
    pa sproti opozarjajo na napačno vnesene podatke.
    \item Spletna aplikacija je vizualno privlačna. Uporablja enoten slog
    (barve, tipografijo in razmike), ki je skladen na vseh pogledih, postavitev
    pa se prilagaja različnim velikostim zaslonov, od namiznih do mobilnih.
    \item Sistem ustrezno varuje vse uporabniške podatke, (gesla pred shranjevanjem zgosti in uporablja časovno omejene žetone JWT).
    \item Sistem ostane odziven tudi pri izvajanju zahtevnejših opravil, (uporablja sinhrono obdelavo).
    \item Arhitektura omogoča hitro in vzdržljivo dodajanje novih funkcionalnosti.
\end{itemize}

\section{Arhitektura rešitve}
Prototip spletne aplokacije je sestavljen iz uporabniškega vmesnika (odjemalec), strežniškega dela aplikacije (zaledni sistem) in ločenega sistema za vrednotenje vozil.
Uporabniški vmesnik je razvit z uporabo knjižnice React~\cite{react-docs} in programskega jezika TypeScript~\cite{typescript} v okolju Node.js različice 24~\cite{nodejs}. Zaledni sistem je razvit z uporabo ogrodja Express~\cite{expressjs-docs}
 v programskem jeziku JavaScript~\cite{flanagan2020javascript} v okolju Node.js različice 24.
Sistem za vrednotenje vozil je implementiran kot ločena storitev v programskem jeziku Python različice 3.11~\cite{python3.11}.
Na sliki~\ref{fig:architecture} je prikaz arhitekture.
\begin{figure}[H]
    \centering
    \includegraphics[width=0.9\textwidth]{arhitektura.png}
    \caption{Arhitektura prototipa spletne aplikacije}
    \label{fig:architecture}
\end{figure}
Za izvajanje posameznih komponent je uporabljena platforma Docker~\cite{docker} 
z vsebniki za, Zaledni sistem, Postgres(16), redis(7), mailhog in sistem za vrednotenje vozil.\\
Razvoj aplikacije je potekal v razvojnem okolju Visual Studio Code. Za avtomatizacijo preverjanja kode, izvajanja testov ter gradnje vsebnikov Docker je bil uporabljen sistem GitHub Actions~\cite{github-actions-docs}, medtem ko so se periodična opravila izvajala s pomočjo časovno načrtovanih opravil (cron).
Za gostovanje aplikacije je bil uporabljen virtualni strežnik ponudnika Hetzner~\cite{hetzner}. Namestitev in upravljanje aplikacije na strežniku je bilo izvedeno s pomočjo platforme Dokploy~\cite{dokploy}.

\section{Uporabne tehnologije}
Poleg temeljnih tehnologij, ki so relativno znane in jih zato ne predstavljamo posebej,
 smo med razvojem uporabili
še nekaj manjših ampak uporabnih knjižnic in orodij, za ožje probleme. V
nadaljevanju jih na kratko predstavimo.

\subsection{Upravljanje obrazcev s knjižnico React Hook Form}
Uporabniški vmesnik vsebuje več obsežnejših obrazcev za izpolnjevanje. Ročno upravljanje takšnih
obrazcev v knjižnici React~\cite{react-docs} zahteva, da za vsako polje posebej vzdržujemo stanje, obravnavamo
dogodke ob spremembi vrednosti in sami zapišemo logiko preverjanja vnosov, kar hitro privede
do obsežne in težko vzdržljive kode.\\
 Zato smo uporabili knjižnico React Hook Form~\cite{react-hook-form-docs}.
Knjižnica stanje obrazca upravlja prek nenadzorovanih komponent (angl.\ \textit{uncontrolled
components}) in referenc, kar pomeni, da se ob vsakem pritisku tipke ne izriše celoten obrazec,
temveč le tisti deli, ki jih sprememba dejansko zadeva. Posamezno polje prijavimo v obrazec,
 pri čemer hkrati navedemo še pravila preverjanja vnosa (angl.\
\textit{validation}), na primer obveznost polja, najmanjšo dolžino ali regularni izraz. Napake
so nato na voljo v objektu \texttt{errors}, iz katerega jih neposredno izpišemo ob pripadajočem
polju. Preverjanje vnosov na strani odjemalca uporabnika opozori na napake, še preden se
zahtevek sploh pošlje na strežnik, kar zmanjša število nepotrebnih zahtevkov.

\subsection{Predpomnjenje strežniškega stanja s knjižnico TanStack Query}
Del podatkov, ki jih prikazuje uporabniški vmesnik, izvira iz zalednega sistema in se
razmeroma redko spreminja. Značilen primer so seznami znamk in modelov vozil, ki jih
potrebujemo v skoraj vsakem obrazcu in filtru. Če bi te podatke ob vsakem prikazu komponente
znova pridobivali s strežnika, bi obremenjevali zaledni sistem in podaljševali
odzivni čas vmesnika.\\
Zato smo uporabili knjižnico TanStack Query~\cite{tanstack-query-docs}. Ta odgovore strežnika shrani v predpomnilnik,
 kjer je vsak odgovor označen s ključem poizvedbe (angl.\ \textit{query
key}). Ko komponenta zahteva podatke z že znanim ključem, jih knjižnica vrne takoj iz
predpomnilnika, zahtevek na strežnik pa izvede po potrebi v ozadju, prikaz pa osveži šele, ko
prispe nov odgovor. Čas, po katerem se shranjeni podatki obravnavajo kot zastareli, določimo za vsak ključ posebej.

\subsection{Testiranje zalednega sistema}
\label{sec:testiranje-be}
Zaledni sistem smo pokrili z avtomatiziranimi testi, ki jih izvajamo z ogrodjem
Jest~\cite{jest-docs}. Jest je testno ogrodje za okolje Node.js~\cite{nodejs}, ki v enem samem paketu združuje
izvajalnik testov, knjižnico trditev (angl.\ \textit{assertions}) in podporo za simulacijo
odvisnosti (angl.\ \textit{mocking}).\\
Ker večina logike zalednega sistema ni izpostavljena kot posamezna funkcija, temveč kot pot
vmesnika API, smo ogrodju Jest dodali še knjižnico Supertest~\cite{supertest-docs}. Ta omogoča
pošiljanje zahtevkov HTTP neposredno na aplikacijo Express, ne da bi strežnik dejansko poslušal
na omrežnih vratih. V testu tako opišemo celotno pot zahtevka: ga pošljemo na določeno
pot, priložimo žeton JWT ter nato preverimo statusno kodo in vsebino odgovora.
Testi se izvajajo samodejno na cevovodu neprekinjene integracije ob vsakem zahtevku za
združitev, kar podrobneje opisujemo v razdelku~\ref{subsec:ci}.

\subsection{Približno ujemanje nizov s knjižnico RapidFuzz}
\label{subsec:rapidfuzz}
Imena znamk in modelov vozil so pri različnih spletnih virih zapisana neenotno, zato jih z
natančnim ujemanjem nizov ni mogoče zanesljivo povezati z imeni v podatkovni bazi. Za takšne
primere smo uporabili knjižnico RapidFuzz~\cite{rapidfuzz-docs}, ki v jeziku
Python omogoča približno ujemanje nizov (angl.\ \textit{fuzzy string
matching}). Knjižnica podobnost med nizoma izrazi kot oceno med 0 in 100, izračunano na podlagi
urejevalne razdalje (angl.\ \textit{edit distance}), pri čemer so vse zahtevnejše operacije
implementirane v jeziku C++~\cite{cpp-standard} in so zato bistveno hitrejše od enakovrednih
izvedb v jeziku Python. Konkreten postopek ujemanja in izbrane pragove opisujemo v
razdelku~\ref{subsec:algoritem}.

\section{Načrtovanje podatkovne baze}

\subsection{Načrtovanje podatkovne baze aplikacije}
Pri načrtovanju podatkovne baze smo morali natančno premisliti, kako shranjujemo
oglase vozil, saj je to ključnega pomena za učinkovito delovanje odločitvenega
modela. Previsoka stopnja normalizacije bi lahko upočasnila poizvedbe, prenizka pa bi
vodila do podvajanja podatkov in z zagotavljanjem skladnosti v bazi~\cite{elmasri2015fundamentals}. Zasnovani
podatkovni načrt je predstavljen na sliki~\ref{fig:er-diagram}.\\

\begin{figure}[H]
    \centering
    \includegraphics[width=0.9\textwidth]{ErDiagramMain.png}
    \caption{ER-diagram glavnih entitet aplikacije.}
    \label{fig:er-diagram}
\end{figure}

Entiteti \texttt{Brand} in \texttt{CarModel} sta izločeni v ločeni tabeli, ker med njima obstaja
hierarhično razmerje ena-proti-mnogo: ena znamka (\texttt{Brand}) ima lahko poljubno število
modelov (\texttt{CarModel}), vsak model pa pripada natanko eni znamki. Shranjevanje znamk
in modelov kot vrednosti tipa \texttt{enum} ali prostega niza bi bilo neprimerno
iz dveh razlogov. Število možnih vrednosti je namreč preveliko za statično definicijo
z \texttt{enum}, prosti niz pa bi omogočal podvojene ali nekonsistentne vnose
(na primer \enquote{Audi A3} in \enquote{audi a3 Sportback} bi bila obravnavana kot različni
entiteti, čeprav gre za isti model). Z relacijsko vezavo prek tujih ključev je
zagotovljen referenčni nadzor nad vnesenimi vrednostmi ter enotnost podatkov.\\

Entiteta \texttt{CarSpecifications} je namerno ločena od entitete \texttt{CarListing}, ker vsebuje
tehnične podrobnosti vozila, ki niso relevantne pri vsakem poizvedovanju. Pri prikazu
seznama oglasov (npr. v podatkovni tabeli) zadostujejo le osnovni podatki iz entitete
\texttt{CarListing} (cena, prevoženi kilometri in letnik), podrobne specifikacije pa se naložijo le
ob ogledu posameznega oglasa. Ta ločitev zmanjša količino podatkov, ki se prenesejo
iz podatkovne baze pri pogostih poizvedbah, in s tem izboljša odzivnost aplikacije.\\

Entiteta \texttt{CarListingLike} predstavlja asociativno tabelo, ki realizira razmerje
mnogo-proti-mnogo med uporabniki in oglasi: en uporabnik lahko všečka
poljubno število oglasov, en oglas pa je lahko všeč pri poljubnem številu
uporabnikov. Tabela vsebuje tuji ključ na entiteto \texttt{User} in tuji ključ na entiteto
\texttt{CarListing}, kar zagotavlja referenčno integriteto in preprečuje podvojene vnose.
Entiteta \texttt{User} hrani ključne podatke o registriranih uporabnikih sistema.

\subsection{Načrtovanje podatkovne baze za ovrednotenje vozil}
V okviru iste podatkovne baze smo dodali ločeno shemo \texttt{market}, ki vsebuje
dve tabeli. Namenjeni zbiranju in obdelavi tržnih podatkov za sistem vrednotenja
vozil~\ref{fig:er-diagram-evaluation}..\\
\begin{figure}[H]
    \centering
    \includegraphics[width=0.9\textwidth]{ErDiagramEvaluation.png}
    \caption{ER-diagram sistema za vrednotenje vozil.}
    \label{fig:er-diagram-evaluation}
\end{figure}
Tabela \texttt{RawMarketListing} služi kot neobdelan arhiv vseh oglasov, ki so pridobljeni
iz zunanjih virov – javno dostopnih zbirk podatkov, spletnih vmesnikov API in z
zajemanjem spletnih strani. Tabela vsebuje naslednja polja:
\begin{itemize}
    \item \texttt{source} določa izvor podatkov (npr.\ ime portala ali vmesnika API),
    \item \texttt{externalId} je identifikator oglasa na izvornem sistemu, ki je določen statično v kodi in preprečuje podvojene vnose,
    \item \texttt{url} hrani neposredno povezavo do oglasa,
    \item \texttt{scraped\_at} je časovni žig trenutka zajema podatkov,
    \item \texttt{raw\_payload} vsebuje izvorni odgovor v obliki JSON brez kakršnekoli obdelave.
\end{itemize}
Izvorni odgovor je namerno shranjen v celoti, kar omogoča
kasnejšo ponovno obdelavo brez ponovnega zajema podatkov, če se normalizacijska
logika spremeni.\\
Tabela \texttt{NormalizedCarListing} vsebuje obdelano in normalizirano obliko
oglasov, ki so vezani na tabelo \texttt{RawMarketListing} prek tujega ključa
\texttt{raw\_listing\_id}. Razmerje med tabelama je ena-proti-mnogo: vsak surovi oglas
ima lahko več normaliziranih zapisov, katere razlikuje primarni ključ. Isto vozilo, pa razlikuje
tuji ključ od \texttt{RawMarketListing}. Polja so poenotena in tipizirana; cena, prevoženi
kilometri in letnik so shranjeni kot \texttt{integer}, znamka, model, vrsta goriva,
tip menjalnika, oblika karoserije in država pa kot \texttt{string}. Takšna
poenotena struktura omogoča neposredno primerjavo oglasov iz različnih virov pri
izračunu tržne vrednosti vozila.\\
Ločitev na dve tabeli predstavlja cevovoda surovih in normaliziranih podatkov
(angl.\ \textit{raw/normalized pipeline})~\cite{kleppmann2017designing}: surovi
podatki ostanejo nedotaknjeni kot revizijska sled in osnova za morebitno ponovno
obdelavo, normalizirana oblika pa je tista, ki se uporablja v analitičnih poizvedbah.

\section{Arhitektura zalednega sistema}
Zaledni sistem smo oblikovali z modularno arhitekturo, ki omogoča jasno ločitev odgovornosti (angl.\ \textit{separation of concerns}), 
kar olajša branje, vzdrževanje in nadaljnjo razširiteljnost kode~\cite{separation-concerns}. V naslednji sliki je
bolj podrobno prikazana arhitektura zalednega sistema~\ref{fig:backend-architecture}.
\begin{figure}[H]
    \centering
    \includegraphics[width=0.9\textwidth]{zaledni-sistem.png}
    \caption{Arhitektura zalednega sistema}
    \label{fig:backend-architecture}
\end{figure}
\subsection{Modul za poizvedbe v podatkovni bazi}
Za implementacijo podatkovnih modelov in poizvedb na podatkovno bazo smo uporabili knjižnico Sequelize~\cite{sequelize-orm}. Sequelize~\cite{sequelize-orm} je vmesnik za 
spreminjanje objektno-relacijskih preslikav (angl.\ \textit{Object-Relational Mapping}), ki omogoča preslikavo objektov iz programske kode v 
relacijske strukture podatkovne baze~\cite{sequelize-orm}. Tak pristop izboljšuje razvoj, saj omogoča programerjem, da namesto neposrednih poizvedb SQL delajo z objekti,
 kar zmanjša napake in poveča berljivost kode. Prednosti pristopa vključujejo tipsko varnost, 
avtomatično generiranje poizvedb za različne podatkovne baze ter intuitivnejše upravljanje relacij med entitetami.\\
Podatkovni modeli so v Sequelize definirani kot razredi, ki predstavljajo tabele v bazi podatkov~\cite{sequelize-orm}.
 Vsak model ima atribute, ki ustrezajo stolpcem v tabeli, ter metode za izvajanje temeljnih operacij s podatki (CRUD – Create, Read, Update, Delete). 
 Ključna pri oblikovanju modelov je tudi postavitev indeksov na poljih, ki se pogosto pojavljajo pri poizvedbah (npr. ID vozila, izvor oglasa ali časovni žigi).
  Dobro zasnovani indeksi so bistvenega pomena za zagotavljanje hitrih odgovorov poizvedb pri velikih količinah podatkov~\cite{database-indexing}.
   Primer postavitve povezave z bazo podatkov je v kodi~\ref{code:db-connection}.


\begin{center}
\captionof{listing}{Koda za vzpostavitev povezave z bazo podatkov.}
\label{code:db-connection}
\begin{minted}{javascript}
import { Sequelize } from 'sequelize';

const sequelize = new Sequelize(process.env.DB_NAME, process.env.DB_USER, process.env.DB_PASSWORD, {
  host: process.env.DB_HOST,
  port: process.env.DB_PORT,
  dialect: 'postgres',
  logging: false,
});

export default sequelize;
\end{minted}
\end{center}


\subsection{Avtentikacija in avtorizacija uporabnikov}
\label{sec:avtentikacija}
Za avtentikacijo uporabnikov smo implementirali sistem, ki temelji na žetonih JWT (JSON Web Token), kar omogoča varen in enostaven prenos podatkov o uporabnikovi pristnosti~\cite{jwt-rfc7519}. 
Žeton JWT je strukturiran dokument, ki je sestavljen iz treh delov: glave (angl.\ \textit{header}), ki vsebuje metapodatke o vrsti kodiranja; \textit{podatkov} (angl.\ \textit{payload}), ki nosijo informacije o uporabniku; in \textit{digitalnega podpisa}, ki zagotavlja celovitost in pristnost žetona~\cite{jwt-rfc7519}.
Postopek avtentikacije je razdeljen na dve fazi. V fazi registracije uporabnik vnese svoje geslo, ki ga pred shranitvijo v podatkovno bazo zgostimo z algoritmom bcrypt~\cite{bcrypt-password}. Bcrypt je specializirana kriptografska funkcija, namenjena zgoščevanju gesel. 
 V fazi prijave uporabnik pošlje svoje geslo in e-poštni naslov. Zaledni sistem pridobi shranjeno zgoščeno geslo iz baze in ga primerja z novim zgoščenim geslom.
  Če se gesli ujemata, zaledni sistem generira žeton JWT z informacijami o uporabniku in ga vrne
 uporabniškemu vmesniku, kjer ga uporabnik shrani v brskalniku največkrat v shrambo seje (angl. \ \textit{localStorage}).\\
Za zaščito poti vmesnika API smo implementirali sredinski sistem za preverjanje žetonov. Vmesnik je funkcija, ki se izvede pred vsakim zahtevkom in preveri,
 ali je uporabnik v zahtevku poslal veljaven žeton JWT. Dodatna varnostna plast je nastavljena s časovnimi omejitvami žetona, ki zagotavlja,
 da je žeton veljaven samo za določen čas. Po preteku roka veljavnosti se mora uporabnik ponovno avtentificirati in pridobiti nov žeton,
  kar zmanjšuje tveganje ob morebitnem razkritju starejših žetonov~\cite{jwt-security-best-practices} Koda za prijavo uporabnika je prikazana v kodi~\ref{code:login},
  prijava in registracija je pa sta na sliki~\ref{fig:registration-login}.

\begin{center}
\captionof{listing}{Koda za prijavo uporabnika.}
\label{code:login}
\begin{minted}{javascript}
export const loginUserService = async ({ email, password }) => {
  try {
    if (!email || !password) throw new MissingUserDataError();

    const user = await User.findOne({ where: { email } });

    if (!user) throw new AppError('Invalid email or password', 401);
    
    const isPasswordValid = await bcrypt.compare(password, user.password);

    if (!isPasswordValid) throw new AppError('Invalid email or password', 401);

    const token = jwt.sign({ id: user.id, email: user.email }, process.env.JWT_SECRET, {
      expiresIn: '7d',
    });

    await emailQueue.add('login', { type: 'login', to: user.email, userName: user.firstName });

    return { user, token };
  } catch (error) {
    throw error;
  }
};
\end{minted}
\end{center}

\begin{figure}[H]
    \centering
    \includegraphics[width=0.9\textwidth]{registracija-prijava.png}
    \caption{Prijava in registracija na prototipu.}
    \label{fig:registration-login}
\end{figure}

\subsection{Asinhrona obdelava podatkov}
Za naloge, ki zahtevajo obsežnejšo obdelavo podatkov ali imajo daljši čas
izvajanja, smo implementirali asinhrono obdelavo z uporabo sistema čakalnih
vrst nalog. Ta pristop omogoča, da strežnik takoj odgovori na zahtevek,
medtem ko se zahtevnejše operacije izvajajo v ozadju kot ločeni procesi. Na ta
način se izboljša splošna odzivnost
aplikacije~\cite{kleppmann2017designing}. Značilne naloge, ki zahtevajo
asinhrono obdelavo, so večinoma zajem podatkov iz zunanjih virov,
normalizacijo in transformacijo zajetih podatkov ter pošiljanje e-poštnih
obvestil uporabnikom~\cite{kleppmann2017designing} \\
Za upravljanje čakalnih vrst nalog v okolju Node.js smo izbrali knjižnico
BullMQ~\cite{bullmq-docs}, ki zagotavlja zanesljivo upravljanje tovrstnih nalog z vgrajeno
podporo za ponovne poskuse (angl.\ \textit{retry}), časovne omejitve (angl.\
\textit{timeout}) in obravnavo napak~\cite{bullmq-docs}. Sistem temelji na
kombinaciji podatkovnega skladišča Redis~\cite{redis-data-store} in delovnih procesov. Redis v tej
arhitekturi služi kot posrednik sporočil (angl.\ \textit{message broker}), ki
hrani stanje nalog in omogoča koordinacijo med posameznimi delovnimi procesi
~\cite{redis-data-store}.\\
Konfiguracija sistema zahteva tri ključne korake. Najprej se vzpostavi
povezava s strežnikom Redis~\cite{redis-data-store}. Nato se registrirajo
delovni procesi, ki poslušajo dodeljene vrste nalog, jih obdelujejo in
ustrezno posodabljajo stanje v skladišču Redis. S takšno zasnovo je
zagotovljeno, da se dolgotrajne operacije ne izvajajo na glavni niti
aplikacije, kar omogoča sočasno obdelavo in izboljšano razširljivost sistema.


\section{Arhitektura uporabniškega vmesnika}
Uporabniški vmesnik smo zasnovali iz več komponent, ki skupaj sestavljajo celoto. Tudi tu smo upoštevali načela
modularnosti, ki omogočajo ponovno uporabo komponent ter lažje vzdrževanje kode~\cite{modular-programming}.
Vsako komponento smo oblikovali tako, da opravlja natančno določeno nalogo in prikazuje omejen del vmesnika.
Komponente smo namenoma ohranili majhne, saj majhen obseg naravno usmerja
razvijalca k temu, da posamezna komponenta ostane osredotočena na eno samo funkcionalnost.
Komponente si med seboj izmenjujejo podatke po hierarhiji nadrejenih in podrejenih komponent. Nadrejena
komponenta posreduje podatke podrejenim komponentam prek lastnosti (angl.\ \textit{props}), podrejena
komponenta pa nadrejeni sporoča spremembe prek povratnih funkcij, ki jih nadrejena komponenta posreduje
navzdol.
Kot ogrodje za izgradnjo uporabniškega vmesnika smo izbrali React~\cite{react-docs}. V naslednjih podpoglavjih podrobneje
predstavimo posamezne dele te arhitekture. V sliki~\ref{fig:frontend-architecture} je prikazana arhitektura ter postavitev
odjemalskega sistema.\\
\begin{figure}[H]
    \centering
    \includegraphics[width=0.9\textwidth]{odjemalec-arhitektura.jpg}
    \caption{Arhitektura odjemalčevega sistema}
    \label{fig:frontend-architecture}
\end{figure}
\subsection{Odjemalska shramba}
Določene podatke si mora odjemalec zapomniti, kar pomeni, da jih mora shraniti na svoji strani. Prednost
odjemalske shrambe je, da odjemalcu ni treba nenehno pošiljati novih zahtevkov na strežnik, kar občutno
zmanjša njegovo obremenitev. Značilen primer uporabe je shranjevanje žetona JWT, ki ga uporabnik prejme
ob prijavi in ga nato pri vsakem zahtevku pošilje v glavi~\cite{jwt-rfc7519}. Na ta način imamo podatke
o prijavljenem uporabniku ves čas na voljo, dokler se ta ne odjavi ali dokler žeton JWT ne poteče.\\
Odjemalska shramba je uporabna tudi pri večkoračnih postopkih, saj si lahko zapomnimo, do katerega koraka
je uporabnik prišel. Če uporabnik namerno ali nenamerno zapre stran, ostane njegovo napredovanje ohranjeno
in lahko z delom nadaljuje. Pri shranjevanju podatkov na odjemalcu moramo biti pozorni, da hranimo le
neobčutljive podatke ter da stanje pravilno posodabljamo ob vsaki spremembi. Spletni brskalniki za trajno
hrambo podatkov na strani odjemalca ponujajo več mehanizmov, med drugim \texttt{localStorage} in
\texttt{sessionStorage}~\cite{webstorage-spec}. V nasprotnem primeru se lahko stanje na odjemalčevi strani razlikuje od stanja
na strežniku, kar lahko vodi do napak.\\
Za odjemalsko shrambo smo uporabili knjižnico Zustand~\cite{zustand-docs}, ki je preprosta, kompaktna in dobro vzdrževana.
Zustand~\cite{zustand-docs} temelji na minimalističnem pristopu k upravljanju stanja in za osnovno delovanje ne zahteva ovojnih
komponent (angl.\ \textit{provider}), kar ga loči od nekaterih drugih uveljavljenih rešitev, kot je recimo Redux~\cite{redux-docs}. Stanje je organizirano v obliki t.\,i.\ shramb (angl.\ \textit{store}), do katerih
dostopamo prek preprostih kavljev (angl.\ \textit{hooks}), pri čemer se ponovno izrišejo le tiste komponente,
ki dejansko uporabljajo spremenjeni del stanja. Zaradi majhne velikosti in neodvisnosti od ogrodja React
je knjižnica primerna tudi za projekte, pri katerih želimo ohraniti nizko zahtevnost in dobro zmogljivost.\\

\subsection{Zahtevki API}
Za pošiljanje zahtevkov smo uporabili knjižnico Axios, ki je ena izmed najbolj razširjenih knjižnic za
pošiljanje zahtevkov HTTP v okolju JavaScript~\cite{flanagan2020javascript,axios-docs}. Nad njo smo izdelali ovojnico \texttt{api},
ki iz okoljskih spremenljivk (angl.\ \textit{environment variables}) prebere naslov strežnika in ga uporabi
kot osnovni naslov URL za vse zahtevke. Tak pristop nam omogoča, da osnovni naslov in skupne nastavitve
določimo na enem mestu, s čimer se izognemo ponavljanju in poenostavimo morebitne kasnejše spremembe. V kodi~\ref{code:api-wrapper}
 je prikazana ovojnica, ki jo nato uvozimo v vsako datoteko, kjer definiramo funkcije za pošiljanje zahtevkov.

\begin{center}
\captionof{listing}{Ovojnica Axios za ustvarjanje odjemalca API.}
\label{code:api-wrapper}
\begin{minted}{javascript}
import axios from "axios";

const api = axios.create({
  baseURL: import.meta.env.VITE_API_BASE_URL,
  withCredentials: true,
});

export default api;
\end{minted}
\end{center}

Ovojnico smo nato uporabili pri vsaki definiciji funkcije za pošiljanje zahtevka. Funkcija vrne polje
\texttt{res.data}, ki vsebuje dejanski odgovor strežnika. Axios namreč vrne objekt s številnimi dodatnimi
podatki, ki niso vedno potrebni, zato zaradi preglednosti kode vrnemo le vsebino odgovora~\cite{axios-docs}. Po potrebi lahko
dostopamo tudi do statusne kode strežnika, glave odgovora in konfiguracije zahtevka. Primer definicije funkcije za
pridobivanje znamk vozil je prikazan v kodi~\ref{code:api-call}.

\begin{center}
\captionof{listing}{Primer definicije funkcije za pridobivanje znamk vozil.}
\label{code:api-call}
\begin{minted}{javascript}
export const fetchCarBrands = async () => {
  const res = await api.get<{ brandNames: string[] }>("/api/cars/brands");
  return res.data;
};
\end{minted}
\end{center}


\subsection{Razvojno orodje Vite}
Vsako izvorno kodo je treba pred prikazom v spletnem brskalniku ustrezno pripraviti oziroma zgraditi. To nalogo
v našem projektu opravlja orodje Vite~\cite{vite-docs}, ki omogoča hitro postavitev razvojnega okolja in izgradnjo končne različice
aplikacije. Za razliko od starejših orodij za združevanje kode (angl.\ \textit{bundler}) (kot
je Webpack~\cite{webpack-docs}, ki mora ob vsaki spremembi znova zgraditi večji del kode), Vite tega ne počne in je zato bistveno
hitrejši.\\
Delovanje orodja Vite~\cite{vite-docs} temelji na dveh ključnih načelih. Prvo je izraba nativne podpore modulom ECMAScript~\cite{esmodules-mdn}
(angl.\ \textit{ES modules}) v sodobnih brskalnikih. Ker sodobni brskalniki neposredno razumejo uvožene module,
 Vite med razvojem kode ne združuje, temveč posamezne module dostavi brskalniku po potrebi~\cite{vite-docs,esmodules-mdn}. Drugo načelo je uporaba orodja esbuild~\cite{esbuild-docs} za predhodno izgradnjo odvisnosti JavaScript~\cite{flanagan2020javascript}.
Orodje je napisano v jeziku Go~\cite{golang-docs} in je od 10 do 100-krat hitrejše od orodij za združevanje, ki so napisana
v jeziku JavaScript.\\
Pomembna prednost orodja Vite je zmožnost sprotnega osveževanja modulov (angl.\ \textit{Hot Module Replacement}),
pri katerem se osveži le tisti del kode, ki se je spremenil, brez izgube trenutnega stanja in podatkov aplikacije.
Ta zmožnost je pri razvoju izjemno pomembna, saj občutno pospeši povratno zanko med spremembo in prikazom ter
tako poveča produktivnost.

\section{Sistem za vrednotenje vozil}
Sistem za vrednotenje vozil smo zasnovali kot samostojno storitev, ki smo jo razvili v programskem jeziku Python.
 Pri razvoju smo z osnovnim jezikom postopoma dodajali pakete glede na potrebe. Za komunikacijo z zalednim
  sistemom smo potrebovali spletni strežnik, saj je zaledni sistem edini, ki ima dostop do sistema za vrednotenje vozil.
Za spletni strežnik in vmesnik API smo uporabili ogrodje FastAPI, ki omogoča enostavno izdelavo vmesnikov API v
 jeziku Python~\cite{fastapi-docs}. Na izseku kode~\ref{code:api-python} je prikazan primer definicije vmesnika API, ki vrne seznam znamk vozil iz podatkovne baze~\ref{code:api-python}.
\begin{center}
\captionof{listing}{Primer definicije vmesnika API z ogrodjem FastAPI}
\label{code:api-python}
\begin{minted}{python}
@app.get("/car-brands")
def get_car_brands():
    return getCarBrands()
\end{minted}
\end{center}
Podatkovno bazo smo ločili v samostojno shemo z lastnimi modeli, s čimer 
 smo se izognili morebitni zmedi pri poimenovanju in strukturi.
Za migracijo tabel in podatkov smo dodali paket Alembic, ki je zmogljivo orodje za upravljanje
 in posodabljanje podatkovnih baz~\cite{alembic-docs}. Alembic sledi shemi podatkovne baze in samodejno 
 ustvarja migracijske skripte, ki jih je mogoče verzionirati skupaj z izvorno kodo, kar zagotavlja ponovljivo 
 in nadzorovano spreminjanje strukture podatkovne baze.
Pri tem projektu smo se prvič srečali s tehnologijo cron~\cite{cron-man}. Ta omogoča samodejno
 izvajanje ukazov ob vnaprej določenih časovnih intervalih na strežniku ali v drugem okolju.
  Cron smo uporabili za zagon programov za samodejni zajem podatkov s spletnih strani za prodajo vozil. 
  Ti programi se na strežniku izvajajo dnevno oziroma tedensko.

\subsection{Algoritem za primerjavo vozil}
\label{subsec:algoritem}
\begin{figure}[H]
\centering
\includegraphics[width=0.8\textwidth]{algoritem-slika.png}
\caption{Algoritem za primerjavo vozil}
\label{fig:algorithm}
\end{figure}
Opis algoritma v nadaljevanju sledi korakom s slike~\ref{fig:algorithm}, od leve proti desni in od zgornje
vrstice k spodnji.
Začeli smo z zajemom podatkov o oglasov vozil, bodisi s samodejnim zajemom s spletnih strani za prodajo vozil
ali pa z dostopom do javno dostopnih virov. Sam postopek zajema smo podrobneje opisali v
poglavju~\ref{sec:zajem-podatkov}, njegovo samodejno periodično izvajanje pa temelji na tehnologiji cron; tu se
osredotočamo na nadaljnjo obdelavo zajetih podatkov.\\
Zajete podatke najprej shranimo v surovi, nestrukturirani obliki, tako kot pridejo z vira. S tem zajem ločimo od
obdelave: če postopek normalizacije pozneje popravimo ali v njem odkrijemo napako, ga lahko ponovimo nad
shranjenimi surovimi podatki.
Sledi normalizacija, s katero surove podatke pretvorimo v enotno obliko. Za vsak spletni vir smo morali
napisati ločen postopek normalizacije.\\
Za vsako vozilo normaliziramo in shranimo ceno, prevožene kilometre, znamko, model, moč motorja, vrsto menjalnika in vrsto goriva.
Pri normalizaciji smo vse te podatke sproti izpisovali in poskrbeli za zanesljiv sistem obvladovanja napak.
Zlasti v začetni fazi razvoja namreč pogosto prihaja do napak, za njihovo učinkovito odpravljanje pa je potreben
podroben izpis poteka izvajanja~\cite{logging-practices}.\\
Poleg razlik v strukturi se posamezni viri med seboj razlikujejo tudi v zapisu posameznih vrednosti. Cene so lahko zapisane z ločilom tisočic, 
ki jih je treba pred pretvorbo v število odstraniti. Prevoženi kilometri so pri enem viru zapisani kot celo število, pri drugem pa kot niz z enoto 
(na primer \enquote{123.456~km}). Datum prve registracije je pri enem viru zapisan v obliki \textit{mesec/leto}, pri drugem pa \textit{leto/mesec}. Podobno je bilo 
treba poenotiti kategorije, za tip menjalnika in vrsto goriva, saj so viri te podatke podajali v različnih jezikih.
Poseben korak v tej fazi je ujemanje znamke in modela vozila, ki je hkrati največji izziv normalizacije, saj sta ista znamka ali model na različnih virih pogosto zapisana drugače
 (okrajšave, dodatne besede, jezikovne razlike -- na primer nemško \enquote{3er} ali \enquote{Klasse}/\enquote{class}).
 Zato smo zgradili sistem ujemanja, ki surovo ime najprej poskuša natančno ujeti s kanoničnim imenom v podatkovni bazi, nato preveri
  tabelo že znanih psevdonimov (vodeno posebej za vsak vir), šele nazadnje pa uporabi približno (angl. \emph{fuzzy}) ujemanje nizov~\cite{fuzzy-string-matching}
   s knjižnico \texttt{RapidFuzz} (razdelek~\ref{subsec:rapidfuzz}). Prag ujemanja smo nastavili na 90~\% za znamke in 75~\% za modele.
    Nižji prag za modele smo izbrali zato, ker so imena modelov krajša in bolj podvržena lažnim ujemanjem. Vsako uspešno približno ujemanje se samodejno shrani kot nov psevdonim za dani vir,
     s čimer se sistem sproti uči. Tako se vsak naslednji pojav istega zapisa razreši že v fazi iskanja psevdonimov, s čimer se izognemo počasnejšemu približnemu ujemanju. Del kode, ki izvaja ujemanje znamke, je prikazan v kodi~\ref{code:brand-matching-python}.

\begin{center}
\captionof{listing}{Ujemanje znamke s dejanskim in približnim ujemanjem}
\label{code:brand-matching-python}
\begin{minted}{python}
def match_brand(raw: str, source: str):
    normalized = normalize_text(raw)
    brands = _get_all_brands()

    for brand_id, canonical_name in brands:
        if normalize_text(canonical_name) == normalized:
            return brand_id, canonical_name

    cached = _lookup_brand_alias(normalized, source)
    if cached:
        return cached

    brand_names = [normalize_text(b[1]) for b in brands]
    result = process.extractOne(normalized, brand_names, scorer=fuzz.WRatio)

    if result is None or result[1] < 90:
        print(f"[matching] No brand match for '{raw}', best score: {result[1] if result else 0}")
        return None

    brand_id, canonical_name = brands[result[2]]
    _store_brand_alias(normalized, brand_id, canonical_name, source)
    return brand_id, canonical_name
\end{minted}
\end{center}

S tem je prva faza s slike~\ref{fig:algorithm} zaključena: po normalizaciji so
vsi podatki o oglasih vozil poenoteni v enotno strukturo (znamka, model, letnik,
prevožena kilometrina, cena itd.), kar je predpogoj za primerjavo posameznih
oglasov.
Sledi druga faza, ocena tržne vrednosti. Njen prvi korak je vzeti vhodne parametre uporabnika.
Implementirali smo API \texttt{/valuate}, ki za
vhodne parametre sprejme znamko, model, letnik in prevoženo kilometrino vozila,
za katero uporabnik želi oceno tržne vrednosti. V tej fazi smo se namenoma omejili
zgolj na navedene štiri parametre, saj je bil primarni cilj dokazati, da
algoritem deluje na konceptualni ravni (angl.\ \emph{proof of concept}), preden
bi ga razširili na dodatne spremenljivke (npr.\ opremljenost, vrsto goriva,
menjalnik, stanje vozila).\\
 Dodajanje večjega števila parametrov brez zadostne
količine normaliziranih podatkov bi namreč vzorec primerljivih vozil dodatno
skrčilo, zaradi česar bi jih hitro ostalo premalo za smiselno statistično oceno.\\
Po svoji zasnovi algoritem ustreza pristopu primerljivih prodaj (angl.\
\emph{comparable sales approach}), ki je uveljavljena metoda vrednotenja nepremičnin in
vozil~\cite{cma-appraisal}. Cene predmeta pri tem pristopu ne ocenjujemo na
podlagi enega samega primerka, temveč na podlagi množice podobnih, dejansko
prodanih (oz.\ oglaševanih) primerkov, ki jim glede na stopnjo podobnosti
pripišemo različno utež. V statističnem smislu gre za obliko metode
$k$-najbližjih sosedov (angl.\ \emph{k-nearest neighbors}, k-NN)~\cite{knn-regression},
 pri kateri namesto ostre meje
\enquote{sosed / ni sosed} uporabimo zvezno utežno funkcijo, ki pada proti nič,
ko se primerjano vozilo bolj razlikuje od ocenjevanega.\\
Rezultat algoritma pri tem ni zgolj ena sama cena, temveč cenovni razpon z
dodatnimi pokazatelji kakovosti ocene; podrobneje ga opišemo pri zadnjem koraku
v razdelku~\ref{subsec:uteceno-povprecje}.

\subsection{Filtriranje kandidatov iz podatkovne baze}
\label{subsec:filtriranje}
Tukaj je opisan korak iskanja primerljivih vozil s slike~\ref{fig:algorithm}.
Izvede se s poizvedbo SQL, ki iz tabele
\texttt{market.normalized\_market\_listing} izlušči kandidate za primerjavo.
Znamko in model filtriramo strogo (po enakosti), saj najmočneje
določata segment in bazno ceno vozila.\\
Numerični spremenljivki pa filtriramo z dovoljenim odstopanjem: pri kilometrini
znaša $\pm 40\,000$~km (konstanta \texttt{KILOMETERS\_DIFFERENCE}), pri letniku
pa $\pm 3$~leta (konstanta \texttt{YEAR\_DIFFERENCE}). Navedeni meji nista bili
izbrani naključno, temveč predstavljata kompromis med natančnostjo in velikostjo
vzorca. Če meji postavimo ožje, so vrnjena vozila bolj primerljiva glede na ocenjevano vozilo, a je 
hkrati manjši vzorec, iz katerega algoritem izračuna
oceno, kar poveča statistično negotovost. Gre za znano razmerje med
pristranskostjo in varianco (angl.\ \emph{bias--variance trade-off})~\cite{bias-variance}.
Ker je podatkovna baza normaliziranih podatkov vsaj v tej fazi razvoja še razmeroma
majhna, smo se zavestno odločili za širši interval, da algoritem sploh pridobi
dovolj velik vzorec za smiselno oceno. Ob prihodnji rasti podatkovne baze bi
bilo smiselno ta interval zožiti oziroma ga določiti dinamično (v\
odvistnosti od števila zadetkov ožjega intervala), s čimer bi izboljšali
natančnost ocene, brez da bi ta izhajala iz premajhnega vzorca.

\subsection{Točkovanje podobnosti posameznega vozila}
\label{subsec:tockovanje}
Tukaj opišemo korak uteževanja podobnosti s slike~\ref{fig:algorithm}.
Za vsako vozilo iz filtrirane množice izračunamo oceno podobnosti (angl.\
\emph{score}) glede na ocenjevano vozilo. Uporabili smo dve delni oceni, (prvo za
kilometrino in drugo za letnik), ki ju nato združimo v skupno oceno. Koda~\ref{code:car-scoring-python} predstavlja ta del izračuna.
\begin{center}
\captionof{listing}{Izračun ocene podobnosti posameznega vozila}
\label{code:car-scoring-python}
\begin{minted}{python}
def get_scores_of_each_car(data_compare: dict, cars: list) -> list[tuple]:
    results = []
    for car in cars:
        absolute_dif = abs(int(car["kilometers"]) - int(data_compare["kilometers"]))
        kilometers_score = max(0, 100 * (1 - absolute_dif / KILOMETERS_DIFFERENCE))

        absolute_dif = abs(int(car["registration_year"]) - int(data_compare["year"]))
        registration_year_score = max(0, 100 * (1 - absolute_dif / YEAR_DIFFERENCE))

        total_score = kilometers_score * 0.6 + registration_year_score * 0.4
        results.append((float(car["price"]), total_score))
    return results
\end{minted}
\end{center}

Obe delni oceni sta zasnovani po istem načelu: vozilo, ki je identično
ocenjevanemu, prejme največ 100 točk, vozilo na robu
dovoljenega intervala (z odstopanjem natanko $40\,000$~km ali 3 leta) pa 0 točk. Med
tema skrajnostma ocena linearno pada, funkcija \texttt{max(0, ...)} zagotavlja,
da ocena nikoli ne postane negativna. Ker vozila zunaj dovoljenega intervala
zaradi filtriranja v poizvedbi SQL sploh ne pridejo do tega koraka, gre pri
tem predvsem za dodatno varnostno mejo.\\
Za linearno padajočo funkcijo (namesto eksponentne) smo se odločili
iz dveh razlogov. Najprej zaradi enostavnosti in jasnosti, saj je linearna
funkcija bolj razumljiva in preverljiva, kar je za prikaz delovanja koncepta, pomembnejše od mogočega izboljšanja
natančnosti. Drugi razlog enostavno umerjanje, saj z eno samo konstanto natančno
določimo, kje ocena doseže ničlo.\\
Skupno oceno podobnosti izračunamo kot uteženo vsoto obeh delnih ocen, pri čemer
kilometrini pripišemo višjo utež ($60\,\%$) kot letniku ($40\,\%$).
Vozili istega letnika se tako bolj razlikujeta po dejanski obrabi, če je prvo
prevozilo $50\,000$~km, drugo pa $200\,000$~km. Vozili s podobno
kilometrino, a nekoliko različnim letnikom pa ostajata bolj primerljivi.

\subsection{Izračun ocenjene cene z uteženim povprečjem}
\label{subsec:uteceno-povprecje}
Tukaj sta opisana zadnja koraka druge faze s slike~\ref{fig:algorithm} in sicer
izračun ocene in njegov rezultat.
Ko vsa primerljiva vozila ovrednotimo z oceno podobnosti, končno ocenjeno ceno
izračunamo kot uteženo povprečje njihovih cen, pri čemer je utež vsakega vozila
prav njegova izračunana ocena podobnosti:

\begin{equation}
  \label{eq:weighted-mean}
  \texttt{estimated\_price} =
  \frac{\sum_i \texttt{price}_i \cdot \texttt{score}_i}{\sum_i \texttt{score}_i}
\end{equation}

Uporaba uteženega povprečja namesto navadnega aritmetičnega je ključna: če bi
izračunali navadno povprečje cen vseh vozil znotraj intervala, bi vozilo z roba intervala
($39\,000$ več km), prispevalo enako kot vozilo, ki ima le $500$ več km. To
je neustrezno, saj je slednje veliko boljši pokazatelj dejanske cene. Z
uteževanjem po vrednosti \texttt{total\_score} pa bolj podobna vozila sorazmerno
bolj vplivajo na končno oceno, vozila na robu intervala pa prispevajo le malo.
Standardna oblika uteženega povprečja (angl.\ \emph{weighted
mean})~\cite{weighted-mean}, poskrbi da vsota uteži v imenovalcu pravilno
določi rezultate, ne glede na število vozil ali velikost njihovih ocen.\\
Poleg ocenjene cene algoritem vrne tudi indeks zaupanja (angl.\
\emph{confidence}), ki je izračunan po enačbi:

\begin{equation}
  \label{eq:confidence}
  \texttt{confidence} = \min\!\left(1{.}0,\ \frac{n}{20}\right),
\end{equation}

pri čemer je $n$ število vozil, najdenih znotraj filtriranih meja. To pove
uporabniku stopnjo zaupanja v ceno, ki smo mu jo posredovali. Če je število
primerjalnih vozil veliko, je indeks zaupanja bližje 1, sicer pa
je indeks zaupanja bližje 0. Tudi tu je vrednost 20 določena hevristično in bi jo lahko spremenili\\
Rezultat algoritma tako ni zgolj ena sama cena. Vmesnik \texttt{/valuate} vrne
ocenjeno (pričakovano) prodajno ceno, poleg nje pa še spodnjo mejo ugodne cene
(\texttt{goodPrice}), zgornjo mejo precenjene cene (\texttt{badPrice}) in iz
njiju izpeljano razmerje \texttt{ratio}, ki pove, kako razpršene so cene
primerljivih vozil na trgu. Pri tem sta pomembni dve podrobnosti. Prvič, pri
izračunani ceni upoštevamo popust, saj se rabljena vozila praviloma prodajo
ceneje, kot so oglaševana. Drugič, razmerje
\texttt{ratio} med \texttt{badPrice} in \texttt{goodPrice} je poleg indeksa
zaupanja drugi pokazatelj kakovosti ocene: meri širino cenovnega razpona
primerljivih vozil in s tem, koliko prostora ima uporabnik pri
pogajanju.

\section{Infrastruktura in avtomatizacija}
Poleg implementacije posameznih komponent je bil pomemben del
razvoja tudi vzpostavitev infrastrukture, ki zagotavlja konsistentno delovanje
sistema v različnih fazah njegovega življenjskega cikla; od lokalnega razvoja,
prek avtomatiziranega preverjanja kode, do namestitve v produkcijsko okolje.
Namen tovrstne avtomatizacije je zmanjšati število ročnih posegov, ki so
podvrženi napakam, in zagotoviti, da se aplikacija v vseh okoljih obnaša enako,
ne glede na razlike v konfiguraciji računalnika ali
strežnika~\cite{devops-handbook}. Tak pristop temelji na načelu usklajenosti
okolij (angl.\ \emph{environment parity}), po katerem naj bodo razvojno, testno
in produkcijsko okolje čim bolj podobna, s čimer zmanjšamo tveganje za napake, ki se
pojavijo šele ob namestitvi~\cite{twelve-factor}.

\subsection{Postavitev razvojnega okolja z vsebnikih}
Za vsebnikizacijo razvojnega okolja smo uporabili orodje Docker Compose~\cite{docker-compose-docs}, ki
omogoča enostavnejšo vzpostavitev več med seboj povezanih vsebnikov. Orodje deluje tako, da v datoteki
\texttt{docker-compose.yml} opišemo vse potrebne storitve
in njihove medsebojne odvisnosti, kar omogoča zagon celotnega okolja z enim
samim ukazom. To je še posebej uporabno pri razvoju in testiranju. Za lokalno
preizkušanje pošiljanja elektronske pošte smo uporabili orodje MailHog~\cite{mailhog}, ki
odhodno pošto prestreže in jo prikaže v spletnem vmesniku, brez
dejanskega pošiljanja na pravi naslov. Za zaledni sistem in sistem za vrednotenje vozil smo napisali
ločeni datoteki \texttt{Dockerfile}, v katerih določimo, jezik
ter okolje sistema ter katere ukaze naj vsebnik ob zagonu izvede.\\
Ker se storitve med seboj zaganjajo z različnimi hitrostmi, smo za pravilno
zaporedje zagona uporabili \texttt{depends\_on} v kombinaciji z
definicijami \texttt{healthcheck}. Brez tega bi se zaledni sistem lahko poskusil prezgodaj
povezati s podatkovno bazo, kar bi
ob zagonu okolja povzročilo napako. Odvisne storitve, zato vsebujejo ukaz za preverjanje pripravljenosti (\texttt{pg\_isready}
pri sistemu PostgreSQL oziroma \texttt{redis-cli ping} pri sistemu Redis~\cite{redis-data-store})
 in se zaženejo šele, ko je preverjanje uspešno.\\
Za shranjevanje podatkov smo uporabili poimenovane nosilce podatkov (angl.\
\emph{named volumes}): prvi skrbi za trajnost podatkovne baze med ponovnimi
zagoni vsebnikov, drugi pa predpomni mapo \texttt{node\_modules} znotraj
zalednega vsebnika, saj bi jo sicer prepisala mapa gostitelja. Za
osveževanje ob spremembah
(angl.\ \emph{hot reload}), uporabljamo prevezavo map (angl.\ \emph{bind mount})
iz gostiteljevega datotečnega sistema v vsebnik, tako da se spremembe v izvorni
kodi takoj odražajo v delovanju aplikacije.\\
Odjemalskega sistema v lokalno razvojno okolje namenoma nismo vključili, saj
njegova vključitev zaradi nenehnih ponovnih izgradenj slike ob vsaki spremembi
bistveno upočasni razvoj. Med razvojem, ga tako poganjamo neposredno na
gostitelju.\\
Za produkcijsko okolje uporabljamo ločeno datoteko
\texttt{docker-compose.prod\allowbreak.yml}, ki namesto lokalne gradnje slik uporablja
vnaprej zgrajene slike iz registra vsebnikov (angl.\ \emph{container registry}).
Storitve poveže v skupno omrežje ter doda poimenovan nosilec podatkov za
shranjevanje naloženih datotek.

\subsection{Neprekinjena integracija}
\label{subsec:ci}
Neprekinjena integracija (angl.\ \emph{continuous integration}) je praksa,
pri kateri se ob vsaki spremembi kode samodejno izvedejo preverjanja (testiranje,
izgradnja, statična analiza kode). Na ta način napake odkrijemo, še preden se
zlijejo v glavno vejo ali dosežejo druge razvijalce~\cite{fowler-ci}. Namesto ročnega testiranja razvijalca pred združitvijo sprememb, 
teste izvede strežnik, vedno na enak način.\\
Za implementacijo neprekinjene integracije smo uporabili orodje GitHub Actions~\cite{github-actions-docs}. Njegovo delovanje opredeljujejo nadlednji pojmi:

\begin{itemize}
    \item \textbf{Potek dela} (angl.\ \emph{workflow}) je določen v datoteki \texttt{.yml}
    v mapi \texttt{.github\allowbreak/workflows/}, ki določa, kdaj naj se
    (npr.\ ob dogodku \texttt{pull\_request} ali \texttt{push}) sprožijo določena
    opravila.
    \item \textbf{Opravilo} (angl.\ \emph{job}) se izvede na svežem virtualnem
    strežniku (npr.\ \texttt{ubuntu-latest}), ki ga GitHub ob zagonu ustvari in
    po zaključku izbriše.
    \item \textbf{Koraki} (angl.\ \emph{steps}) znotraj opravila so zaporedni
    ukazi. To so lahko neposredni lupinski ukazi (npr.\ \texttt{run: pnpm test})
    ali vnaprej pripravljeni in ponovno uporabni gradniki (angl.\ \emph{actions}),
    ki jih razvije GitHub ali skupnost. Tako gradnik
    \texttt{actions/checkout} prenese izvorno kodo, gradnik
    \texttt{pnpm/action-setup} pa namesti upravitelja paketov. Namesto
    lastnoročno napisane skripte (za namestitev okolja Node, predpomnjenje
    odvisnosti in nastavitev spremenljivke \texttt{PATH}) tako uporabimo že obstoječi
    gradnik.
    \item \textbf{Storitve} (angl.\ \emph{services}) so pomožni vsebniki, ki
    tečejo vzporedno z opravilom. Uporabni so, takrat ko testi potrebujejo pravo
    podatkovno bazo namesto njene simulacije (angl.\ \emph{mock}).
\end{itemize}

Potek dela v datoteki \texttt{ci.yml}, smo določili tako, da se ob
vsakem novem zahtevku za združitev (angl.\ \emph{pull request}) izvedejo vsi testi
zalednega sistema ob tem se preveri, ali se uspešno izvedejo. Potek dela v datoteki
\texttt{docker-publish.yml} je zasnovan še korak dlje: ob zahtevku za združitev se
zgolj preveri, ali se slika Docker uspešno zgradi, ob združitvi v glavno vejo pa se
sliko tudi dejansko zgradi, digitalno podpiše in objavi. Ta korak že sodi na področje
neprekinjene dostave (angl.\ \emph{continuous delivery}), torej samodejne
distribucije zgrajene programske rešitve~\cite{humble-cd}. V naslednji sliki je prikazan 
primer izgradnje in izvajanja testov na cevovodu~\ref{fig:ci-cd}.

\begin{figure}[H]
\centering
\includegraphics[width=0.8\textwidth]{ci-cd.png}
\caption{Samodejno izvajanje testov nad vejo na cevovodu}
\label{fig:ci-cd}
\end{figure}

\subsection{Namestitev in gostovanje}
\label{sec:namestitev}
Za gostovanje smo izbrali eno izmed cenovno najugodnejših možnosti in sicer
najeti navidezni zasebni strežnik (angl.\ \emph{virtual private server})
pri ponudniku Hetzner~\cite{hetzner}. Strežnik razpolaga s 4~GB delovnega pomnilnika in dvema
procesorskima jedroma, kar zadostuje za zagon vseh treh storitev. Po prvi prijavi v
strežnik smo najprej posodobili operacijski sistem na najnovejšo različico,
nato pa namestili orodji Docker in Docker Compose.\\
Za upravljanje aplikacije smo uporabili orodje Dokploy~\cite{dokploy}, ki omogoča enostavno
uporabo datotek \texttt{docker-compose} v produkcijskem okolju in hkrati močno
poenostavi postopek posodabljanja aplikacije. V njem lahko nastavimo tudi
opravila, ki se izvajajo periodično, (podobno kot storitev \texttt{cron}). Na
strežnik je treba namestiti še povezovalno komponento med orodjema Dokploy in
Docker, nato pa Dokploy sam poskrbi za izgradnjo aplikacije: iz repozitorija
GitHub zajame najnovejšo različico kode in zgradi pripadajočo sliko neposredno
na strežniku. Tak postopek omogoča, da se vsaka sprememba v produkcijskem okolju
odrazi v nekaj minutah.\\
Dokploy je odprtokodno in brezplačno orodje, najem strežnika pri ponudniku
Hetzner pa predstavlja nizek mesečni strošek (trenutno 6.50€ na mesec).
Levo na sliki~\ref{fig:hezyner-dokploy} je upravljalna stran, desno je pa uporabniški vmesnik Dokploy.
\begin{figure}[H]
\centering
\includegraphics[width=0.9\textwidth]{primerjava-hetzner-dokploy.png}
\caption{Upravljalna stran Hetzner in Dokploy}
\label{fig:hezyner-dokploy}
\end{figure}

% =====================================================================
%  Poglavje: Analiza in vrednotenje rešitve
%  Vključeno iz diploma-predloga.tex prek % =====================================================================
%  Poglavje: Analiza in vrednotenje rešitve
%  Vključeno iz diploma-predloga.tex prek % =====================================================================
%  Poglavje: Analiza in vrednotenje rešitve
%  Vključeno iz diploma-predloga.tex prek \input{podporno/analiza.tex}
%
%  Številke so vpisane iz izhodov skript v CarStockMarketEvaluation/analysis/.
%  Preostale oznake \res{...} (rdeče) je treba še izpolniti:
%    - obremenitveni test (benchmarks/loadtest.py)
%    - sledljivost testov, Jest coverage, Lighthouse
%    - tabeli udeležencev in najdenih težav
%  Slike (PDF) iz analysis/out/ prekopiraj v mapo slike/.
% =====================================================================

\newcommand{\res}[1]{\textcolor{red}{\textbf{[#1]}}}

\chapter{Analiza in vrednotenje rešitve}
\label{ch:analiza}

V tem poglavju sistematično pregleadamo in analiziramo razvito rešitev. Najprej iz zahtev in
izhodišč, ki smo si jih zastavili na začetku, izpeljemo raziskovalna vprašanja in za vsako določimo merilo
uspešnosti (razdelek~\ref{sec:rq}). Nato opišemo merilno okolje in postopke
(razdelek~\ref{sec:metodologija}) ter natančno definiramo vse uporabljene mere
(razdelek~\ref{sec:mere}). Sledijo rezultati: značilnosti zbranih
podatkov~(\ref{sec:podatki}), natančnost ocene
vrednosti~(\ref{sec:natancnost}), občutljivost algoritma na
parametre~(\ref{sec:obcutljivost}), latenca posameznega
vrednotenja~(\ref{sec:latenca}) in obnašanje pod
obremenitvijo~(\ref{sec:obremenitev}). Zatem preverimo funkcionalne in
nefunkcionalne zahteve (\ref{sec:funkcionalne}, \ref{sec:nefunkcionalne}) ter
predstavimo uporabniški preizkus~(\ref{sec:uporabniki}) in simulacijo
uporabe~(\ref{sec:simulacija}). Poglavje sklenemo z obravnavo groženj
veljavnosti~(\ref{sec:groznje}) in povzetkom ugotovitev glede na
zastavljena izhodišča~(\ref{sec:povzetek-analize}).

% =====================================================================
\section{Raziskovalna vprašanja in merila uspešnosti}
\label{sec:rq}

Na začetku diplomskega dela (poglavje~\ref{ch:uvod}) smo si zastavili, da je
mogoče uporabniku na podlagi zbranih tržnih podatkov in primerjave s podobnimi
vozili ponuditi \emph{uporabno} in \emph{dovolj natančno} oceno tržne vrednosti
vozila in da vključitev takega sistema izboljša uporabniško izkušnjo. Iz tega
izhodišča in iz nefunkcionalnih zahtev (razdelek~\ref{sec:zahteve}) izpeljemo
štiri raziskovalna vprašanja (RV1--RV4) in za vsako določimo merilo uspešnosti.

\emph{Mediana absolutne odstotne napake} (MdAPE) je osrednja mera natančnosti v
tem poglavju; nanjo se sklicujejo merila uspešnosti in vsi nadaljnji razdelki.
Izračunamo jo tako, da za vsak testni primer določimo absolutno odstotno napako,
torej za koliko odstotkov se ocenjena cena razlikuje od dejanske, nato pa čez
vse primere vzamemo mediano teh vrednosti. MdAPE $=15\,\%$ tako pomeni, da je pri
polovici vozil ocena znotraj $15\,\%$ od dejanske cene, pri drugi polovici pa je
odstopanje večje; manjša kot je vrednost, natančnejše so ocene. Mediano
uporabljamo namesto povprečja, ker je manj občutljiva na posamezne skrajne
napake. Natančne definicije MdAPE in preostalih mer so v
razdelku~\ref{sec:mere}.

Raziskovalna vprašanja so:

\begin{itemize}
  \item \textbf{RV1 --- natančnost ocene.} Kako natančno sistem oceni tržno
        vrednost vozila v primerjavi z dejanskimi oglasi na trgu?
        \emph{Merilo:} MdAPE je pri vozilih z vsaj $10$ primerljivimi
        vozili manjša od $15\,\%$, ocena pa je bistveno natančnejša od
        preproste (naivne) referenčne metode.
  \item \textbf{RV2 --- odvisnost natančnosti od podatkov.} Od česa je
        natančnost ocene najbolj odvisna (velikost vzorca primerljivih vozil,
        razpršenost ponudbe, segment vozila)?
  \item \textbf{RV3 --- hitrost in razširljivost.} Ali sistem zadošča
        nefunkcionalnim zahtevam glede odzivnosti tudi pod sočasno
        obremenitvijo, in kje je zgornja meja trenutne postavitve?
  \item \textbf{RV4 --- uporabnost.} Ali ciljni uporabniki znajo brez navodil
        opraviti glavne naloge (registracija, objava oglasa, vrednotenje)?
\end{itemize}
% =====================================================================
\section{Metodologija in merilno okolje}
\label{sec:metodologija}

\paragraph{Merilno okolje.}
Meritve hitrosti (razdelka~\ref{sec:latenca} in~\ref{sec:obremenitev}) so
izvedene v \emph{produkcijskem} okolju, torej na navideznem zasebnem strežniku
z dvema procesorskima jedroma in $4$~GB pomnilnika ter s tehnologijami in
njihovimi različicami, opisanimi v poglavju~\ref{ch:razvoj}
(razdelek~\ref{sec:namestitev}). Analitične in merilne skripte tečejo nad isto
podatkovno bazo kot je na produkciji.

\paragraph{Posnetek podatkov.}
Meritve v tem poglavju temeljijo na posnetku podatkovne baze iz septembra
$2026$, ki obsega približno $86.000$ normaliziranih oglasov
(razdelek~\ref{sec:podatki}). Ker se posnetek med analizo ne spreminja, so
rezultati ponovljivi.

\paragraph{Orodja.}
Za meritve in analizo v tem poglavju smo napisali namenske skripte, ločene od
produkcijske kode; ena od njih je dokumentirana kopija ocenjevalnega
algoritma, parametrizirana za eksperimente. Statistično obdelavo in izris
grafov smo izvedli s knjižnicama \texttt{NumPy}~\cite{harris2020numpy} in
\texttt{Matplotlib}~\cite{hunter2007matplotlib}.

\paragraph{Natančnost.}
Natančnost ocene (razdelek~\ref{sec:natancnost}) merimo z naključno izbranim
normaliziranim oglasom: njegove atribute (znamka, model, letnik, kilometrina)
uporabimo kot vhod v algoritem, sam oglas pa izločimo iz primerjalne množice. Nato izračunamo oceno in jo
primerjamo z \emph{dejansko oglaševano} ceno tega oglasa. Postopek ponovimo za
$3000$ naključnih oglasov.

% =====================================================================
\section{Uporabljene mere}
\label{sec:mere}

\paragraph{Latenca}
 je čas, ki preteče od trenutka tik pred oddajo zahtevka do
trenutka, ko odjemalec prejme celoten odgovor. Glede na to, kolikšen del te
poti zajamemo, ločimo tri vrste:
\begin{itemize}
  \item[(a)] \emph{latenca obdelave na strežniku} --- čas izvajanja funkcije
        za vrednotenje, od začetka do konca, brez omrežja;
  \item[(b)] \emph{latenca poizvedbe SQL} --- čas izvajanja poizvedbe v podatkovni bazi,
   od začetka do konca, brez omrežja;
  \item[(c)] \emph{latenca od konca do konca} (angl.\ \emph{end-to-end}) ---
        (a) skupaj z omrežnim obhodom med odjemalcem in strežnikom ter
        obdelavo zahtevka HTTP; to je čas, ki ga zazna uporabnik.
\end{itemize}
V razdelku~\ref{sec:latenca} merimo vrsti (a) in (b), v
razdelku~\ref{sec:obremenitev} pa vrsto (c).

\paragraph{Percentili.}
\textbf{Percentil} $p_k$ je vrednost, pod katero pade $k\,\%$ meritev. Poročamo
$p_{50}$ (mediana), $p_{95}$ in $p_{99}$.

\paragraph{Prepustnost, sočasnost, nasičenje.}
Obnašanje sistema pod obremenitvijo opišemo s tremi količinami:
\begin{itemize}
  \item \textbf{prepustnost} --- število uspešno obdelanih zahtevkov na
        sekundo;
  \item \textbf{sočasnost} --- število odjemalcev, ki hkrati držijo odprt
        zahtevek;
  \item \textbf{nasičenje} (\emph{koleno}) --- točka, v kateri dodatna
        sočasnost ne poveča več prepustnosti, latenca pa začne naraščati
        približno linearno.
\end{itemize}
Notranjo skladnost meritev preverimo z \textbf{Littlovim zakonom}~\cite{little1961},
\begin{equation}
  \label{eq:little}
  L = \lambda \cdot W,
\end{equation}
kjer je $L$ povprečno število zahtevkov v sistemu, $\lambda$ prepustnost in
$W$ povprečna latenca.

\paragraph{Napaka ocene.}
Za oglas $i$ z dejansko ceno $p_i$ in ocenjeno ceno $\hat p_i$ definiramo
absolutno odstotno napako
\begin{equation}
  \mathrm{APE}_i = \frac{\lvert \hat p_i - p_i \rvert}{p_i}.
\end{equation}
\textbf{MAPE} je povprečje $\mathrm{APE}_i$, \textbf{MdAPE} pa njihova mediana;
MdAPE je robustnejša na osamelce in je naša glavna mera~\cite{hyndman2006}.
Poročamo tudi \textbf{MAE} (povprečna absolutna napaka) in \textbf{RMSE}
(koren povprečne kvadratne napake) v evrih ter \emph{predznačeno} odstotno
napako, ki pokaže sistematičen odmik ocene navzgor ali navzdol.

\paragraph{Pokritost.}
\textbf{Pokritost} je delež poizvedb, za katere sistem vrne oceno, torej za
katere najde vsaj eno (oziroma vsaj $5$ oz.\ $10$) primerljivo vozilo.

\paragraph{Raznolikost ponudbe.}
\textbf{Raznolikost} (v izvorni kodi \texttt{ratio}) je relativna širina
cenovnega razpona primerljivih vozil,
\begin{equation}
  \label{eq:raznolikost}
  \mathrm{raznolikost} = \frac{P_{75} - P_{25}}{\hat{p}},
\end{equation}
kjer sta $P_{25}$ in $P_{75}$ (v kodi \texttt{goodPrice} in \texttt{badPrice})
uteženi $25.$ oziroma $75.$ percentil cen primerljivih vozil, $\hat{p}$
(v kodi \texttt{estimatedPrice}) pa ocenjena cena vozila. Vrednost izražamo v
odstotkih; enoten zapis uporabljamo v vseh tabelah tega poglavja.

% =====================================================================
\section{Značilnosti zbranih podatkov}
\label{sec:podatki}
\paragraph{Viri in obseg.}
Velika večina oglasov je bila zajeta z največjega slovenskega spletnega portala
za prodajo rabljenih vozil (razdelek~\ref{sec:zajem-podatkov}); manjši del
predstavlja vzorec z večjega nemškega oglasnega portala, ki smo ga uporabili za
preverjanje delovanja cevovoda na drugačni strukturi vira. Zajem je zajel
javno dostopne oglase za nekomercialno raziskovalno rabo, z omejeno hitrostjo
zahtevkov in brez shranjevanja kontaktnih ali osebnih podatkov prodajalcev.
Po normalizaciji baza obsega $85.846$ oglasov ($83.493$ s slovenskega in
$2.353$ z nemškega portala) za $10$ znamk vozil; razčlenitev po znamki je v
tabeli~\ref{tab:podatki-znamke}.

\begin{table}[htbp]
  \centering
  \caption{Število normaliziranih oglasov po znamki.}
  \label{tab:podatki-znamke}
  % vir: analysis/out/dataset_per_brand.csv
  \begin{tabular}{lr lr}
    \toprule
    \textbf{Znamka} & \textbf{Oglasi} & \textbf{Znamka} & \textbf{Oglasi} \\
    \midrule
    Volkswagen    & 21.191 & Peugeot     & 7.513 \\
    Audi          & 11.842 & Citroën     & 6.729 \\
    BMW           & 10.724 & Ford        & 5.783 \\
    Renault       & 9.615  & Toyota      & 2.648 \\
    Mercedes-Benz & 9.355  & Alfa Romeo  & 446 \\
    \bottomrule
  \end{tabular}
\end{table}

\paragraph{Porazdelitve ključnih atributov.}
Slika~\ref{fig:porazdelitve} prikazuje porazdelitve oglaševanih cen, prevožene
kilometrine in letnika. Porazdelitev cen je izrazito desno asimetrična (dolg
rep dragih vozil), zato jo prikazujemo na logaritemski osi. Mediana cene je
$14.890$~EUR (povprečje $18.137$~EUR), mediana kilometrine $159.950$~km,
mediana letnika pa $2017$; podrobnosti so v tabeli~\ref{tab:podatki-porazdelitve}.
Zaradi berljivosti so na grafih odrezane skrajne vrednosti (nekaj oglasov ima
očitno napačno ceno, kilometrino ali letnik); njihovo število je označeno v
kotu vsakega grafa, celoten razpon pa je razviden iz stolpca \emph{maks.}\ v
tabeli~\ref{tab:podatki-porazdelitve} in iz obravnave osamelcev v
tabeli~\ref{tab:osamelci}.

\begin{table}[htbp]
  \centering
  \caption{Povzetek porazdelitev ključnih numeričnih atributov.}
  \label{tab:podatki-porazdelitve}
  % vir: analysis/out/dataset_distribution_summary.json
  \setlength{\tabcolsep}{5pt}
  \begin{tabular}{lrrrrr}
    \toprule
    & \textbf{$p_{5}$} & \textbf{mediana} & \textbf{povprečje} & \textbf{$p_{95}$} & \textbf{maks.} \\
    \midrule
    cena [EUR]        & 1.550  & 14.890  & 18.137  & 49.014  & 374.000 \\
    kilometrina [km]  & 26.000 & 159.950 & 169.464 & 334.000 & 1.974.001 \\
    letnik            & 2005   & 2017    & 2016    & 2024    & 2029 \\
    moč [kW]          & 56     & 110     & 117     & 221     & 920 \\
    \bottomrule
  \end{tabular}
\end{table}

\begin{figure}[htbp]
  \centering
  \begin{subfigure}{0.48\textwidth}
    \includegraphics[width=\textwidth]{dataset_price_hist.pdf}
    \caption{oglaševana cena (log)}
  \end{subfigure}\hfill
  \begin{subfigure}{0.48\textwidth}
    \includegraphics[width=\textwidth]{dataset_km_hist.pdf}
    \caption{prevoženi kilometri}
  \end{subfigure}
  \\[1ex]
  \begin{subfigure}{0.48\textwidth}
    \includegraphics[width=\textwidth]{dataset_year_hist.pdf}
    \caption{letnik prve registracije}
  \end{subfigure}\hfill
  \begin{subfigure}{0.48\textwidth}
    \includegraphics[width=\textwidth]{dataset_brand_bar.pdf}
    \caption{oglasi po znamki}
  \end{subfigure}
  \caption{Porazdelitve ključnih atributov v normalizirani bazi.}
  \label{fig:porazdelitve}
  % vir: analysis/out/dataset_*_hist.pdf, dataset_brand_bar.pdf
\end{figure}

\paragraph{Pokritost po segmentih.}
Za pošteno oceno mora imeti algoritem dovolj primerljivih vozil v ožjem
segmentu (ista znamka in model, podoben letnik in kilometrina).
Slika~\ref{fig:pokritost} prikazuje toplotno karto za znamki Renault in
Volkswagen: vrstice so modeli, stolpci letniki, barva pa
število oglasov v tej kombinaciji. Vidno je, da so pogosti modeli srednjih
letnikov (npr.\ Volkswagen Golf ali Renault Clio) gosto pokriti, redki modeli
in skrajni letniki pa le redko, kar neposredno pojasni razlike v zanesljivosti
ocen (razdelek~\ref{sec:natancnost}).

\begin{figure}[htbp]
  \centering
  \includegraphics[width=\textwidth]{dataset_coverage_heatmap.pdf}
  \caption{Pokritost: število oglasov po modelu in letniku za znamki
  Renault in Volkswagen (po ena znamka na vrstico).}
  \label{fig:pokritost}
  % vir: analysis/out/dataset_coverage_heatmap.pdf
\end{figure}

\paragraph{Osamelci in kakovost podatkov.}
Postopek normalizacije trenutno ne izvaja filtriranja osamelcev, zato baza
vsebuje tudi vozila z očitno napačnimi vrednostmi.
Tabela~\ref{tab:osamelci} prikazuje njihov delež. Kar $10{,}7\,\%$ oglasov ima
vsaj eno vrednost zunaj smiselnih mej, največ pri kilometrini ($9{,}9\,\%$;
najvišja zabeležena vrednost je skoraj dva milijona kilometrov) in letniku
($8{,}1\,\%$; nekaj oglasov ima celo letnik v prihodnosti). Preprost filter po
interkvartilnem razponu (IQR) na ceni znotraj skupine
$(\text{znamka},\text{model})$ bi odstranil dodatnih $2{,}2\,\%$ oglasov.
Vpliv doslednega filtriranja osamelcev na natančnost obravnavamo v
razdelku~\ref{sec:obcutljivost}, uvedbo takega filtra pa v
poglavju~\ref{ch:sklep}.

\begin{table}[H]
  \centering
  \caption{Osamelci v normalizirani bazi ($85.846$ oglasov).}
  \label{tab:osamelci}
  % vir: analysis/out/dataset_outliers.csv
  \begin{tabularx}{\textwidth}{@{}X r r@{}}
    \toprule
    \textbf{Merilo} & \textbf{Št.\ vrstic} & \textbf{Delež} \\
    \midrule
    cena zunaj $[500,\ 200.000]$~EUR         & 427   & $0{,}5\,\%$ \\
    kilometrina zunaj $[1000,\ 500.000]$~km  & 8.495 & $9{,}9\,\%$ \\
    letnik zunaj $[1990,\ 2027]$             & 6.979 & $8{,}1\,\%$ \\
    vsaj ena vrednost zunaj mej              & 9.169 & $10{,}7\,\%$ \\
    dodatno po filtru IQR na ceni            & 1.878 & $2{,}2\,\%$ \\
    \bottomrule
  \end{tabularx}
\end{table}

% =====================================================================
\section{Natančnost ocene vrednosti}
\label{sec:natancnost}

To je osrednja meritev diplomskega dela: neposredno preverja zastavljeno
izhodišče o \emph{dovolj natančni} oceni (RV1). Protokol je opisan v
razdelku~\ref{sec:metodologija}, rezultate pa izdela skripta
\texttt{analysis/accuracy\_cv.py} nad vzorcem $3000$ oglasov (od tega jih je
$2986$ dobilo oceno).

\paragraph{Merilo resnice in namerni odmik ocene.}
Kot merilo resnice uporabljamo \emph{oglaševano} ceno, saj dejanske prodajne
cene niso javno dostopne. Sistem pa namenoma ne ocenjuje oglaševane cene,
temveč \emph{pričakovano prodajno ceno}: ker se rabljena vozila v pogajanjih
praviloma prodajo pod oglaševano ceno, oceno zniža s stopenjskim popustom
($9$--$15\,\%$ glede na cenovni razred). Zato je pri primerjavi z oglaševano
ceno pričakovan rahel \emph{negativni} odmik ocene, kar je zaželeno in ne
pomeni slabše ocene. To lastnost upoštevamo pri branju vseh rezultatov v
nadaljevanju: napako merimo dvakrat --- proti oglaševani ceni z vključenim
popustom (kar dejansko vrne sistem) in brez popusta (čista primerjalna
ocena).

\paragraph{Skupni rezultat.}
Brez popusta znaša MdAPE predlaganega pristopa $13{,}5\,\%$, MAPE $23{,}8\,\%$,
MAE $2798$~EUR in RMSE $5151$~EUR; z vključenim namernim popustom je MdAPE
proti oglaševani ceni $14{,}5\,\%$ (mediana predznačene napake $-6{,}8\,\%$).
Pokritost je $99{,}5\,\%$ pri $n \ge 1$ in $95{,}7\,\%$ pri $n \ge 10$.
Slika~\ref{fig:pred-vs-actual} prikazuje napovedano ceno proti dejanski
($R^2 = 0{,}87$), slika~\ref{fig:ape-hist} pa porazdelitev odstotne napake.

\begin{figure}[htbp]
  \centering
  \begin{subfigure}{0.48\textwidth}
    \includegraphics[width=\textwidth]{accuracy_pred_vs_actual.pdf}
    \caption{napovedano proti dejanskemu}
    \label{fig:pred-vs-actual}
  \end{subfigure}\hfill
  \begin{subfigure}{0.48\textwidth}
    \includegraphics[width=\textwidth]{accuracy_ape_hist.pdf}
    \caption{porazdelitev napake}
    \label{fig:ape-hist}
  \end{subfigure}
  \caption{Natančnost ocene vrednosti (leave-one-out, brez popusta).}
  % vir: analysis/out/accuracy_pred_vs_actual.pdf, accuracy_ape_hist.pdf
\end{figure}

\paragraph{Primerjava z osnovnimi modeli.}
Da presodimo, ali je predlagani pristop boljši od trivialnega, ga primerjamo s
štirimi referenčnimi modeli na istem protokolu in isti množici kandidatov
(tabela~\ref{tab:baselines}, slika~\ref{fig:baselines}): B0~globalna mediana
cene za dani model, B1~navadno povprečje cen v filtrirnem oknu, B1m~mediana cen
v oknu, B2~linearna regresija $\text{cena} \sim \text{letnik} + \text{km}$,
B3~metoda $k$~najbližjih sosedov ($k = 10$) po standardiziranih letniku in
kilometrini, ter B4~naš pristop (uteženo povprečje s podobnostnim jedrom, brez
popusta).

\begin{table}[htbp]
  \centering
  \caption{Primerjava natančnosti z osnovnimi modeli (leave-one-out).
  MdAPE in MAPE sta v odstotkih, MAE in RMSE v evrih.}
  \label{tab:baselines}
  % vir: analysis/out/accuracy_overall.csv
  \setlength{\tabcolsep}{5pt}
  \begin{tabular}{l r r r r}
    \toprule
    \textbf{Model} & \textbf{MdAPE} & \textbf{MAPE} & \textbf{MAE} & \textbf{RMSE} \\
    \midrule
    B0 globalna mediana       & 37{,}8 & 95{,}7 & 7395 & 10883 \\
    B1 povprečje v oknu       & 14{,}2 & 24{,}9 & 2930 & 5333 \\
    B1m mediana v oknu        & 13{,}1 & 22{,}5 & 2877 & 5561 \\
    B2 linearna regresija     & 12{,}9 & 22{,}4 & 2677 & 5032 \\
    B3 $k$-najbližjih sosedov & \textbf{11{,}2} & 21{,}5 & 2497 & 4883 \\
    B4 naš pristop            & 13{,}5 & 23{,}8 & 2798 & 5151 \\
    \bottomrule
  \end{tabular}
\end{table}

\begin{figure}[htbp]
  \centering
  \includegraphics[width=0.7\textwidth]{accuracy_baselines.pdf}
  \caption{MdAPE osnovnih modelov in predlaganega pristopa.}
  \label{fig:baselines}
  % vir: analysis/out/accuracy_baselines.pdf
\end{figure}

Predlagani pristop je \emph{bistveno} boljši od naivne globalne mediane
($13{,}5\,\%$ proti $37{,}8\,\%$) in primerljiv s preostalimi metodami, ki se
gibljejo med $11$ in $14\,\%$. Ni pa najboljši: metoda $k$~najbližjih sosedov
doseže $11{,}2\,\%$, tudi navadna mediana istega okna ($13{,}1\,\%$) je
malenkost natančnejša. To je pričakovano za zasnovo, katere cilj je
\emph{razložljivost}: za razliko od naštetih metod predlagani pristop poleg
točkovne ocene vrne tudi cenovni razpon in indeks zaupanja, uteži pa so
neposredno interpretabilne. Rezultat hkrati potrjuje smiselnost prehoda na
metode strojnega učenja, ki jo predlagamo v poglavju~\ref{ch:sklep}: če že
preprost $k$-NN preseže sedanji pristop, je od zmogljivejših modelov mogoče
pričakovati še boljše rezultate.

\paragraph{Odvisnost od velikosti vzorca in segmenta (RV2).}
Slika~\ref{fig:mdape-n} prikazuje MdAPE glede na število primerljivih vozil
$n$: napaka je pri $n < 10$ okoli $20\,\%$, pri $n \ge 20$ pa se ustali pri
približno $13\,\%$. To empirično podpira izbiro konstante $20$ v indeksu
zaupanja (enačba~\ref{eq:confidence}) --- pri toliko primerljivih vozilih je
ocena že blizu svoje najboljše natančnosti. Slika~\ref{fig:mdape-brand} prikazuje MdAPE po znamki: najbolje se
odrežejo pogosti, gosto pokriti modeli (Audi $9{,}8\,\%$, Toyota $10{,}8\,\%$,
Volkswagen $12{,}1\,\%$), najslabše pa Mercedes-Benz ($24{,}9\,\%$), pri
katerem se pri normalizaciji več različnih modelov zlije v isto kanonično
ime, kar poslabša izbor primerljivih vozil.
Tabela~\ref{tab:natancnost-segmenti} kaže še razčlenitev po starosti,
kilometrini in raznolikosti ponudbe.

\begin{figure}[htbp]
  \centering
  \begin{subfigure}{0.48\textwidth}
    \includegraphics[width=\textwidth]{accuracy_mdape_vs_n.pdf}
    \caption{glede na velikost vzorca $n$}
    \label{fig:mdape-n}
  \end{subfigure}\hfill
  \begin{subfigure}{0.48\textwidth}
    \includegraphics[width=\textwidth]{accuracy_mdape_by_brand.pdf}
    \caption{po znamki}
    \label{fig:mdape-brand}
  \end{subfigure}
  \caption{Odvisnost natančnosti ocene od podatkov.}
  % vir: analysis/out/accuracy_mdape_vs_n.pdf, accuracy_mdape_by_brand.pdf
\end{figure}

\begin{table}[htbp]
  \centering
  \caption{Natančnost predlaganega pristopa po segmentih (MdAPE).}
  \label{tab:natancnost-segmenti}
  % vir: analysis/out/accuracy_by_age.csv, accuracy_by_km.csv, accuracy_by_ratio.csv
  \begin{tabularx}{\textwidth}{@{}l X r r@{}}
    \toprule
    \textbf{Razčlenitev} & \textbf{Koš} & \textbf{Primeri} & \textbf{MdAPE [\%]} \\
    \midrule
    starost [leta] & 0--3   & 367 & 9{,}0 \\
                   & 4--7   & 921 & 10{,}6 \\
                   & 8--12  & 932 & 14{,}5 \\
                   & $13+$  & 766 & 22{,}9 \\
    \midrule
    kilometrina    & $<50.000$        & 280  & 9{,}6 \\
                   & $50.000$--$100.000$  & 449  & 11{,}9 \\
                   & $100.000$--$150.000$ & 634  & 12{,}4 \\
                   & $150.000$--$200.000$ & 629  & 13{,}7 \\
                   & $>200.000$       & 1021 & 16{,}7 \\
    \midrule
    raznolikost    & $0$--$0{,}10$    & 88   & 7{,}5 \\
                   & $0{,}10$--$0{,}25$ & 1002 & 8{,}4 \\
                   & $0{,}25$--$0{,}50$ & 1613 & 15{,}5 \\
                   & $>0{,}50$        & 447  & 29{,}7 \\
    \bottomrule
  \end{tabularx}
\end{table}

Vzorec pokaže jasne zakonitosti: napaka narašča s starostjo vozila ($9\,\%$ pri
najnovejših, $23\,\%$ pri vozilih, starejših od $13$~let), s prevoženo
kilometrino in --- najmočneje --- z \emph{raznolikostjo ponudbe}. Pri modelih z
enotno ponudbo (raznolikost pod $0{,}1$) je MdAPE le $7{,}5\,\%$, pri razpršeni
ponudbi (nad $0{,}5$) pa $29{,}7\,\%$. To potrjuje, da je raznolikost, ki jo
sistem uporabniku tudi prikaže, dober vgrajen pokazatelj zanesljivosti
posamezne ocene.

\paragraph{Vpliv popusta.}
Slika~\ref{fig:popust} prikazuje predznačeno napako z in brez popusta. Brez
popusta je ocena rahlo optimistična (mediana $+3{,}8\,\%$ nad oglaševano ceno),
z namernim popustom pa se premakne na $-6{,}8\,\%$ pod oglaševano ceno. Ta
premik je skladen z zasnovo sistema, ki cilja pričakovano prodajno in ne
oglaševane cene; dokončno umerjanje popusta bi bilo mogoče šele s podatki o
dejanskih prodajah (razdelek~\ref{sec:groznje}).

\begin{figure}[htbp]
  \centering
  \includegraphics[width=0.7\textwidth]{accuracy_discount_bias.pdf}
  \caption{Predznačena napaka proti oglaševani ceni, z in brez namernega popusta.}
  \label{fig:popust}
  % vir: analysis/out/accuracy_discount_bias.pdf
\end{figure}

\paragraph{Časovni razrez.}
Pri napovedovanju najnovejših $10\,\%$ oglasov zgolj iz starejših $90\,\%$
znaša MdAPE predlaganega pristopa $13{,}0\,\%$ (tabela~\ref{tab:temporal}), kar
je celo malenkost bolje od rezultata leave-one-out in potrjuje, da ocena ne
izhaja iz uhajanja informacije.

\begin{table}[htbp]
  \centering
  \caption{Natančnost pri časovnem razrezu (napoved najnovejših $10\,\%$).
  MdAPE in MAPE sta v odstotkih, MAE in RMSE v evrih.}
  \label{tab:temporal}
  % vir: analysis/out/accuracy_temporal.csv
  \setlength{\tabcolsep}{5pt}
  \begin{tabular}{l r r r r}
    \toprule
    \textbf{Model} & \textbf{MdAPE} & \textbf{MAPE} & \textbf{MAE} & \textbf{RMSE} \\
    \midrule
    B0 globalna mediana       & 35{,}9 & 74{,}5 & 7596 & 11222 \\
    B1m mediana v oknu        & 12{,}8 & 19{,}7 & 3088 & 5805 \\
    B3 $k$-najbližjih sosedov & 10{,}5 & 18{,}2 & 2667 & 4890 \\
    B4 naš pristop            & 13{,}0 & 20{,}6 & 2994 & 5347 \\
    \bottomrule
  \end{tabular}
\end{table}

% =====================================================================
\section{Občutljivost algoritma na parametre}
\label{sec:obcutljivost}

Ocenjevalni algoritem vsebuje več hevristično določenih konstant: dovoljeno
odstopanje kilometrine ($\pm 40.000$~km) in letnika ($\pm 3$~leta), uteži
delnih ocen ($0{,}6$ in $0{,}4$), oblika padajoče funkcije in stopenjski
popust. V tem razdelku vsako spreminjamo ob nespremenjenih ostalih in merimo
MdAPE ter pokritost (vzorec $2000$ oglasov, leave-one-out). Rezultate izdela
skripta \texttt{analysis/sensitivity.py}.

\paragraph{Okno primerljivosti.}
Tabela~\ref{tab:sens-okno} in slika~\ref{fig:sens-km} kažeta, da je algoritem
na širino \emph{kilometrinskega} okna skoraj neobčutljiv: MdAPE se med
$\pm 10.000$~km in $\pm 80.000$~km giblje le med $13{,}2$ in $13{,}6\,\%$.
Nasprotno je odvisen od \emph{letnikovega} okna: pri $\pm 1$~letu je MdAPE
$12{,}6\,\%$, pri $\pm 5$~letih pa $15{,}2\,\%$. Ožje okno da bolj primerljiva
vozila, a manjši vzorec (nižja pokritost); gre za znano razmerje med
pristranskostjo in varianco~\cite{bias-variance}. Toplotna karta
(slika~\ref{fig:sens-heatmap}) potrjuje, da natančnost določa skoraj izključno
letnikovo okno. Sedanja izbira $\pm 3$~leta je razumen kompromis med
natančnostjo in pokritostjo; morebitna zožitev na $\pm 2$~leti bi MdAPE
znižala za približno $0{,}3$~odstotne točke ob še vedno visoki pokritosti.

\begin{table}[htbp]
  \centering
  \caption{Občutljivost na širino okna primerljivosti. Pokritost je delež
  poizvedb z vsaj $5$ primerljivimi vozili.}
  \label{tab:sens-okno}
  % vir: analysis/out/sensitivity_km.csv, sensitivity_year.csv
  \setlength{\tabcolsep}{5pt}
  \begin{tabular}{l r r r}
    \toprule
    \textbf{Parameter} & \textbf{MdAPE [\%]} & \textbf{Pokritost [\%]} & \textbf{Mediana $n$} \\
    \midrule
    kilometrina $\pm 10.000$~km & 13{,}2 & 92{,}3 & 53 \\
    kilometrina $\pm 40.000$~km & 13{,}4 & 98{,}1 & 196 \\
    kilometrina $\pm 80.000$~km & 13{,}6 & 99{,}0 & 362 \\
    \midrule
    letnik $\pm 1$~leto  & 12{,}6 & 95{,}6 & 97 \\
    letnik $\pm 3$~leta  & 13{,}4 & 98{,}1 & 196 \\
    letnik $\pm 5$~let   & 15{,}2 & 98{,}5 & 262 \\
    \bottomrule
  \end{tabular}
\end{table}

\begin{figure}[htbp]
  \centering
  \begin{subfigure}{0.48\textwidth}
    \includegraphics[width=\textwidth]{sensitivity_km.pdf}
    \caption{okno kilometrine}
  \end{subfigure}\hfill
  \begin{subfigure}{0.48\textwidth}
    \includegraphics[width=\textwidth]{sensitivity_year.pdf}
    \caption{okno letnika}
  \end{subfigure}
  \caption{Občutljivost na širino okna primerljivosti (MdAPE in pokritost).}
  \label{fig:sens-km}
  % vir: analysis/out/sensitivity_km.pdf, sensitivity_year.pdf
\end{figure}

\begin{figure}[htbp]
  \centering
  \includegraphics[width=0.7\textwidth]{sensitivity_km_year_heatmap.pdf}
  \caption{MdAPE [\%] po mreži dovoljenih odstopanj kilometrine in letnika.}
  \label{fig:sens-heatmap}
  % vir: analysis/out/sensitivity_km_year_heatmap.pdf
\end{figure}

\paragraph{Uteži, oblika jedra in popust.}
Tabela~\ref{tab:sens-ostalo} povzema preostale parametre. Razmerje uteži med
kilometrino in letnikom na natančnost praktično ne vpliva ($13{,}4$ do
$13{,}6\,\%$ za vsa preizkušena razmerja), prav tako ne zamenjava linearne
padajoče funkcije z eksponentno ($13{,}4$ proti $13{,}3\,\%$). Edina parametra,
ki opazno premakneta rezultat, sta torej širina letnikovega okna in namerni
popust: ta predznačeno napako proti oglaševani ceni premakne iz $+3{,}8\,\%$ v
$-6{,}8\,\%$, MdAPE pa zviša z $13{,}4$ na $14{,}5\,\%$. Ta premik je posledica
zavestne odločitve, da sistem ocenjuje pričakovano prodajno in ne oglaševane
cene, in ne slabše ocene.

\begin{table}[htbp]
  \centering
  \caption{Občutljivost na uteži, obliko jedra in popust.}
  \label{tab:sens-ostalo}
  % vir: analysis/out/sensitivity_weights.csv, sensitivity_kernel.csv, sensitivity_discount.csv
  \begin{tabularx}{\textwidth}{@{}l X r >{\raggedleft\arraybackslash}p{2.3cm}@{}}
    \toprule
    \textbf{Parameter} & \textbf{Nastavitev} & \textbf{MdAPE [\%]} & \textbf{Predznačena napaka [\%]} \\
    \midrule
    uteži (km / letnik) & $0{,}5 / 0{,}5$              & 13{,}4 & $+3{,}7$ \\
    uteži (km / letnik) & $0{,}6 / 0{,}4$ (privzeto)   & 13{,}4 & $+3{,}8$ \\
    uteži (km / letnik) & $0{,}8 / 0{,}2$              & 13{,}6 & $+3{,}8$ \\
    \midrule
    jedro               & linearno (privzeto)         & 13{,}4 & $+3{,}8$ \\
    jedro               & eksponentno                 & 13{,}3 & $+3{,}8$ \\
    \midrule
    popust              & brez                        & 13{,}4 & $+3{,}8$ \\
    popust              & stopenjski (privzeto)        & 14{,}5 & $-6{,}8$ \\
    \bottomrule
  \end{tabularx}
\end{table}

\paragraph{Sklep o robustnosti.}
Algoritem je robusten: na večino svojih konstant se skoraj ne odziva, kar
pomeni, da njihova hevristična izbira ni sporna. Njegova natančnost je odvisna
predvsem od kakovosti in gostote podatkov (razdelek~\ref{sec:natancnost}), ne
od finega nastavljanja parametrov.

% =====================================================================
\section{Latenca posameznega vrednotenja}
\label{sec:latenca}

Merimo latenco vrste (a) in (b) iz razdelka~\ref{sec:mere}. Meritev izvede
\texttt{benchmarks/bench\_algo.py}: $2000$ neposrednih klicev funkcije nad
naključnimi realnimi vhodi (brez HTTP), po $100$ ogrevalnih klicih, ki se ne
merijo. Ločeno merimo čas poizvedbe SQL, čas točkovanja v Pythonu in čas
serializacije odgovora.

\paragraph{Ozko grlo: pregled celotne tabele.}
Prvotna izvedba poizvedbe ni imela ustreznega indeksa.
Izsek~\ref{lst:explain-before} prikazuje izpis \texttt{EXPLAIN (ANALYZE,
BUFFERS)}: načrt je \texttt{Seq Scan}, baza pri vsakem vrednotenju prebere
celotno tabelo in zavrže $85.273$ vrstic, ki ne ustrezajo pogojem
(\texttt{Rows Removed by Filter}). Ker so podatki v celoti v predpomnilniku
(\texttt{shared hit}, brez \texttt{read}), je čas stabilen, a raste linearno z
velikostjo tabele.

\begin{center}
\captionof{listing}{\texttt{EXPLAIN (ANALYZE, BUFFERS)} pred dodajanjem indeksa.}
\label{lst:explain-before}
\begin{minted}{text}
Seq Scan on normalized_market_listing
    (cost=0.00..3314.15 rows=12 width=97)
    (actual time=0.555..6.847 rows=573 loops=1)
  Filter: ((kilometers >= 18500) AND (kilometers <= 98500)
       AND (registration_year >= 2019) AND (registration_year <= 2025)
       AND (brand_name = 'Citroën') AND (model_name = 'C3'))
  Rows Removed by Filter: 85273
  Buffers: shared hit=1168
Execution Time: 6.880 ms
\end{minted}
\end{center}

\paragraph{Učinek sestavljenega indeksa.}
Po dodajanju sestavljenega indeksa na
$(\text{brand\_name},\ \text{model\_name},\ \text{registration\_year},\
\text{kilometers})$ se načrt spremeni v \texttt{Bitmap Index Scan}
(izsek~\ref{lst:explain-after}): baza prebere le nekaj strani indeksa in
pripadajoče vrstice, čas izvajanja pa pade z $6{,}88$~ms na $0{,}69$~ms.
Latenca poizvedbe SQL se v mediani zmanjša z $9{,}7$~ms na $2{,}0$~ms
(tabela~\ref{tab:indeks-pred-po}, slika~\ref{fig:indeks-pred-po}). Pomembneje
od trenutnega pohitrenja je, da postane čas poizvedbe neodvisen od velikosti
tabele: namesto $O(\text{tabela})$ postane $O(\text{zadetki})$.

\begin{table}[htbp]
  \centering
  \caption{Latenca poizvedbe SQL pred in po dodajanju indeksa, v milisekundah.}
  \label{tab:indeks-pred-po}
  % vir: benchmarks/report.md (results_algo_before/after.json)
  \begin{tabular}{lrrrr}
    \toprule
    & \textbf{p50} & \textbf{p95} & \textbf{p99} & \textbf{Povprečje} \\
    \midrule
    pred (Seq Scan)  & 9{,}7 & 21{,}0 & 24{,}8 & 11{,}3 \\
    po (Index Scan)  & 2{,}0 & 8{,}6  & 12{,}5 & 3{,}0 \\
    \bottomrule
  \end{tabular}
\end{table}

\begin{figure}[htbp]
  \centering
  \includegraphics[width=0.6\textwidth]{latency_index_before_after.pdf}
  \caption{Latenca poizvedbe SQL pred in po dodajanju sestavljenega indeksa.}
  \label{fig:indeks-pred-po}
  % vir: analysis/out/latency_index_before_after.pdf
\end{figure}

\begin{center}
\captionof{listing}{\texttt{EXPLAIN (ANALYZE, BUFFERS)} po dodajanju indeksa.}
\label{lst:explain-after}
\begin{minted}{text}
Bitmap Heap Scan on normalized_market_listing
    (cost=5.84..50.49 rows=12 width=97)
    (actual time=0.202..0.615 rows=573 loops=1)
  Recheck Cond: (brand_name = 'Citroën' AND model_name = 'C3'
       AND registration_year BETWEEN 2019 AND 2025
       AND kilometers BETWEEN 18500 AND 98500)
  Buffers: shared hit=195
  ->  Bitmap Index Scan on idx_nml_valuation
        (actual time=0.166..0.166 rows=573 loops=1)
        Buffers: shared hit=6
Execution Time: 0.688 ms
\end{minted}
\end{center}

\paragraph{Razčlenitev latence po komponentah.}
Tabela~\ref{tab:latenca-vrednotenja} in slika~\ref{fig:latency-breakdown}
prikazujeta latenco celotnega vrednotenja po dodajanju indeksa. Tipično
vrednotenje traja $2{,}0$~ms, v najslabšem odstotku primerov (p99) pa
$13{,}5$~ms. Skoraj celotna latenca ($\approx 88\,\%$ povprečja) odpade na
poizvedbo SQL; točkovanje v Pythonu ($0{,}4$~ms povprečno) in serializacija
odgovora ($\approx 0{,}01$~ms) sta zanemarljiva. Poizvedba v povprečju vrne
$287$ primerljivih vozil (mediana $175$, največ $1495$).

\begin{table}[htbp]
  \centering
  \caption{Latenca posameznega vrednotenja po dodajanju indeksa
  ($N = 2000$), v milisekundah.}
  \label{tab:latenca-vrednotenja}
  % vir: benchmarks/report.md (results_algo.json)
  \begin{tabular}{lrrrr}
    \toprule
    \textbf{Komponenta} & \textbf{p50} & \textbf{p95} & \textbf{p99} & \textbf{Povprečje} \\
    \midrule
    Skupaj                & 2{,}0  & 9{,}7 & 13{,}5 & 3{,}1 \\
    Poizvedba SQL         & 1{,}8  & 8{,}4 & 11{,}5 & 2{,}7 \\
    Točkovanje (Python)   & 0{,}2  & 1{,}2 & 1{,}8  & 0{,}4 \\
    Serializacija (JSON)  & 0{,}01 & 0{,}02 & 0{,}03 & 0{,}01 \\
    \bottomrule
  \end{tabular}
\end{table}

\begin{figure}[htbp]
  \centering
  \includegraphics[width=0.6\textwidth]{latency_breakdown.pdf}
  \caption{Razčlenitev latence posameznega vrednotenja po komponentah.}
  \label{fig:latency-breakdown}
  % vir: analysis/out/latency_breakdown.pdf
\end{figure}

Vrednotenje se torej izvede daleč pod nefunkcionalno zahtevo $2$~s
(razdelek~\ref{sec:zahteve}); prostora za rast podatkovne baze je ob
prisotnosti indeksa dovolj.

% =====================================================================
\section{Obnašanje pod obremenitvijo}
\label{sec:obremenitev}

Merimo end-to-end latenco (vrsta (c)) in prepustnost končne točke
\texttt{POST /valuate} pri naraščajoči sočasnosti. Obremenitev ustvarja
skripta \texttt{benchmarks/loadtest.py} z več vzporednimi odjemalci, ki v
fiksnem $20$-sekundnem oknu neprekinjeno pošiljajo zahtevke z naključnimi
realnimi vhodi. Test se samodejno prekine ob deležu napak nad $5\,\%$ ali
$p_{95}$ nad $5$~s. Med testom vzorčimo tudi zasedenost baze
(\texttt{pg\_stat\_activity}).

\begin{table}[htbp]
  \centering
  \caption{Prepustnost in latenca pri naraščajoči sočasnosti.
  Latence so v milisekundah, $L$ je ocena po Littlovem zakonu.}
  \label{tab:obremenitev}
  % vir: benchmarks/results_load.json (loadtest.py --db-sample)
  \begin{tabular}{r r r r r r r}
    \toprule
    \textbf{Sočasnost} & \textbf{Prep.\ [z/s]} & \textbf{p50} & \textbf{p95} & \textbf{p99}
      & \textbf{Napake} & \textbf{$L$} \\
    \midrule
    1  & \res{x} & \res{x} & \res{x} & \res{x} & \res{x} & \res{x} \\
    5  & \res{x} & \res{x} & \res{x} & \res{x} & \res{x} & \res{x} \\
    10 & \res{x} & \res{x} & \res{x} & \res{x} & \res{x} & \res{x} \\
    20 & \res{x} & \res{x} & \res{x} & \res{x} & \res{x} & \res{x} \\
    40 & \res{x} & \res{x} & \res{x} & \res{x} & \res{x} & \res{x} \\
    \bottomrule
  \end{tabular}
\end{table}

\paragraph{Rezultat.}
\res{Prepustnost narašča približno linearno do sočasnosti X, nato se ustali pri
približno X zahtevkih na sekundo. Vse latence ostanejo pod nefunkcionalno
zahtevo 2~s.}

\paragraph{Skladnost meritve (Littlov zakon).}
\res{Stolpec $L$ (prepustnost krat povprečna latenca) ostaja blizu nastavljene
sočasnosti, kar potrjuje notranjo skladnost meritve.}

\paragraph{Dokaz ozkega grla.}
Teoretični strop postavitve je $1000 / S \approx 324$ vrednotenj na sekundo,
kjer je $S = 3{,}1$~ms povprečni čas ene evalvacije. Dejanska meja je nižja,
ker ves promet serializira \emph{ena sama deljena povezava do baze}
(\texttt{app/database.py}) in en sam strežniški proces.
\res{To potrjuje vzorčenje baze med testom: število aktivnih poizvedb kljub
naraščajoči sočasnosti odjemalcev ne preseže X.}

% =====================================================================
\section{Preverjanje funkcionalnih zahtev}
\label{sec:funkcionalne}

Tabela~\ref{tab:sledljivost} povezuje vsako funkcionalno zahtevo iz
razdelka~\ref{sec:zahteve} z načinom preverjanja in izidom.

\begin{table}[htbp]
  \centering
  \caption{Sledljivost funkcionalnih zahtev do preverjanja.}
  \label{tab:sledljivost}
  % vir: ročni scenariji + imena Jest testov v car-stock-market-be
  \begin{tabularx}{\textwidth}{@{}X X c@{}}
    \toprule
    \textbf{Zahteva} & \textbf{Način preverjanja} & \textbf{Izid} \\
    \midrule
    Registracija uporabnika              & \res{test / scenarij} & \res{OK} \\
    Prijava in odjava                    & \res{test / scenarij} & \res{OK} \\
    Objava oglasa (le registrirani)      & \res{test / scenarij} & \res{OK} \\
    Vrednotenje vozila po atributih      & \res{test / scenarij} & \res{OK} \\
    Všečkanje oglasov                    & \res{test / scenarij} & \res{OK} \\
    Ažuriranje lastnih/všečkanih oglasov & \res{test / scenarij} & \res{OK} \\
    Pregled oglasov (anonimni)           & \res{test / scenarij} & \res{OK} \\
    \bottomrule
  \end{tabularx}
\end{table}

\paragraph{Pokritost testov.}
Avtomatizirani testi zalednega sistema (Jest + Supertest,
razdelek~\ref{sec:testiranje-be}) obsegajo \res{X} testov v \res{X} zbirkah in
pokrivajo \res{X}\,\% stavkov ter \res{X}\,\% vej. Testi se izvedejo ob vsakem
zahtevku za združitev na cevovodu neprekinjene integracije
(razdelek~\ref{subsec:ci}).

% =====================================================================
\section{Preverjanje nefunkcionalnih zahtev}
\label{sec:nefunkcionalne}

Tabela~\ref{tab:nfr} za vsako nefunkcionalno zahtevo iz
razdelka~\ref{sec:zahteve} navaja merilo, izmerjeno vrednost in sodbo.

\begin{table}[htbp]
  \centering
  \caption{Preverjanje nefunkcionalnih zahtev.}
  \label{tab:nfr}
  \setlength{\tabcolsep}{4pt}
  \begin{tabularx}{\textwidth}{@{}>{\raggedright\arraybackslash}X
                                 >{\raggedright\arraybackslash}X
                                 >{\raggedright\arraybackslash}X c@{}}
    \toprule
    \textbf{Zahteva} & \textbf{Merilo} & \textbf{Izmerjeno} & \textbf{Sodba} \\
    \midrule
    Odziv vmesnika na dejanje uporabnika & $< 100$~ms & \res{x}~ms  & \res{OK} \\
    Zahtevki HTTP (razen vrednotenja)    & $< 300$~ms & \res{x}~ms ($p_{95}$) & \res{OK} \\
    Vrednotenje vozila                   & $< 2$~s    & $13{,}5$~ms ($p_{99}$, brez omrežja) & OK \\
    Vizualna privlačnost / dostopnost    & Lighthouse & \res{x} / \res{x} & \res{OK} \\
    Prilagodljiva postavitev             & preizkus na več velikostih zaslona & \res{da} & \res{OK} \\
    Varovanje uporabniških podatkov      & bcrypt, JWT z rokom & \res{opis} & \res{OK} \\
    Odzivnost pri zahtevnejših opravilih & asinhrona obdelava & \res{opis} & \res{OK} \\
    \bottomrule
  \end{tabularx}
\end{table}

\paragraph{Meritev odzivnosti vmesnika in endpointov.}
\res{Odzivnost uporabniškega vmesnika smo izmerili z orodjem Lighthouse na treh
ključnih straneh; ocene za performanco so X, za dostopnost pa X. Latenco
glavnih endpointov zalednega sistema smo izmerili z orodjem autocannon:
$p_{95}$ je v vseh primerih pod X~ms, torej pod zahtevo 300~ms.}

\paragraph{Varnost.}
Gesla so pred shranjevanjem zgoščena z algoritmom bcrypt
(razdelek~\ref{sec:avtentikacija}), žetoni JWT imajo rok veljavnosti $7$~dni,
poti vmesnika API pa ščiti sredinski preverjevalnik žetonov. \res{Orodje
\texttt{npm audit} ob zadnjem pregledu ni javilo kritičnih ranljivosti v
odvisnostih.}

% =====================================================================
\section{Uporabniški preizkus}
\label{sec:uporabniki}

Izvedli smo formativni uporabniški preizkus~\cite{nielsen1994} z namenom
odkriti težave uporabnosti, ne pa statistično ovrednotiti izdelka.

\paragraph{Udeleženci in postopek.}
V preizkus smo vključili \res{N} udeležencev iz kroga družine in bližnjih
prijateljev, katerih starost in računalniška pismenost sta se med seboj
razlikovali (tabela~\ref{tab:udelezenci}). Vsak udeleženec je prejel enak
seznam nalog: registracija, objava oglasa, všečkanje oglasa, izbira
slovenskega jezika v nastavitvah in vrednotenje poljubnega vozila. Med
izvajanjem je udeleženec razmišljal na glas, opazovalec pa je beležil, ali je
nalogo dokončal samostojno, in vsako opaženo težavo.

\begin{table}[htbp]
  \centering
  \caption{Udeleženci uporabniškega preizkusa.}
  \label{tab:udelezenci}
  \begin{tabular}{c c c l}
    \toprule
    \textbf{\#} & \textbf{Starost (razpon)} & \textbf{Rač.\ pismenost (1--5)} & \textbf{Naprava} \\
    \midrule
    1 & \res{x} & \res{x} & \res{x} \\
    2 & \res{x} & \res{x} & \res{x} \\
    3 & \res{x} & \res{x} & \res{x} \\
    \res{\dots} & & & \\
    \bottomrule
  \end{tabular}
\end{table}

\paragraph{Rezultati.}
Vsi udeleženci so vse naloge dokončali brez pomoči, kar potrjuje, da so glavne
poti v vmesniku dovolj jasne (RV4). Kljub temu so odkrili nekaj
pomanjkljivosti; tabela~\ref{tab:tezave} jih razvršča po resnosti (lestvica
$0$--$4$ po Nielsenu~\cite{nielsen1994}) in navaja, ali so bile odpravljene.

\begin{table}[htbp]
  \centering
  \caption{Težave, odkrite med uporabniškim preizkusom
  (resnost po lestvici $0$--$4$).}
  \label{tab:tezave}
  \setlength{\tabcolsep}{5pt}
  \begin{tabularx}{\textwidth}{@{}>{\raggedright\arraybackslash}X c c@{}}
    \toprule
    \textbf{Težava} & \textbf{Resnost} & \textbf{Status} \\
    \midrule
    Premajhni naslovi podstrani v glavi glavne strani       & \res{x} & odpravljeno \\
    Manjkajoče polje za ponovni vnos gesla pri registraciji & \res{x} & odpravljeno \\
    Velika števila brez ločil za tisočice                   & \res{x} & odpravljeno \\
    Manjkajoč opis pod slikami oblik vozil v obrazcu        & \res{x} & odpravljeno \\
    Premalo vozil, dostopnih za vrednotenje                 & \res{x} & \res{v načrtu} \\
    \bottomrule
  \end{tabularx}
\end{table}

Omejitve preizkusa (majhen in priložnostni vzorec, brez standardiziranega
vprašalnika, brez merjenja časa po nalogah) obravnavamo v
razdelku~\ref{sec:groznje}.

% =====================================================================
\section{Simulacija uporabe}
\label{sec:simulacija}

Za ponazoritev delovanja rešitve povzemamo tipičen potek uporabe. Uporabnik je
ob vstopu usmerjen na glavno stran (slika~\ref{fig:glavna-stran}), kjer sta
izpostavljeni prijava in registracija ter pregled novih in vseh objavljenih
oglasov. Registriran uporabnik lahko objavi oglas prek vodene forme
(slika~\ref{fig:objava-oglasa}), ovrednoti poljubno vozilo, všečka oglase in si
ogleda svoj profil z všečkanimi vozili (slika~\ref{fig:liked-cars}).

\begin{figure}[H]
  \centering
  
\includegraphics[width=0.9\textwidth]{glavna-stran.png}
  \caption{Glavna stran prototipa}
  \label{fig:glavna-stran}
\end{figure}
\begin{figure}[H]
  \centering
  \includegraphics[width=0.9\textwidth]{objava-oglasa.png}
  \caption{Vodena forma za objavo oglasa}
  \label{fig:objava-oglasa}
\end{figure}
\begin{figure}[H]
  \centering
  \includegraphics[width=0.9\textwidth]{liked-cars.png}
  \caption{Stran uporabnika s prikazanimi všečkanimi vozili}
  \label{fig:liked-cars}
\end{figure}

\paragraph{Primerjava ocen z obstoječimi portali.}
Za ilustracijo smo delovanje sistema primerjali s ponudbo na dveh spletnih
portalih na treh vozilih (tabela~\ref{tab:vrednotenje}). Sistematično in
kvantitativno oceno natančnosti podaja razdelek~\ref{sec:natancnost}; ta
primerjava je zgolj kvalitativna ponazoritev.

\begin{table}[htbp]
  \centering
  \caption{Primerjava ocenjenih cen s ponudbo na spletnih portalih (v evrih).}
  \label{tab:vrednotenje}
  \setlength{\tabcolsep}{5pt}
  \begin{tabular}{lccc}
    \toprule
    & \textbf{Audi A3} & \textbf{BMW ser.\ 2} & \textbf{Citroën Berlingo} \\
    \midrule
    slov.\ portal (razpon)     & 4.000--6.000  & 9.000--12.000  & 7.500--10.000 \\
    nem.\ portal (razpon)      & 4.500--9.000  & pribl.\ 11.100 & 8.000--12.000 \\
    \midrule
    Sistem --- dobra cena      & 4.500  & 7.600   & 7.600  \\
    Sistem --- sredinska cena  & 6.350  & 9.300   & 9.941  \\
    Sistem --- slaba cena      & 7.600  & 10.500  & 12.100 \\
    Raznolikost                & $49\,\%$ & $31\,\%$ & $45\,\%$ \\
    \bottomrule
  \end{tabular}
\end{table}

Pri vozilih z nizko raznolikostjo ponudbe (Audi~A3, BMW serije~2) se sredinska
ocena dobro ujema s cenami na obeh portalih. Pri modelu Citroën Berlingo je
raznolikost bistveno višja, zato je razpon med dobro in slabo oceno širši,
sredinska ocena pa manj zanesljiva --- kar je skladno z ugotovitvijo iz
razdelka~\ref{sec:natancnost}, da napaka strmo narašča z raznolikostjo ponudbe.
Pri interpretaciji je treba upoštevati, da sistem ocenjuje pričakovano prodajno
in ne oglaševano ceno.

% =====================================================================
\section{Grožnje veljavnosti}
\label{sec:groznje}

\paragraph{Notranja veljavnost.}
Kot merilo resnice pri merjenju natančnosti uporabljamo \emph{oglaševano} ceno,
ne dejanske prodajne cene, ki ni javno dostopna. Namerni popust, s katerim
sistem cilja pričakovano prodajno ceno, je zato ocenjen in ne umerjen na
dejanskih prodajah. Isti fizični oglas se lahko pojavi pri več virih; prečno
preverjanje izloči zapise z istim izvornim oglasom, ne more pa zaznati
podvojitev med viri.

\paragraph{Zunanja veljavnost.}
Podatki pokrivajo $10$ znamk in pretežno slovenski trg (z manjšim nemškim
vzorcem); znamka z največ oglasi na trgu (okoli $8000$ oglasov) v šifrantu ni
zajeta in je iz analize izpuščena. Uporabniški preizkus je temeljil na majhnem
in priložnostnem vzorcu.

\paragraph{Konstruktna veljavnost.}
Pojem \enquote{tržna vrednost} operacionaliziramo kot uteženo povprečje
oglaševanih cen primerljivih vozil, znižano za pričakovani pogajalski popust;
\enquote{raznolikost ponudbe} pa kot relativno širino interkvartilnega razpona
teh cen. Drugačna operacionalizacija bi lahko dala drugačne rezultate.

\paragraph{Zanesljivost.}
Meritve latence in obremenitve so bile izvedene na deljenem produkcijskem
strežniku, kjer lahko na rezultate vpliva ozadje. Meritve natančnosti so
ponovljive nad fiksnim posnetkom baze.

% =====================================================================
\section{Povzetek ugotovitev}
\label{sec:povzetek-analize}

Tabela~\ref{tab:rv-povzetek} povzema odgovore na raziskovalna vprašanja glede
na vnaprej postavljena merila.

\begin{table}[htbp]
  \centering
  \caption{Povzetek odgovorov na raziskovalna vprašanja.}
  \label{tab:rv-povzetek}
  \begin{tabularx}{\textwidth}{@{}l X c@{}}
    \toprule
    \textbf{RV} & \textbf{Glavni rezultat} & \textbf{Sodba} \\
    \midrule
    RV1 natančnost & MdAPE $13{,}5\,\%$ ($9$--$10\,\%$ za dobro pokrita vozila),
                     bistveno bolje od naivne ocene ($37{,}8\,\%$), primerljivo z
                     drugimi metodami & izpolnjeno \\
    RV2 odvisnost  & napaka pada z velikostjo vzorca (ustali se pri $n\approx20$),
                     narašča s starostjo, kilometrino in raznolikostjo ponudbe
                     ($7{,}5\,\%$ do $29{,}7\,\%$) & pojasnjeno \\
    RV3 hitrost    & vrednotenje $2$~ms (p50), $13{,}5$~ms (p99); indeks odpravi
                     odvisnost od velikosti baze; \res{meja X~z/s} & \res{izpolnjeno} \\
    RV4 uporabnost & vse naloge dokončane brez pomoči; \res{X} težav odpravljenih & izpolnjeno \\
    \bottomrule
  \end{tabularx}
\end{table}

\paragraph{Sodba o zastavljenih izhodiščih.}
Rezultati \emph{potrjujejo} izhodišča, ki smo si jih zastavili na začetku. Sistem na podlagi štirih
atributov vozila vrne oceno vrednosti z mediano odstotne napake $13{,}5\,\%$
(brez namernega popusta), pri dobro pokritih vozilih pa $9$--$10\,\%$, kar je
znotraj zastavljenega praga $15\,\%$ in bistveno bolje od naivne metode.
Predlagani pristop se med primerljivo natančnimi metodami ne uvršča na prvo
mesto --- preprost $k$-najbližjih sosedov doseže $11{,}2\,\%$ --- je pa edini,
ki poleg točkovne ocene vrne tudi razložljiv cenovni razpon in indeks
zaupanja, ta pa se izkaže za dober pokazatelj dejanske zanesljivosti
posamezne ocene. Glavne omejitve so odvisnost od enega prevladujočega
podatkovnega vira, nižja natančnost pri starejših vozilih in modelih z
razpršeno ponudbo ter merilo resnice, ki temelji na oglaševanih in ne
prodajnih cenah. Konkretne smeri za izboljšave --- vključno s prehodom na
metode strojnega učenja, ki jih ta analiza dodatno utemeljuje --- navajamo v
poglavju~\ref{ch:sklep}.

%
%  Številke so vpisane iz izhodov skript v CarStockMarketEvaluation/analysis/.
%  Preostale oznake \res{...} (rdeče) je treba še izpolniti:
%    - obremenitveni test (benchmarks/loadtest.py)
%    - sledljivost testov, Jest coverage, Lighthouse
%    - tabeli udeležencev in najdenih težav
%  Slike (PDF) iz analysis/out/ prekopiraj v mapo slike/.
% =====================================================================

\newcommand{\res}[1]{\textcolor{red}{\textbf{[#1]}}}

\chapter{Analiza in vrednotenje rešitve}
\label{ch:analiza}

V tem poglavju sistematično pregleadamo in analiziramo razvito rešitev. Najprej iz zahtev in
izhodišč, ki smo si jih zastavili na začetku, izpeljemo raziskovalna vprašanja in za vsako določimo merilo
uspešnosti (razdelek~\ref{sec:rq}). Nato opišemo merilno okolje in postopke
(razdelek~\ref{sec:metodologija}) ter natančno definiramo vse uporabljene mere
(razdelek~\ref{sec:mere}). Sledijo rezultati: značilnosti zbranih
podatkov~(\ref{sec:podatki}), natančnost ocene
vrednosti~(\ref{sec:natancnost}), občutljivost algoritma na
parametre~(\ref{sec:obcutljivost}), latenca posameznega
vrednotenja~(\ref{sec:latenca}) in obnašanje pod
obremenitvijo~(\ref{sec:obremenitev}). Zatem preverimo funkcionalne in
nefunkcionalne zahteve (\ref{sec:funkcionalne}, \ref{sec:nefunkcionalne}) ter
predstavimo uporabniški preizkus~(\ref{sec:uporabniki}) in simulacijo
uporabe~(\ref{sec:simulacija}). Poglavje sklenemo z obravnavo groženj
veljavnosti~(\ref{sec:groznje}) in povzetkom ugotovitev glede na
zastavljena izhodišča~(\ref{sec:povzetek-analize}).

% =====================================================================
\section{Raziskovalna vprašanja in merila uspešnosti}
\label{sec:rq}

Na začetku diplomskega dela (poglavje~\ref{ch:uvod}) smo si zastavili, da je
mogoče uporabniku na podlagi zbranih tržnih podatkov in primerjave s podobnimi
vozili ponuditi \emph{uporabno} in \emph{dovolj natančno} oceno tržne vrednosti
vozila in da vključitev takega sistema izboljša uporabniško izkušnjo. Iz tega
izhodišča in iz nefunkcionalnih zahtev (razdelek~\ref{sec:zahteve}) izpeljemo
štiri raziskovalna vprašanja (RV1--RV4) in za vsako določimo merilo uspešnosti.

\emph{Mediana absolutne odstotne napake} (MdAPE) je osrednja mera natančnosti v
tem poglavju; nanjo se sklicujejo merila uspešnosti in vsi nadaljnji razdelki.
Izračunamo jo tako, da za vsak testni primer določimo absolutno odstotno napako,
torej za koliko odstotkov se ocenjena cena razlikuje od dejanske, nato pa čez
vse primere vzamemo mediano teh vrednosti. MdAPE $=15\,\%$ tako pomeni, da je pri
polovici vozil ocena znotraj $15\,\%$ od dejanske cene, pri drugi polovici pa je
odstopanje večje; manjša kot je vrednost, natančnejše so ocene. Mediano
uporabljamo namesto povprečja, ker je manj občutljiva na posamezne skrajne
napake. Natančne definicije MdAPE in preostalih mer so v
razdelku~\ref{sec:mere}.

Raziskovalna vprašanja so:

\begin{itemize}
  \item \textbf{RV1 --- natančnost ocene.} Kako natančno sistem oceni tržno
        vrednost vozila v primerjavi z dejanskimi oglasi na trgu?
        \emph{Merilo:} MdAPE je pri vozilih z vsaj $10$ primerljivimi
        vozili manjša od $15\,\%$, ocena pa je bistveno natančnejša od
        preproste (naivne) referenčne metode.
  \item \textbf{RV2 --- odvisnost natančnosti od podatkov.} Od česa je
        natančnost ocene najbolj odvisna (velikost vzorca primerljivih vozil,
        razpršenost ponudbe, segment vozila)?
  \item \textbf{RV3 --- hitrost in razširljivost.} Ali sistem zadošča
        nefunkcionalnim zahtevam glede odzivnosti tudi pod sočasno
        obremenitvijo, in kje je zgornja meja trenutne postavitve?
  \item \textbf{RV4 --- uporabnost.} Ali ciljni uporabniki znajo brez navodil
        opraviti glavne naloge (registracija, objava oglasa, vrednotenje)?
\end{itemize}
% =====================================================================
\section{Metodologija in merilno okolje}
\label{sec:metodologija}

\paragraph{Merilno okolje.}
Meritve hitrosti (razdelka~\ref{sec:latenca} in~\ref{sec:obremenitev}) so
izvedene v \emph{produkcijskem} okolju, torej na navideznem zasebnem strežniku
z dvema procesorskima jedroma in $4$~GB pomnilnika ter s tehnologijami in
njihovimi različicami, opisanimi v poglavju~\ref{ch:razvoj}
(razdelek~\ref{sec:namestitev}). Analitične in merilne skripte tečejo nad isto
podatkovno bazo kot je na produkciji.

\paragraph{Posnetek podatkov.}
Meritve v tem poglavju temeljijo na posnetku podatkovne baze iz septembra
$2026$, ki obsega približno $86.000$ normaliziranih oglasov
(razdelek~\ref{sec:podatki}). Ker se posnetek med analizo ne spreminja, so
rezultati ponovljivi.

\paragraph{Orodja.}
Za meritve in analizo v tem poglavju smo napisali namenske skripte, ločene od
produkcijske kode; ena od njih je dokumentirana kopija ocenjevalnega
algoritma, parametrizirana za eksperimente. Statistično obdelavo in izris
grafov smo izvedli s knjižnicama \texttt{NumPy}~\cite{harris2020numpy} in
\texttt{Matplotlib}~\cite{hunter2007matplotlib}.

\paragraph{Natančnost.}
Natančnost ocene (razdelek~\ref{sec:natancnost}) merimo z naključno izbranim
normaliziranim oglasom: njegove atribute (znamka, model, letnik, kilometrina)
uporabimo kot vhod v algoritem, sam oglas pa izločimo iz primerjalne množice. Nato izračunamo oceno in jo
primerjamo z \emph{dejansko oglaševano} ceno tega oglasa. Postopek ponovimo za
$3000$ naključnih oglasov.

% =====================================================================
\section{Uporabljene mere}
\label{sec:mere}

\paragraph{Latenca}
 je čas, ki preteče od trenutka tik pred oddajo zahtevka do
trenutka, ko odjemalec prejme celoten odgovor. Glede na to, kolikšen del te
poti zajamemo, ločimo tri vrste:
\begin{itemize}
  \item[(a)] \emph{latenca obdelave na strežniku} --- čas izvajanja funkcije
        za vrednotenje, od začetka do konca, brez omrežja;
  \item[(b)] \emph{latenca poizvedbe SQL} --- čas izvajanja poizvedbe v podatkovni bazi,
   od začetka do konca, brez omrežja;
  \item[(c)] \emph{latenca od konca do konca} (angl.\ \emph{end-to-end}) ---
        (a) skupaj z omrežnim obhodom med odjemalcem in strežnikom ter
        obdelavo zahtevka HTTP; to je čas, ki ga zazna uporabnik.
\end{itemize}
V razdelku~\ref{sec:latenca} merimo vrsti (a) in (b), v
razdelku~\ref{sec:obremenitev} pa vrsto (c).

\paragraph{Percentili.}
\textbf{Percentil} $p_k$ je vrednost, pod katero pade $k\,\%$ meritev. Poročamo
$p_{50}$ (mediana), $p_{95}$ in $p_{99}$.

\paragraph{Prepustnost, sočasnost, nasičenje.}
Obnašanje sistema pod obremenitvijo opišemo s tremi količinami:
\begin{itemize}
  \item \textbf{prepustnost} --- število uspešno obdelanih zahtevkov na
        sekundo;
  \item \textbf{sočasnost} --- število odjemalcev, ki hkrati držijo odprt
        zahtevek;
  \item \textbf{nasičenje} (\emph{koleno}) --- točka, v kateri dodatna
        sočasnost ne poveča več prepustnosti, latenca pa začne naraščati
        približno linearno.
\end{itemize}
Notranjo skladnost meritev preverimo z \textbf{Littlovim zakonom}~\cite{little1961},
\begin{equation}
  \label{eq:little}
  L = \lambda \cdot W,
\end{equation}
kjer je $L$ povprečno število zahtevkov v sistemu, $\lambda$ prepustnost in
$W$ povprečna latenca.

\paragraph{Napaka ocene.}
Za oglas $i$ z dejansko ceno $p_i$ in ocenjeno ceno $\hat p_i$ definiramo
absolutno odstotno napako
\begin{equation}
  \mathrm{APE}_i = \frac{\lvert \hat p_i - p_i \rvert}{p_i}.
\end{equation}
\textbf{MAPE} je povprečje $\mathrm{APE}_i$, \textbf{MdAPE} pa njihova mediana;
MdAPE je robustnejša na osamelce in je naša glavna mera~\cite{hyndman2006}.
Poročamo tudi \textbf{MAE} (povprečna absolutna napaka) in \textbf{RMSE}
(koren povprečne kvadratne napake) v evrih ter \emph{predznačeno} odstotno
napako, ki pokaže sistematičen odmik ocene navzgor ali navzdol.

\paragraph{Pokritost.}
\textbf{Pokritost} je delež poizvedb, za katere sistem vrne oceno, torej za
katere najde vsaj eno (oziroma vsaj $5$ oz.\ $10$) primerljivo vozilo.

\paragraph{Raznolikost ponudbe.}
\textbf{Raznolikost} (v izvorni kodi \texttt{ratio}) je relativna širina
cenovnega razpona primerljivih vozil,
\begin{equation}
  \label{eq:raznolikost}
  \mathrm{raznolikost} = \frac{P_{75} - P_{25}}{\hat{p}},
\end{equation}
kjer sta $P_{25}$ in $P_{75}$ (v kodi \texttt{goodPrice} in \texttt{badPrice})
uteženi $25.$ oziroma $75.$ percentil cen primerljivih vozil, $\hat{p}$
(v kodi \texttt{estimatedPrice}) pa ocenjena cena vozila. Vrednost izražamo v
odstotkih; enoten zapis uporabljamo v vseh tabelah tega poglavja.

% =====================================================================
\section{Značilnosti zbranih podatkov}
\label{sec:podatki}
\paragraph{Viri in obseg.}
Velika večina oglasov je bila zajeta z največjega slovenskega spletnega portala
za prodajo rabljenih vozil (razdelek~\ref{sec:zajem-podatkov}); manjši del
predstavlja vzorec z večjega nemškega oglasnega portala, ki smo ga uporabili za
preverjanje delovanja cevovoda na drugačni strukturi vira. Zajem je zajel
javno dostopne oglase za nekomercialno raziskovalno rabo, z omejeno hitrostjo
zahtevkov in brez shranjevanja kontaktnih ali osebnih podatkov prodajalcev.
Po normalizaciji baza obsega $85.846$ oglasov ($83.493$ s slovenskega in
$2.353$ z nemškega portala) za $10$ znamk vozil; razčlenitev po znamki je v
tabeli~\ref{tab:podatki-znamke}.

\begin{table}[htbp]
  \centering
  \caption{Število normaliziranih oglasov po znamki.}
  \label{tab:podatki-znamke}
  % vir: analysis/out/dataset_per_brand.csv
  \begin{tabular}{lr lr}
    \toprule
    \textbf{Znamka} & \textbf{Oglasi} & \textbf{Znamka} & \textbf{Oglasi} \\
    \midrule
    Volkswagen    & 21.191 & Peugeot     & 7.513 \\
    Audi          & 11.842 & Citroën     & 6.729 \\
    BMW           & 10.724 & Ford        & 5.783 \\
    Renault       & 9.615  & Toyota      & 2.648 \\
    Mercedes-Benz & 9.355  & Alfa Romeo  & 446 \\
    \bottomrule
  \end{tabular}
\end{table}

\paragraph{Porazdelitve ključnih atributov.}
Slika~\ref{fig:porazdelitve} prikazuje porazdelitve oglaševanih cen, prevožene
kilometrine in letnika. Porazdelitev cen je izrazito desno asimetrična (dolg
rep dragih vozil), zato jo prikazujemo na logaritemski osi. Mediana cene je
$14.890$~EUR (povprečje $18.137$~EUR), mediana kilometrine $159.950$~km,
mediana letnika pa $2017$; podrobnosti so v tabeli~\ref{tab:podatki-porazdelitve}.
Zaradi berljivosti so na grafih odrezane skrajne vrednosti (nekaj oglasov ima
očitno napačno ceno, kilometrino ali letnik); njihovo število je označeno v
kotu vsakega grafa, celoten razpon pa je razviden iz stolpca \emph{maks.}\ v
tabeli~\ref{tab:podatki-porazdelitve} in iz obravnave osamelcev v
tabeli~\ref{tab:osamelci}.

\begin{table}[htbp]
  \centering
  \caption{Povzetek porazdelitev ključnih numeričnih atributov.}
  \label{tab:podatki-porazdelitve}
  % vir: analysis/out/dataset_distribution_summary.json
  \setlength{\tabcolsep}{5pt}
  \begin{tabular}{lrrrrr}
    \toprule
    & \textbf{$p_{5}$} & \textbf{mediana} & \textbf{povprečje} & \textbf{$p_{95}$} & \textbf{maks.} \\
    \midrule
    cena [EUR]        & 1.550  & 14.890  & 18.137  & 49.014  & 374.000 \\
    kilometrina [km]  & 26.000 & 159.950 & 169.464 & 334.000 & 1.974.001 \\
    letnik            & 2005   & 2017    & 2016    & 2024    & 2029 \\
    moč [kW]          & 56     & 110     & 117     & 221     & 920 \\
    \bottomrule
  \end{tabular}
\end{table}

\begin{figure}[htbp]
  \centering
  \begin{subfigure}{0.48\textwidth}
    \includegraphics[width=\textwidth]{dataset_price_hist.pdf}
    \caption{oglaševana cena (log)}
  \end{subfigure}\hfill
  \begin{subfigure}{0.48\textwidth}
    \includegraphics[width=\textwidth]{dataset_km_hist.pdf}
    \caption{prevoženi kilometri}
  \end{subfigure}
  \\[1ex]
  \begin{subfigure}{0.48\textwidth}
    \includegraphics[width=\textwidth]{dataset_year_hist.pdf}
    \caption{letnik prve registracije}
  \end{subfigure}\hfill
  \begin{subfigure}{0.48\textwidth}
    \includegraphics[width=\textwidth]{dataset_brand_bar.pdf}
    \caption{oglasi po znamki}
  \end{subfigure}
  \caption{Porazdelitve ključnih atributov v normalizirani bazi.}
  \label{fig:porazdelitve}
  % vir: analysis/out/dataset_*_hist.pdf, dataset_brand_bar.pdf
\end{figure}

\paragraph{Pokritost po segmentih.}
Za pošteno oceno mora imeti algoritem dovolj primerljivih vozil v ožjem
segmentu (ista znamka in model, podoben letnik in kilometrina).
Slika~\ref{fig:pokritost} prikazuje toplotno karto za znamki Renault in
Volkswagen: vrstice so modeli, stolpci letniki, barva pa
število oglasov v tej kombinaciji. Vidno je, da so pogosti modeli srednjih
letnikov (npr.\ Volkswagen Golf ali Renault Clio) gosto pokriti, redki modeli
in skrajni letniki pa le redko, kar neposredno pojasni razlike v zanesljivosti
ocen (razdelek~\ref{sec:natancnost}).

\begin{figure}[htbp]
  \centering
  \includegraphics[width=\textwidth]{dataset_coverage_heatmap.pdf}
  \caption{Pokritost: število oglasov po modelu in letniku za znamki
  Renault in Volkswagen (po ena znamka na vrstico).}
  \label{fig:pokritost}
  % vir: analysis/out/dataset_coverage_heatmap.pdf
\end{figure}

\paragraph{Osamelci in kakovost podatkov.}
Postopek normalizacije trenutno ne izvaja filtriranja osamelcev, zato baza
vsebuje tudi vozila z očitno napačnimi vrednostmi.
Tabela~\ref{tab:osamelci} prikazuje njihov delež. Kar $10{,}7\,\%$ oglasov ima
vsaj eno vrednost zunaj smiselnih mej, največ pri kilometrini ($9{,}9\,\%$;
najvišja zabeležena vrednost je skoraj dva milijona kilometrov) in letniku
($8{,}1\,\%$; nekaj oglasov ima celo letnik v prihodnosti). Preprost filter po
interkvartilnem razponu (IQR) na ceni znotraj skupine
$(\text{znamka},\text{model})$ bi odstranil dodatnih $2{,}2\,\%$ oglasov.
Vpliv doslednega filtriranja osamelcev na natančnost obravnavamo v
razdelku~\ref{sec:obcutljivost}, uvedbo takega filtra pa v
poglavju~\ref{ch:sklep}.

\begin{table}[H]
  \centering
  \caption{Osamelci v normalizirani bazi ($85.846$ oglasov).}
  \label{tab:osamelci}
  % vir: analysis/out/dataset_outliers.csv
  \begin{tabularx}{\textwidth}{@{}X r r@{}}
    \toprule
    \textbf{Merilo} & \textbf{Št.\ vrstic} & \textbf{Delež} \\
    \midrule
    cena zunaj $[500,\ 200.000]$~EUR         & 427   & $0{,}5\,\%$ \\
    kilometrina zunaj $[1000,\ 500.000]$~km  & 8.495 & $9{,}9\,\%$ \\
    letnik zunaj $[1990,\ 2027]$             & 6.979 & $8{,}1\,\%$ \\
    vsaj ena vrednost zunaj mej              & 9.169 & $10{,}7\,\%$ \\
    dodatno po filtru IQR na ceni            & 1.878 & $2{,}2\,\%$ \\
    \bottomrule
  \end{tabularx}
\end{table}

% =====================================================================
\section{Natančnost ocene vrednosti}
\label{sec:natancnost}

To je osrednja meritev diplomskega dela: neposredno preverja zastavljeno
izhodišče o \emph{dovolj natančni} oceni (RV1). Protokol je opisan v
razdelku~\ref{sec:metodologija}, rezultate pa izdela skripta
\texttt{analysis/accuracy\_cv.py} nad vzorcem $3000$ oglasov (od tega jih je
$2986$ dobilo oceno).

\paragraph{Merilo resnice in namerni odmik ocene.}
Kot merilo resnice uporabljamo \emph{oglaševano} ceno, saj dejanske prodajne
cene niso javno dostopne. Sistem pa namenoma ne ocenjuje oglaševane cene,
temveč \emph{pričakovano prodajno ceno}: ker se rabljena vozila v pogajanjih
praviloma prodajo pod oglaševano ceno, oceno zniža s stopenjskim popustom
($9$--$15\,\%$ glede na cenovni razred). Zato je pri primerjavi z oglaševano
ceno pričakovan rahel \emph{negativni} odmik ocene, kar je zaželeno in ne
pomeni slabše ocene. To lastnost upoštevamo pri branju vseh rezultatov v
nadaljevanju: napako merimo dvakrat --- proti oglaševani ceni z vključenim
popustom (kar dejansko vrne sistem) in brez popusta (čista primerjalna
ocena).

\paragraph{Skupni rezultat.}
Brez popusta znaša MdAPE predlaganega pristopa $13{,}5\,\%$, MAPE $23{,}8\,\%$,
MAE $2798$~EUR in RMSE $5151$~EUR; z vključenim namernim popustom je MdAPE
proti oglaševani ceni $14{,}5\,\%$ (mediana predznačene napake $-6{,}8\,\%$).
Pokritost je $99{,}5\,\%$ pri $n \ge 1$ in $95{,}7\,\%$ pri $n \ge 10$.
Slika~\ref{fig:pred-vs-actual} prikazuje napovedano ceno proti dejanski
($R^2 = 0{,}87$), slika~\ref{fig:ape-hist} pa porazdelitev odstotne napake.

\begin{figure}[htbp]
  \centering
  \begin{subfigure}{0.48\textwidth}
    \includegraphics[width=\textwidth]{accuracy_pred_vs_actual.pdf}
    \caption{napovedano proti dejanskemu}
    \label{fig:pred-vs-actual}
  \end{subfigure}\hfill
  \begin{subfigure}{0.48\textwidth}
    \includegraphics[width=\textwidth]{accuracy_ape_hist.pdf}
    \caption{porazdelitev napake}
    \label{fig:ape-hist}
  \end{subfigure}
  \caption{Natančnost ocene vrednosti (leave-one-out, brez popusta).}
  % vir: analysis/out/accuracy_pred_vs_actual.pdf, accuracy_ape_hist.pdf
\end{figure}

\paragraph{Primerjava z osnovnimi modeli.}
Da presodimo, ali je predlagani pristop boljši od trivialnega, ga primerjamo s
štirimi referenčnimi modeli na istem protokolu in isti množici kandidatov
(tabela~\ref{tab:baselines}, slika~\ref{fig:baselines}): B0~globalna mediana
cene za dani model, B1~navadno povprečje cen v filtrirnem oknu, B1m~mediana cen
v oknu, B2~linearna regresija $\text{cena} \sim \text{letnik} + \text{km}$,
B3~metoda $k$~najbližjih sosedov ($k = 10$) po standardiziranih letniku in
kilometrini, ter B4~naš pristop (uteženo povprečje s podobnostnim jedrom, brez
popusta).

\begin{table}[htbp]
  \centering
  \caption{Primerjava natančnosti z osnovnimi modeli (leave-one-out).
  MdAPE in MAPE sta v odstotkih, MAE in RMSE v evrih.}
  \label{tab:baselines}
  % vir: analysis/out/accuracy_overall.csv
  \setlength{\tabcolsep}{5pt}
  \begin{tabular}{l r r r r}
    \toprule
    \textbf{Model} & \textbf{MdAPE} & \textbf{MAPE} & \textbf{MAE} & \textbf{RMSE} \\
    \midrule
    B0 globalna mediana       & 37{,}8 & 95{,}7 & 7395 & 10883 \\
    B1 povprečje v oknu       & 14{,}2 & 24{,}9 & 2930 & 5333 \\
    B1m mediana v oknu        & 13{,}1 & 22{,}5 & 2877 & 5561 \\
    B2 linearna regresija     & 12{,}9 & 22{,}4 & 2677 & 5032 \\
    B3 $k$-najbližjih sosedov & \textbf{11{,}2} & 21{,}5 & 2497 & 4883 \\
    B4 naš pristop            & 13{,}5 & 23{,}8 & 2798 & 5151 \\
    \bottomrule
  \end{tabular}
\end{table}

\begin{figure}[htbp]
  \centering
  \includegraphics[width=0.7\textwidth]{accuracy_baselines.pdf}
  \caption{MdAPE osnovnih modelov in predlaganega pristopa.}
  \label{fig:baselines}
  % vir: analysis/out/accuracy_baselines.pdf
\end{figure}

Predlagani pristop je \emph{bistveno} boljši od naivne globalne mediane
($13{,}5\,\%$ proti $37{,}8\,\%$) in primerljiv s preostalimi metodami, ki se
gibljejo med $11$ in $14\,\%$. Ni pa najboljši: metoda $k$~najbližjih sosedov
doseže $11{,}2\,\%$, tudi navadna mediana istega okna ($13{,}1\,\%$) je
malenkost natančnejša. To je pričakovano za zasnovo, katere cilj je
\emph{razložljivost}: za razliko od naštetih metod predlagani pristop poleg
točkovne ocene vrne tudi cenovni razpon in indeks zaupanja, uteži pa so
neposredno interpretabilne. Rezultat hkrati potrjuje smiselnost prehoda na
metode strojnega učenja, ki jo predlagamo v poglavju~\ref{ch:sklep}: če že
preprost $k$-NN preseže sedanji pristop, je od zmogljivejših modelov mogoče
pričakovati še boljše rezultate.

\paragraph{Odvisnost od velikosti vzorca in segmenta (RV2).}
Slika~\ref{fig:mdape-n} prikazuje MdAPE glede na število primerljivih vozil
$n$: napaka je pri $n < 10$ okoli $20\,\%$, pri $n \ge 20$ pa se ustali pri
približno $13\,\%$. To empirično podpira izbiro konstante $20$ v indeksu
zaupanja (enačba~\ref{eq:confidence}) --- pri toliko primerljivih vozilih je
ocena že blizu svoje najboljše natančnosti. Slika~\ref{fig:mdape-brand} prikazuje MdAPE po znamki: najbolje se
odrežejo pogosti, gosto pokriti modeli (Audi $9{,}8\,\%$, Toyota $10{,}8\,\%$,
Volkswagen $12{,}1\,\%$), najslabše pa Mercedes-Benz ($24{,}9\,\%$), pri
katerem se pri normalizaciji več različnih modelov zlije v isto kanonično
ime, kar poslabša izbor primerljivih vozil.
Tabela~\ref{tab:natancnost-segmenti} kaže še razčlenitev po starosti,
kilometrini in raznolikosti ponudbe.

\begin{figure}[htbp]
  \centering
  \begin{subfigure}{0.48\textwidth}
    \includegraphics[width=\textwidth]{accuracy_mdape_vs_n.pdf}
    \caption{glede na velikost vzorca $n$}
    \label{fig:mdape-n}
  \end{subfigure}\hfill
  \begin{subfigure}{0.48\textwidth}
    \includegraphics[width=\textwidth]{accuracy_mdape_by_brand.pdf}
    \caption{po znamki}
    \label{fig:mdape-brand}
  \end{subfigure}
  \caption{Odvisnost natančnosti ocene od podatkov.}
  % vir: analysis/out/accuracy_mdape_vs_n.pdf, accuracy_mdape_by_brand.pdf
\end{figure}

\begin{table}[htbp]
  \centering
  \caption{Natančnost predlaganega pristopa po segmentih (MdAPE).}
  \label{tab:natancnost-segmenti}
  % vir: analysis/out/accuracy_by_age.csv, accuracy_by_km.csv, accuracy_by_ratio.csv
  \begin{tabularx}{\textwidth}{@{}l X r r@{}}
    \toprule
    \textbf{Razčlenitev} & \textbf{Koš} & \textbf{Primeri} & \textbf{MdAPE [\%]} \\
    \midrule
    starost [leta] & 0--3   & 367 & 9{,}0 \\
                   & 4--7   & 921 & 10{,}6 \\
                   & 8--12  & 932 & 14{,}5 \\
                   & $13+$  & 766 & 22{,}9 \\
    \midrule
    kilometrina    & $<50.000$        & 280  & 9{,}6 \\
                   & $50.000$--$100.000$  & 449  & 11{,}9 \\
                   & $100.000$--$150.000$ & 634  & 12{,}4 \\
                   & $150.000$--$200.000$ & 629  & 13{,}7 \\
                   & $>200.000$       & 1021 & 16{,}7 \\
    \midrule
    raznolikost    & $0$--$0{,}10$    & 88   & 7{,}5 \\
                   & $0{,}10$--$0{,}25$ & 1002 & 8{,}4 \\
                   & $0{,}25$--$0{,}50$ & 1613 & 15{,}5 \\
                   & $>0{,}50$        & 447  & 29{,}7 \\
    \bottomrule
  \end{tabularx}
\end{table}

Vzorec pokaže jasne zakonitosti: napaka narašča s starostjo vozila ($9\,\%$ pri
najnovejših, $23\,\%$ pri vozilih, starejših od $13$~let), s prevoženo
kilometrino in --- najmočneje --- z \emph{raznolikostjo ponudbe}. Pri modelih z
enotno ponudbo (raznolikost pod $0{,}1$) je MdAPE le $7{,}5\,\%$, pri razpršeni
ponudbi (nad $0{,}5$) pa $29{,}7\,\%$. To potrjuje, da je raznolikost, ki jo
sistem uporabniku tudi prikaže, dober vgrajen pokazatelj zanesljivosti
posamezne ocene.

\paragraph{Vpliv popusta.}
Slika~\ref{fig:popust} prikazuje predznačeno napako z in brez popusta. Brez
popusta je ocena rahlo optimistična (mediana $+3{,}8\,\%$ nad oglaševano ceno),
z namernim popustom pa se premakne na $-6{,}8\,\%$ pod oglaševano ceno. Ta
premik je skladen z zasnovo sistema, ki cilja pričakovano prodajno in ne
oglaševane cene; dokončno umerjanje popusta bi bilo mogoče šele s podatki o
dejanskih prodajah (razdelek~\ref{sec:groznje}).

\begin{figure}[htbp]
  \centering
  \includegraphics[width=0.7\textwidth]{accuracy_discount_bias.pdf}
  \caption{Predznačena napaka proti oglaševani ceni, z in brez namernega popusta.}
  \label{fig:popust}
  % vir: analysis/out/accuracy_discount_bias.pdf
\end{figure}

\paragraph{Časovni razrez.}
Pri napovedovanju najnovejših $10\,\%$ oglasov zgolj iz starejših $90\,\%$
znaša MdAPE predlaganega pristopa $13{,}0\,\%$ (tabela~\ref{tab:temporal}), kar
je celo malenkost bolje od rezultata leave-one-out in potrjuje, da ocena ne
izhaja iz uhajanja informacije.

\begin{table}[htbp]
  \centering
  \caption{Natančnost pri časovnem razrezu (napoved najnovejših $10\,\%$).
  MdAPE in MAPE sta v odstotkih, MAE in RMSE v evrih.}
  \label{tab:temporal}
  % vir: analysis/out/accuracy_temporal.csv
  \setlength{\tabcolsep}{5pt}
  \begin{tabular}{l r r r r}
    \toprule
    \textbf{Model} & \textbf{MdAPE} & \textbf{MAPE} & \textbf{MAE} & \textbf{RMSE} \\
    \midrule
    B0 globalna mediana       & 35{,}9 & 74{,}5 & 7596 & 11222 \\
    B1m mediana v oknu        & 12{,}8 & 19{,}7 & 3088 & 5805 \\
    B3 $k$-najbližjih sosedov & 10{,}5 & 18{,}2 & 2667 & 4890 \\
    B4 naš pristop            & 13{,}0 & 20{,}6 & 2994 & 5347 \\
    \bottomrule
  \end{tabular}
\end{table}

% =====================================================================
\section{Občutljivost algoritma na parametre}
\label{sec:obcutljivost}

Ocenjevalni algoritem vsebuje več hevristično določenih konstant: dovoljeno
odstopanje kilometrine ($\pm 40.000$~km) in letnika ($\pm 3$~leta), uteži
delnih ocen ($0{,}6$ in $0{,}4$), oblika padajoče funkcije in stopenjski
popust. V tem razdelku vsako spreminjamo ob nespremenjenih ostalih in merimo
MdAPE ter pokritost (vzorec $2000$ oglasov, leave-one-out). Rezultate izdela
skripta \texttt{analysis/sensitivity.py}.

\paragraph{Okno primerljivosti.}
Tabela~\ref{tab:sens-okno} in slika~\ref{fig:sens-km} kažeta, da je algoritem
na širino \emph{kilometrinskega} okna skoraj neobčutljiv: MdAPE se med
$\pm 10.000$~km in $\pm 80.000$~km giblje le med $13{,}2$ in $13{,}6\,\%$.
Nasprotno je odvisen od \emph{letnikovega} okna: pri $\pm 1$~letu je MdAPE
$12{,}6\,\%$, pri $\pm 5$~letih pa $15{,}2\,\%$. Ožje okno da bolj primerljiva
vozila, a manjši vzorec (nižja pokritost); gre za znano razmerje med
pristranskostjo in varianco~\cite{bias-variance}. Toplotna karta
(slika~\ref{fig:sens-heatmap}) potrjuje, da natančnost določa skoraj izključno
letnikovo okno. Sedanja izbira $\pm 3$~leta je razumen kompromis med
natančnostjo in pokritostjo; morebitna zožitev na $\pm 2$~leti bi MdAPE
znižala za približno $0{,}3$~odstotne točke ob še vedno visoki pokritosti.

\begin{table}[htbp]
  \centering
  \caption{Občutljivost na širino okna primerljivosti. Pokritost je delež
  poizvedb z vsaj $5$ primerljivimi vozili.}
  \label{tab:sens-okno}
  % vir: analysis/out/sensitivity_km.csv, sensitivity_year.csv
  \setlength{\tabcolsep}{5pt}
  \begin{tabular}{l r r r}
    \toprule
    \textbf{Parameter} & \textbf{MdAPE [\%]} & \textbf{Pokritost [\%]} & \textbf{Mediana $n$} \\
    \midrule
    kilometrina $\pm 10.000$~km & 13{,}2 & 92{,}3 & 53 \\
    kilometrina $\pm 40.000$~km & 13{,}4 & 98{,}1 & 196 \\
    kilometrina $\pm 80.000$~km & 13{,}6 & 99{,}0 & 362 \\
    \midrule
    letnik $\pm 1$~leto  & 12{,}6 & 95{,}6 & 97 \\
    letnik $\pm 3$~leta  & 13{,}4 & 98{,}1 & 196 \\
    letnik $\pm 5$~let   & 15{,}2 & 98{,}5 & 262 \\
    \bottomrule
  \end{tabular}
\end{table}

\begin{figure}[htbp]
  \centering
  \begin{subfigure}{0.48\textwidth}
    \includegraphics[width=\textwidth]{sensitivity_km.pdf}
    \caption{okno kilometrine}
  \end{subfigure}\hfill
  \begin{subfigure}{0.48\textwidth}
    \includegraphics[width=\textwidth]{sensitivity_year.pdf}
    \caption{okno letnika}
  \end{subfigure}
  \caption{Občutljivost na širino okna primerljivosti (MdAPE in pokritost).}
  \label{fig:sens-km}
  % vir: analysis/out/sensitivity_km.pdf, sensitivity_year.pdf
\end{figure}

\begin{figure}[htbp]
  \centering
  \includegraphics[width=0.7\textwidth]{sensitivity_km_year_heatmap.pdf}
  \caption{MdAPE [\%] po mreži dovoljenih odstopanj kilometrine in letnika.}
  \label{fig:sens-heatmap}
  % vir: analysis/out/sensitivity_km_year_heatmap.pdf
\end{figure}

\paragraph{Uteži, oblika jedra in popust.}
Tabela~\ref{tab:sens-ostalo} povzema preostale parametre. Razmerje uteži med
kilometrino in letnikom na natančnost praktično ne vpliva ($13{,}4$ do
$13{,}6\,\%$ za vsa preizkušena razmerja), prav tako ne zamenjava linearne
padajoče funkcije z eksponentno ($13{,}4$ proti $13{,}3\,\%$). Edina parametra,
ki opazno premakneta rezultat, sta torej širina letnikovega okna in namerni
popust: ta predznačeno napako proti oglaševani ceni premakne iz $+3{,}8\,\%$ v
$-6{,}8\,\%$, MdAPE pa zviša z $13{,}4$ na $14{,}5\,\%$. Ta premik je posledica
zavestne odločitve, da sistem ocenjuje pričakovano prodajno in ne oglaševane
cene, in ne slabše ocene.

\begin{table}[htbp]
  \centering
  \caption{Občutljivost na uteži, obliko jedra in popust.}
  \label{tab:sens-ostalo}
  % vir: analysis/out/sensitivity_weights.csv, sensitivity_kernel.csv, sensitivity_discount.csv
  \begin{tabularx}{\textwidth}{@{}l X r >{\raggedleft\arraybackslash}p{2.3cm}@{}}
    \toprule
    \textbf{Parameter} & \textbf{Nastavitev} & \textbf{MdAPE [\%]} & \textbf{Predznačena napaka [\%]} \\
    \midrule
    uteži (km / letnik) & $0{,}5 / 0{,}5$              & 13{,}4 & $+3{,}7$ \\
    uteži (km / letnik) & $0{,}6 / 0{,}4$ (privzeto)   & 13{,}4 & $+3{,}8$ \\
    uteži (km / letnik) & $0{,}8 / 0{,}2$              & 13{,}6 & $+3{,}8$ \\
    \midrule
    jedro               & linearno (privzeto)         & 13{,}4 & $+3{,}8$ \\
    jedro               & eksponentno                 & 13{,}3 & $+3{,}8$ \\
    \midrule
    popust              & brez                        & 13{,}4 & $+3{,}8$ \\
    popust              & stopenjski (privzeto)        & 14{,}5 & $-6{,}8$ \\
    \bottomrule
  \end{tabularx}
\end{table}

\paragraph{Sklep o robustnosti.}
Algoritem je robusten: na večino svojih konstant se skoraj ne odziva, kar
pomeni, da njihova hevristična izbira ni sporna. Njegova natančnost je odvisna
predvsem od kakovosti in gostote podatkov (razdelek~\ref{sec:natancnost}), ne
od finega nastavljanja parametrov.

% =====================================================================
\section{Latenca posameznega vrednotenja}
\label{sec:latenca}

Merimo latenco vrste (a) in (b) iz razdelka~\ref{sec:mere}. Meritev izvede
\texttt{benchmarks/bench\_algo.py}: $2000$ neposrednih klicev funkcije nad
naključnimi realnimi vhodi (brez HTTP), po $100$ ogrevalnih klicih, ki se ne
merijo. Ločeno merimo čas poizvedbe SQL, čas točkovanja v Pythonu in čas
serializacije odgovora.

\paragraph{Ozko grlo: pregled celotne tabele.}
Prvotna izvedba poizvedbe ni imela ustreznega indeksa.
Izsek~\ref{lst:explain-before} prikazuje izpis \texttt{EXPLAIN (ANALYZE,
BUFFERS)}: načrt je \texttt{Seq Scan}, baza pri vsakem vrednotenju prebere
celotno tabelo in zavrže $85.273$ vrstic, ki ne ustrezajo pogojem
(\texttt{Rows Removed by Filter}). Ker so podatki v celoti v predpomnilniku
(\texttt{shared hit}, brez \texttt{read}), je čas stabilen, a raste linearno z
velikostjo tabele.

\begin{center}
\captionof{listing}{\texttt{EXPLAIN (ANALYZE, BUFFERS)} pred dodajanjem indeksa.}
\label{lst:explain-before}
\begin{minted}{text}
Seq Scan on normalized_market_listing
    (cost=0.00..3314.15 rows=12 width=97)
    (actual time=0.555..6.847 rows=573 loops=1)
  Filter: ((kilometers >= 18500) AND (kilometers <= 98500)
       AND (registration_year >= 2019) AND (registration_year <= 2025)
       AND (brand_name = 'Citroën') AND (model_name = 'C3'))
  Rows Removed by Filter: 85273
  Buffers: shared hit=1168
Execution Time: 6.880 ms
\end{minted}
\end{center}

\paragraph{Učinek sestavljenega indeksa.}
Po dodajanju sestavljenega indeksa na
$(\text{brand\_name},\ \text{model\_name},\ \text{registration\_year},\
\text{kilometers})$ se načrt spremeni v \texttt{Bitmap Index Scan}
(izsek~\ref{lst:explain-after}): baza prebere le nekaj strani indeksa in
pripadajoče vrstice, čas izvajanja pa pade z $6{,}88$~ms na $0{,}69$~ms.
Latenca poizvedbe SQL se v mediani zmanjša z $9{,}7$~ms na $2{,}0$~ms
(tabela~\ref{tab:indeks-pred-po}, slika~\ref{fig:indeks-pred-po}). Pomembneje
od trenutnega pohitrenja je, da postane čas poizvedbe neodvisen od velikosti
tabele: namesto $O(\text{tabela})$ postane $O(\text{zadetki})$.

\begin{table}[htbp]
  \centering
  \caption{Latenca poizvedbe SQL pred in po dodajanju indeksa, v milisekundah.}
  \label{tab:indeks-pred-po}
  % vir: benchmarks/report.md (results_algo_before/after.json)
  \begin{tabular}{lrrrr}
    \toprule
    & \textbf{p50} & \textbf{p95} & \textbf{p99} & \textbf{Povprečje} \\
    \midrule
    pred (Seq Scan)  & 9{,}7 & 21{,}0 & 24{,}8 & 11{,}3 \\
    po (Index Scan)  & 2{,}0 & 8{,}6  & 12{,}5 & 3{,}0 \\
    \bottomrule
  \end{tabular}
\end{table}

\begin{figure}[htbp]
  \centering
  \includegraphics[width=0.6\textwidth]{latency_index_before_after.pdf}
  \caption{Latenca poizvedbe SQL pred in po dodajanju sestavljenega indeksa.}
  \label{fig:indeks-pred-po}
  % vir: analysis/out/latency_index_before_after.pdf
\end{figure}

\begin{center}
\captionof{listing}{\texttt{EXPLAIN (ANALYZE, BUFFERS)} po dodajanju indeksa.}
\label{lst:explain-after}
\begin{minted}{text}
Bitmap Heap Scan on normalized_market_listing
    (cost=5.84..50.49 rows=12 width=97)
    (actual time=0.202..0.615 rows=573 loops=1)
  Recheck Cond: (brand_name = 'Citroën' AND model_name = 'C3'
       AND registration_year BETWEEN 2019 AND 2025
       AND kilometers BETWEEN 18500 AND 98500)
  Buffers: shared hit=195
  ->  Bitmap Index Scan on idx_nml_valuation
        (actual time=0.166..0.166 rows=573 loops=1)
        Buffers: shared hit=6
Execution Time: 0.688 ms
\end{minted}
\end{center}

\paragraph{Razčlenitev latence po komponentah.}
Tabela~\ref{tab:latenca-vrednotenja} in slika~\ref{fig:latency-breakdown}
prikazujeta latenco celotnega vrednotenja po dodajanju indeksa. Tipično
vrednotenje traja $2{,}0$~ms, v najslabšem odstotku primerov (p99) pa
$13{,}5$~ms. Skoraj celotna latenca ($\approx 88\,\%$ povprečja) odpade na
poizvedbo SQL; točkovanje v Pythonu ($0{,}4$~ms povprečno) in serializacija
odgovora ($\approx 0{,}01$~ms) sta zanemarljiva. Poizvedba v povprečju vrne
$287$ primerljivih vozil (mediana $175$, največ $1495$).

\begin{table}[htbp]
  \centering
  \caption{Latenca posameznega vrednotenja po dodajanju indeksa
  ($N = 2000$), v milisekundah.}
  \label{tab:latenca-vrednotenja}
  % vir: benchmarks/report.md (results_algo.json)
  \begin{tabular}{lrrrr}
    \toprule
    \textbf{Komponenta} & \textbf{p50} & \textbf{p95} & \textbf{p99} & \textbf{Povprečje} \\
    \midrule
    Skupaj                & 2{,}0  & 9{,}7 & 13{,}5 & 3{,}1 \\
    Poizvedba SQL         & 1{,}8  & 8{,}4 & 11{,}5 & 2{,}7 \\
    Točkovanje (Python)   & 0{,}2  & 1{,}2 & 1{,}8  & 0{,}4 \\
    Serializacija (JSON)  & 0{,}01 & 0{,}02 & 0{,}03 & 0{,}01 \\
    \bottomrule
  \end{tabular}
\end{table}

\begin{figure}[htbp]
  \centering
  \includegraphics[width=0.6\textwidth]{latency_breakdown.pdf}
  \caption{Razčlenitev latence posameznega vrednotenja po komponentah.}
  \label{fig:latency-breakdown}
  % vir: analysis/out/latency_breakdown.pdf
\end{figure}

Vrednotenje se torej izvede daleč pod nefunkcionalno zahtevo $2$~s
(razdelek~\ref{sec:zahteve}); prostora za rast podatkovne baze je ob
prisotnosti indeksa dovolj.

% =====================================================================
\section{Obnašanje pod obremenitvijo}
\label{sec:obremenitev}

Merimo end-to-end latenco (vrsta (c)) in prepustnost končne točke
\texttt{POST /valuate} pri naraščajoči sočasnosti. Obremenitev ustvarja
skripta \texttt{benchmarks/loadtest.py} z več vzporednimi odjemalci, ki v
fiksnem $20$-sekundnem oknu neprekinjeno pošiljajo zahtevke z naključnimi
realnimi vhodi. Test se samodejno prekine ob deležu napak nad $5\,\%$ ali
$p_{95}$ nad $5$~s. Med testom vzorčimo tudi zasedenost baze
(\texttt{pg\_stat\_activity}).

\begin{table}[htbp]
  \centering
  \caption{Prepustnost in latenca pri naraščajoči sočasnosti.
  Latence so v milisekundah, $L$ je ocena po Littlovem zakonu.}
  \label{tab:obremenitev}
  % vir: benchmarks/results_load.json (loadtest.py --db-sample)
  \begin{tabular}{r r r r r r r}
    \toprule
    \textbf{Sočasnost} & \textbf{Prep.\ [z/s]} & \textbf{p50} & \textbf{p95} & \textbf{p99}
      & \textbf{Napake} & \textbf{$L$} \\
    \midrule
    1  & \res{x} & \res{x} & \res{x} & \res{x} & \res{x} & \res{x} \\
    5  & \res{x} & \res{x} & \res{x} & \res{x} & \res{x} & \res{x} \\
    10 & \res{x} & \res{x} & \res{x} & \res{x} & \res{x} & \res{x} \\
    20 & \res{x} & \res{x} & \res{x} & \res{x} & \res{x} & \res{x} \\
    40 & \res{x} & \res{x} & \res{x} & \res{x} & \res{x} & \res{x} \\
    \bottomrule
  \end{tabular}
\end{table}

\paragraph{Rezultat.}
\res{Prepustnost narašča približno linearno do sočasnosti X, nato se ustali pri
približno X zahtevkih na sekundo. Vse latence ostanejo pod nefunkcionalno
zahtevo 2~s.}

\paragraph{Skladnost meritve (Littlov zakon).}
\res{Stolpec $L$ (prepustnost krat povprečna latenca) ostaja blizu nastavljene
sočasnosti, kar potrjuje notranjo skladnost meritve.}

\paragraph{Dokaz ozkega grla.}
Teoretični strop postavitve je $1000 / S \approx 324$ vrednotenj na sekundo,
kjer je $S = 3{,}1$~ms povprečni čas ene evalvacije. Dejanska meja je nižja,
ker ves promet serializira \emph{ena sama deljena povezava do baze}
(\texttt{app/database.py}) in en sam strežniški proces.
\res{To potrjuje vzorčenje baze med testom: število aktivnih poizvedb kljub
naraščajoči sočasnosti odjemalcev ne preseže X.}

% =====================================================================
\section{Preverjanje funkcionalnih zahtev}
\label{sec:funkcionalne}

Tabela~\ref{tab:sledljivost} povezuje vsako funkcionalno zahtevo iz
razdelka~\ref{sec:zahteve} z načinom preverjanja in izidom.

\begin{table}[htbp]
  \centering
  \caption{Sledljivost funkcionalnih zahtev do preverjanja.}
  \label{tab:sledljivost}
  % vir: ročni scenariji + imena Jest testov v car-stock-market-be
  \begin{tabularx}{\textwidth}{@{}X X c@{}}
    \toprule
    \textbf{Zahteva} & \textbf{Način preverjanja} & \textbf{Izid} \\
    \midrule
    Registracija uporabnika              & \res{test / scenarij} & \res{OK} \\
    Prijava in odjava                    & \res{test / scenarij} & \res{OK} \\
    Objava oglasa (le registrirani)      & \res{test / scenarij} & \res{OK} \\
    Vrednotenje vozila po atributih      & \res{test / scenarij} & \res{OK} \\
    Všečkanje oglasov                    & \res{test / scenarij} & \res{OK} \\
    Ažuriranje lastnih/všečkanih oglasov & \res{test / scenarij} & \res{OK} \\
    Pregled oglasov (anonimni)           & \res{test / scenarij} & \res{OK} \\
    \bottomrule
  \end{tabularx}
\end{table}

\paragraph{Pokritost testov.}
Avtomatizirani testi zalednega sistema (Jest + Supertest,
razdelek~\ref{sec:testiranje-be}) obsegajo \res{X} testov v \res{X} zbirkah in
pokrivajo \res{X}\,\% stavkov ter \res{X}\,\% vej. Testi se izvedejo ob vsakem
zahtevku za združitev na cevovodu neprekinjene integracije
(razdelek~\ref{subsec:ci}).

% =====================================================================
\section{Preverjanje nefunkcionalnih zahtev}
\label{sec:nefunkcionalne}

Tabela~\ref{tab:nfr} za vsako nefunkcionalno zahtevo iz
razdelka~\ref{sec:zahteve} navaja merilo, izmerjeno vrednost in sodbo.

\begin{table}[htbp]
  \centering
  \caption{Preverjanje nefunkcionalnih zahtev.}
  \label{tab:nfr}
  \setlength{\tabcolsep}{4pt}
  \begin{tabularx}{\textwidth}{@{}>{\raggedright\arraybackslash}X
                                 >{\raggedright\arraybackslash}X
                                 >{\raggedright\arraybackslash}X c@{}}
    \toprule
    \textbf{Zahteva} & \textbf{Merilo} & \textbf{Izmerjeno} & \textbf{Sodba} \\
    \midrule
    Odziv vmesnika na dejanje uporabnika & $< 100$~ms & \res{x}~ms  & \res{OK} \\
    Zahtevki HTTP (razen vrednotenja)    & $< 300$~ms & \res{x}~ms ($p_{95}$) & \res{OK} \\
    Vrednotenje vozila                   & $< 2$~s    & $13{,}5$~ms ($p_{99}$, brez omrežja) & OK \\
    Vizualna privlačnost / dostopnost    & Lighthouse & \res{x} / \res{x} & \res{OK} \\
    Prilagodljiva postavitev             & preizkus na več velikostih zaslona & \res{da} & \res{OK} \\
    Varovanje uporabniških podatkov      & bcrypt, JWT z rokom & \res{opis} & \res{OK} \\
    Odzivnost pri zahtevnejših opravilih & asinhrona obdelava & \res{opis} & \res{OK} \\
    \bottomrule
  \end{tabularx}
\end{table}

\paragraph{Meritev odzivnosti vmesnika in endpointov.}
\res{Odzivnost uporabniškega vmesnika smo izmerili z orodjem Lighthouse na treh
ključnih straneh; ocene za performanco so X, za dostopnost pa X. Latenco
glavnih endpointov zalednega sistema smo izmerili z orodjem autocannon:
$p_{95}$ je v vseh primerih pod X~ms, torej pod zahtevo 300~ms.}

\paragraph{Varnost.}
Gesla so pred shranjevanjem zgoščena z algoritmom bcrypt
(razdelek~\ref{sec:avtentikacija}), žetoni JWT imajo rok veljavnosti $7$~dni,
poti vmesnika API pa ščiti sredinski preverjevalnik žetonov. \res{Orodje
\texttt{npm audit} ob zadnjem pregledu ni javilo kritičnih ranljivosti v
odvisnostih.}

% =====================================================================
\section{Uporabniški preizkus}
\label{sec:uporabniki}

Izvedli smo formativni uporabniški preizkus~\cite{nielsen1994} z namenom
odkriti težave uporabnosti, ne pa statistično ovrednotiti izdelka.

\paragraph{Udeleženci in postopek.}
V preizkus smo vključili \res{N} udeležencev iz kroga družine in bližnjih
prijateljev, katerih starost in računalniška pismenost sta se med seboj
razlikovali (tabela~\ref{tab:udelezenci}). Vsak udeleženec je prejel enak
seznam nalog: registracija, objava oglasa, všečkanje oglasa, izbira
slovenskega jezika v nastavitvah in vrednotenje poljubnega vozila. Med
izvajanjem je udeleženec razmišljal na glas, opazovalec pa je beležil, ali je
nalogo dokončal samostojno, in vsako opaženo težavo.

\begin{table}[htbp]
  \centering
  \caption{Udeleženci uporabniškega preizkusa.}
  \label{tab:udelezenci}
  \begin{tabular}{c c c l}
    \toprule
    \textbf{\#} & \textbf{Starost (razpon)} & \textbf{Rač.\ pismenost (1--5)} & \textbf{Naprava} \\
    \midrule
    1 & \res{x} & \res{x} & \res{x} \\
    2 & \res{x} & \res{x} & \res{x} \\
    3 & \res{x} & \res{x} & \res{x} \\
    \res{\dots} & & & \\
    \bottomrule
  \end{tabular}
\end{table}

\paragraph{Rezultati.}
Vsi udeleženci so vse naloge dokončali brez pomoči, kar potrjuje, da so glavne
poti v vmesniku dovolj jasne (RV4). Kljub temu so odkrili nekaj
pomanjkljivosti; tabela~\ref{tab:tezave} jih razvršča po resnosti (lestvica
$0$--$4$ po Nielsenu~\cite{nielsen1994}) in navaja, ali so bile odpravljene.

\begin{table}[htbp]
  \centering
  \caption{Težave, odkrite med uporabniškim preizkusom
  (resnost po lestvici $0$--$4$).}
  \label{tab:tezave}
  \setlength{\tabcolsep}{5pt}
  \begin{tabularx}{\textwidth}{@{}>{\raggedright\arraybackslash}X c c@{}}
    \toprule
    \textbf{Težava} & \textbf{Resnost} & \textbf{Status} \\
    \midrule
    Premajhni naslovi podstrani v glavi glavne strani       & \res{x} & odpravljeno \\
    Manjkajoče polje za ponovni vnos gesla pri registraciji & \res{x} & odpravljeno \\
    Velika števila brez ločil za tisočice                   & \res{x} & odpravljeno \\
    Manjkajoč opis pod slikami oblik vozil v obrazcu        & \res{x} & odpravljeno \\
    Premalo vozil, dostopnih za vrednotenje                 & \res{x} & \res{v načrtu} \\
    \bottomrule
  \end{tabularx}
\end{table}

Omejitve preizkusa (majhen in priložnostni vzorec, brez standardiziranega
vprašalnika, brez merjenja časa po nalogah) obravnavamo v
razdelku~\ref{sec:groznje}.

% =====================================================================
\section{Simulacija uporabe}
\label{sec:simulacija}

Za ponazoritev delovanja rešitve povzemamo tipičen potek uporabe. Uporabnik je
ob vstopu usmerjen na glavno stran (slika~\ref{fig:glavna-stran}), kjer sta
izpostavljeni prijava in registracija ter pregled novih in vseh objavljenih
oglasov. Registriran uporabnik lahko objavi oglas prek vodene forme
(slika~\ref{fig:objava-oglasa}), ovrednoti poljubno vozilo, všečka oglase in si
ogleda svoj profil z všečkanimi vozili (slika~\ref{fig:liked-cars}).

\begin{figure}[H]
  \centering
  
\includegraphics[width=0.9\textwidth]{glavna-stran.png}
  \caption{Glavna stran prototipa}
  \label{fig:glavna-stran}
\end{figure}
\begin{figure}[H]
  \centering
  \includegraphics[width=0.9\textwidth]{objava-oglasa.png}
  \caption{Vodena forma za objavo oglasa}
  \label{fig:objava-oglasa}
\end{figure}
\begin{figure}[H]
  \centering
  \includegraphics[width=0.9\textwidth]{liked-cars.png}
  \caption{Stran uporabnika s prikazanimi všečkanimi vozili}
  \label{fig:liked-cars}
\end{figure}

\paragraph{Primerjava ocen z obstoječimi portali.}
Za ilustracijo smo delovanje sistema primerjali s ponudbo na dveh spletnih
portalih na treh vozilih (tabela~\ref{tab:vrednotenje}). Sistematično in
kvantitativno oceno natančnosti podaja razdelek~\ref{sec:natancnost}; ta
primerjava je zgolj kvalitativna ponazoritev.

\begin{table}[htbp]
  \centering
  \caption{Primerjava ocenjenih cen s ponudbo na spletnih portalih (v evrih).}
  \label{tab:vrednotenje}
  \setlength{\tabcolsep}{5pt}
  \begin{tabular}{lccc}
    \toprule
    & \textbf{Audi A3} & \textbf{BMW ser.\ 2} & \textbf{Citroën Berlingo} \\
    \midrule
    slov.\ portal (razpon)     & 4.000--6.000  & 9.000--12.000  & 7.500--10.000 \\
    nem.\ portal (razpon)      & 4.500--9.000  & pribl.\ 11.100 & 8.000--12.000 \\
    \midrule
    Sistem --- dobra cena      & 4.500  & 7.600   & 7.600  \\
    Sistem --- sredinska cena  & 6.350  & 9.300   & 9.941  \\
    Sistem --- slaba cena      & 7.600  & 10.500  & 12.100 \\
    Raznolikost                & $49\,\%$ & $31\,\%$ & $45\,\%$ \\
    \bottomrule
  \end{tabular}
\end{table}

Pri vozilih z nizko raznolikostjo ponudbe (Audi~A3, BMW serije~2) se sredinska
ocena dobro ujema s cenami na obeh portalih. Pri modelu Citroën Berlingo je
raznolikost bistveno višja, zato je razpon med dobro in slabo oceno širši,
sredinska ocena pa manj zanesljiva --- kar je skladno z ugotovitvijo iz
razdelka~\ref{sec:natancnost}, da napaka strmo narašča z raznolikostjo ponudbe.
Pri interpretaciji je treba upoštevati, da sistem ocenjuje pričakovano prodajno
in ne oglaševano ceno.

% =====================================================================
\section{Grožnje veljavnosti}
\label{sec:groznje}

\paragraph{Notranja veljavnost.}
Kot merilo resnice pri merjenju natančnosti uporabljamo \emph{oglaševano} ceno,
ne dejanske prodajne cene, ki ni javno dostopna. Namerni popust, s katerim
sistem cilja pričakovano prodajno ceno, je zato ocenjen in ne umerjen na
dejanskih prodajah. Isti fizični oglas se lahko pojavi pri več virih; prečno
preverjanje izloči zapise z istim izvornim oglasom, ne more pa zaznati
podvojitev med viri.

\paragraph{Zunanja veljavnost.}
Podatki pokrivajo $10$ znamk in pretežno slovenski trg (z manjšim nemškim
vzorcem); znamka z največ oglasi na trgu (okoli $8000$ oglasov) v šifrantu ni
zajeta in je iz analize izpuščena. Uporabniški preizkus je temeljil na majhnem
in priložnostnem vzorcu.

\paragraph{Konstruktna veljavnost.}
Pojem \enquote{tržna vrednost} operacionaliziramo kot uteženo povprečje
oglaševanih cen primerljivih vozil, znižano za pričakovani pogajalski popust;
\enquote{raznolikost ponudbe} pa kot relativno širino interkvartilnega razpona
teh cen. Drugačna operacionalizacija bi lahko dala drugačne rezultate.

\paragraph{Zanesljivost.}
Meritve latence in obremenitve so bile izvedene na deljenem produkcijskem
strežniku, kjer lahko na rezultate vpliva ozadje. Meritve natančnosti so
ponovljive nad fiksnim posnetkom baze.

% =====================================================================
\section{Povzetek ugotovitev}
\label{sec:povzetek-analize}

Tabela~\ref{tab:rv-povzetek} povzema odgovore na raziskovalna vprašanja glede
na vnaprej postavljena merila.

\begin{table}[htbp]
  \centering
  \caption{Povzetek odgovorov na raziskovalna vprašanja.}
  \label{tab:rv-povzetek}
  \begin{tabularx}{\textwidth}{@{}l X c@{}}
    \toprule
    \textbf{RV} & \textbf{Glavni rezultat} & \textbf{Sodba} \\
    \midrule
    RV1 natančnost & MdAPE $13{,}5\,\%$ ($9$--$10\,\%$ za dobro pokrita vozila),
                     bistveno bolje od naivne ocene ($37{,}8\,\%$), primerljivo z
                     drugimi metodami & izpolnjeno \\
    RV2 odvisnost  & napaka pada z velikostjo vzorca (ustali se pri $n\approx20$),
                     narašča s starostjo, kilometrino in raznolikostjo ponudbe
                     ($7{,}5\,\%$ do $29{,}7\,\%$) & pojasnjeno \\
    RV3 hitrost    & vrednotenje $2$~ms (p50), $13{,}5$~ms (p99); indeks odpravi
                     odvisnost od velikosti baze; \res{meja X~z/s} & \res{izpolnjeno} \\
    RV4 uporabnost & vse naloge dokončane brez pomoči; \res{X} težav odpravljenih & izpolnjeno \\
    \bottomrule
  \end{tabularx}
\end{table}

\paragraph{Sodba o zastavljenih izhodiščih.}
Rezultati \emph{potrjujejo} izhodišča, ki smo si jih zastavili na začetku. Sistem na podlagi štirih
atributov vozila vrne oceno vrednosti z mediano odstotne napake $13{,}5\,\%$
(brez namernega popusta), pri dobro pokritih vozilih pa $9$--$10\,\%$, kar je
znotraj zastavljenega praga $15\,\%$ in bistveno bolje od naivne metode.
Predlagani pristop se med primerljivo natančnimi metodami ne uvršča na prvo
mesto --- preprost $k$-najbližjih sosedov doseže $11{,}2\,\%$ --- je pa edini,
ki poleg točkovne ocene vrne tudi razložljiv cenovni razpon in indeks
zaupanja, ta pa se izkaže za dober pokazatelj dejanske zanesljivosti
posamezne ocene. Glavne omejitve so odvisnost od enega prevladujočega
podatkovnega vira, nižja natančnost pri starejših vozilih in modelih z
razpršeno ponudbo ter merilo resnice, ki temelji na oglaševanih in ne
prodajnih cenah. Konkretne smeri za izboljšave --- vključno s prehodom na
metode strojnega učenja, ki jih ta analiza dodatno utemeljuje --- navajamo v
poglavju~\ref{ch:sklep}.

%
%  Številke so vpisane iz izhodov skript v CarStockMarketEvaluation/analysis/.
%  Preostale oznake \res{...} (rdeče) je treba še izpolniti:
%    - obremenitveni test (benchmarks/loadtest.py)
%    - sledljivost testov, Jest coverage, Lighthouse
%    - tabeli udeležencev in najdenih težav
%  Slike (PDF) iz analysis/out/ prekopiraj v mapo slike/.
% =====================================================================

\newcommand{\res}[1]{\textcolor{red}{\textbf{[#1]}}}

\chapter{Analiza in vrednotenje rešitve}
\label{ch:analiza}

V tem poglavju sistematično pregleadamo in analiziramo razvito rešitev. Najprej iz zahtev in
izhodišč, ki smo si jih zastavili na začetku, izpeljemo raziskovalna vprašanja in za vsako določimo merilo
uspešnosti (razdelek~\ref{sec:rq}). Nato opišemo merilno okolje in postopke
(razdelek~\ref{sec:metodologija}) ter natančno definiramo vse uporabljene mere
(razdelek~\ref{sec:mere}). Sledijo rezultati: značilnosti zbranih
podatkov~(\ref{sec:podatki}), natančnost ocene
vrednosti~(\ref{sec:natancnost}), občutljivost algoritma na
parametre~(\ref{sec:obcutljivost}), latenca posameznega
vrednotenja~(\ref{sec:latenca}) in obnašanje pod
obremenitvijo~(\ref{sec:obremenitev}). Zatem preverimo funkcionalne in
nefunkcionalne zahteve (\ref{sec:funkcionalne}, \ref{sec:nefunkcionalne}) ter
predstavimo uporabniški preizkus~(\ref{sec:uporabniki}) in simulacijo
uporabe~(\ref{sec:simulacija}). Poglavje sklenemo z obravnavo groženj
veljavnosti~(\ref{sec:groznje}) in povzetkom ugotovitev glede na
zastavljena izhodišča~(\ref{sec:povzetek-analize}).

% =====================================================================
\section{Raziskovalna vprašanja in merila uspešnosti}
\label{sec:rq}

Na začetku diplomskega dela (poglavje~\ref{ch:uvod}) smo si zastavili, da je
mogoče uporabniku na podlagi zbranih tržnih podatkov in primerjave s podobnimi
vozili ponuditi \emph{uporabno} in \emph{dovolj natančno} oceno tržne vrednosti
vozila in da vključitev takega sistema izboljša uporabniško izkušnjo. Iz tega
izhodišča in iz nefunkcionalnih zahtev (razdelek~\ref{sec:zahteve}) izpeljemo
štiri raziskovalna vprašanja (RV1--RV4) in za vsako določimo merilo uspešnosti.

\emph{Mediana absolutne odstotne napake} (MdAPE) je osrednja mera natančnosti v
tem poglavju; nanjo se sklicujejo merila uspešnosti in vsi nadaljnji razdelki.
Izračunamo jo tako, da za vsak testni primer določimo absolutno odstotno napako,
torej za koliko odstotkov se ocenjena cena razlikuje od dejanske, nato pa čez
vse primere vzamemo mediano teh vrednosti. MdAPE $=15\,\%$ tako pomeni, da je pri
polovici vozil ocena znotraj $15\,\%$ od dejanske cene, pri drugi polovici pa je
odstopanje večje; manjša kot je vrednost, natančnejše so ocene. Mediano
uporabljamo namesto povprečja, ker je manj občutljiva na posamezne skrajne
napake. Natančne definicije MdAPE in preostalih mer so v
razdelku~\ref{sec:mere}.

Raziskovalna vprašanja so:

\begin{itemize}
  \item \textbf{RV1 --- natančnost ocene.} Kako natančno sistem oceni tržno
        vrednost vozila v primerjavi z dejanskimi oglasi na trgu?
        \emph{Merilo:} MdAPE je pri vozilih z vsaj $10$ primerljivimi
        vozili manjša od $15\,\%$, ocena pa je bistveno natančnejša od
        preproste (naivne) referenčne metode.
  \item \textbf{RV2 --- odvisnost natančnosti od podatkov.} Od česa je
        natančnost ocene najbolj odvisna (velikost vzorca primerljivih vozil,
        razpršenost ponudbe, segment vozila)?
  \item \textbf{RV3 --- hitrost in razširljivost.} Ali sistem zadošča
        nefunkcionalnim zahtevam glede odzivnosti tudi pod sočasno
        obremenitvijo, in kje je zgornja meja trenutne postavitve?
  \item \textbf{RV4 --- uporabnost.} Ali ciljni uporabniki znajo brez navodil
        opraviti glavne naloge (registracija, objava oglasa, vrednotenje)?
\end{itemize}
% =====================================================================
\section{Metodologija in merilno okolje}
\label{sec:metodologija}

\paragraph{Merilno okolje.}
Meritve hitrosti (razdelka~\ref{sec:latenca} in~\ref{sec:obremenitev}) so
izvedene v \emph{produkcijskem} okolju, torej na navideznem zasebnem strežniku
z dvema procesorskima jedroma in $4$~GB pomnilnika ter s tehnologijami in
njihovimi različicami, opisanimi v poglavju~\ref{ch:razvoj}
(razdelek~\ref{sec:namestitev}). Analitične in merilne skripte tečejo nad isto
podatkovno bazo kot je na produkciji.

\paragraph{Posnetek podatkov.}
Meritve v tem poglavju temeljijo na posnetku podatkovne baze iz septembra
$2026$, ki obsega približno $86.000$ normaliziranih oglasov
(razdelek~\ref{sec:podatki}). Ker se posnetek med analizo ne spreminja, so
rezultati ponovljivi.

\paragraph{Orodja.}
Za meritve in analizo v tem poglavju smo napisali namenske skripte, ločene od
produkcijske kode; ena od njih je dokumentirana kopija ocenjevalnega
algoritma, parametrizirana za eksperimente. Statistično obdelavo in izris
grafov smo izvedli s knjižnicama \texttt{NumPy}~\cite{harris2020numpy} in
\texttt{Matplotlib}~\cite{hunter2007matplotlib}.

\paragraph{Natančnost.}
Natančnost ocene (razdelek~\ref{sec:natancnost}) merimo z naključno izbranim
normaliziranim oglasom: njegove atribute (znamka, model, letnik, kilometrina)
uporabimo kot vhod v algoritem, sam oglas pa izločimo iz primerjalne množice. Nato izračunamo oceno in jo
primerjamo z \emph{dejansko oglaševano} ceno tega oglasa. Postopek ponovimo za
$3000$ naključnih oglasov.

% =====================================================================
\section{Uporabljene mere}
\label{sec:mere}

\paragraph{Latenca}
 je čas, ki preteče od trenutka tik pred oddajo zahtevka do
trenutka, ko odjemalec prejme celoten odgovor. Glede na to, kolikšen del te
poti zajamemo, ločimo tri vrste:
\begin{itemize}
  \item[(a)] \emph{latenca obdelave na strežniku} --- čas izvajanja funkcije
        za vrednotenje, od začetka do konca, brez omrežja;
  \item[(b)] \emph{latenca poizvedbe SQL} --- čas izvajanja poizvedbe v podatkovni bazi,
   od začetka do konca, brez omrežja;
  \item[(c)] \emph{latenca od konca do konca} (angl.\ \emph{end-to-end}) ---
        (a) skupaj z omrežnim obhodom med odjemalcem in strežnikom ter
        obdelavo zahtevka HTTP; to je čas, ki ga zazna uporabnik.
\end{itemize}
V razdelku~\ref{sec:latenca} merimo vrsti (a) in (b), v
razdelku~\ref{sec:obremenitev} pa vrsto (c).

\paragraph{Percentili.}
\textbf{Percentil} $p_k$ je vrednost, pod katero pade $k\,\%$ meritev. Poročamo
$p_{50}$ (mediana), $p_{95}$ in $p_{99}$.

\paragraph{Prepustnost, sočasnost, nasičenje.}
Obnašanje sistema pod obremenitvijo opišemo s tremi količinami:
\begin{itemize}
  \item \textbf{prepustnost} --- število uspešno obdelanih zahtevkov na
        sekundo;
  \item \textbf{sočasnost} --- število odjemalcev, ki hkrati držijo odprt
        zahtevek;
  \item \textbf{nasičenje} (\emph{koleno}) --- točka, v kateri dodatna
        sočasnost ne poveča več prepustnosti, latenca pa začne naraščati
        približno linearno.
\end{itemize}
Notranjo skladnost meritev preverimo z \textbf{Littlovim zakonom}~\cite{little1961},
\begin{equation}
  \label{eq:little}
  L = \lambda \cdot W,
\end{equation}
kjer je $L$ povprečno število zahtevkov v sistemu, $\lambda$ prepustnost in
$W$ povprečna latenca.

\paragraph{Napaka ocene.}
Za oglas $i$ z dejansko ceno $p_i$ in ocenjeno ceno $\hat p_i$ definiramo
absolutno odstotno napako
\begin{equation}
  \mathrm{APE}_i = \frac{\lvert \hat p_i - p_i \rvert}{p_i}.
\end{equation}
\textbf{MAPE} je povprečje $\mathrm{APE}_i$, \textbf{MdAPE} pa njihova mediana;
MdAPE je robustnejša na osamelce in je naša glavna mera~\cite{hyndman2006}.
Poročamo tudi \textbf{MAE} (povprečna absolutna napaka) in \textbf{RMSE}
(koren povprečne kvadratne napake) v evrih ter \emph{predznačeno} odstotno
napako, ki pokaže sistematičen odmik ocene navzgor ali navzdol.

\paragraph{Pokritost.}
\textbf{Pokritost} je delež poizvedb, za katere sistem vrne oceno, torej za
katere najde vsaj eno (oziroma vsaj $5$ oz.\ $10$) primerljivo vozilo.

\paragraph{Raznolikost ponudbe.}
\textbf{Raznolikost} (v izvorni kodi \texttt{ratio}) je relativna širina
cenovnega razpona primerljivih vozil,
\begin{equation}
  \label{eq:raznolikost}
  \mathrm{raznolikost} = \frac{P_{75} - P_{25}}{\hat{p}},
\end{equation}
kjer sta $P_{25}$ in $P_{75}$ (v kodi \texttt{goodPrice} in \texttt{badPrice})
uteženi $25.$ oziroma $75.$ percentil cen primerljivih vozil, $\hat{p}$
(v kodi \texttt{estimatedPrice}) pa ocenjena cena vozila. Vrednost izražamo v
odstotkih; enoten zapis uporabljamo v vseh tabelah tega poglavja.

% =====================================================================
\section{Značilnosti zbranih podatkov}
\label{sec:podatki}
\paragraph{Viri in obseg.}
Velika večina oglasov je bila zajeta z največjega slovenskega spletnega portala
za prodajo rabljenih vozil (razdelek~\ref{sec:zajem-podatkov}); manjši del
predstavlja vzorec z večjega nemškega oglasnega portala, ki smo ga uporabili za
preverjanje delovanja cevovoda na drugačni strukturi vira. Zajem je zajel
javno dostopne oglase za nekomercialno raziskovalno rabo, z omejeno hitrostjo
zahtevkov in brez shranjevanja kontaktnih ali osebnih podatkov prodajalcev.
Po normalizaciji baza obsega $85.846$ oglasov ($83.493$ s slovenskega in
$2.353$ z nemškega portala) za $10$ znamk vozil; razčlenitev po znamki je v
tabeli~\ref{tab:podatki-znamke}.

\begin{table}[htbp]
  \centering
  \caption{Število normaliziranih oglasov po znamki.}
  \label{tab:podatki-znamke}
  % vir: analysis/out/dataset_per_brand.csv
  \begin{tabular}{lr lr}
    \toprule
    \textbf{Znamka} & \textbf{Oglasi} & \textbf{Znamka} & \textbf{Oglasi} \\
    \midrule
    Volkswagen    & 21.191 & Peugeot     & 7.513 \\
    Audi          & 11.842 & Citroën     & 6.729 \\
    BMW           & 10.724 & Ford        & 5.783 \\
    Renault       & 9.615  & Toyota      & 2.648 \\
    Mercedes-Benz & 9.355  & Alfa Romeo  & 446 \\
    \bottomrule
  \end{tabular}
\end{table}

\paragraph{Porazdelitve ključnih atributov.}
Slika~\ref{fig:porazdelitve} prikazuje porazdelitve oglaševanih cen, prevožene
kilometrine in letnika. Porazdelitev cen je izrazito desno asimetrična (dolg
rep dragih vozil), zato jo prikazujemo na logaritemski osi. Mediana cene je
$14.890$~EUR (povprečje $18.137$~EUR), mediana kilometrine $159.950$~km,
mediana letnika pa $2017$; podrobnosti so v tabeli~\ref{tab:podatki-porazdelitve}.
Zaradi berljivosti so na grafih odrezane skrajne vrednosti (nekaj oglasov ima
očitno napačno ceno, kilometrino ali letnik); njihovo število je označeno v
kotu vsakega grafa, celoten razpon pa je razviden iz stolpca \emph{maks.}\ v
tabeli~\ref{tab:podatki-porazdelitve} in iz obravnave osamelcev v
tabeli~\ref{tab:osamelci}.

\begin{table}[htbp]
  \centering
  \caption{Povzetek porazdelitev ključnih numeričnih atributov.}
  \label{tab:podatki-porazdelitve}
  % vir: analysis/out/dataset_distribution_summary.json
  \setlength{\tabcolsep}{5pt}
  \begin{tabular}{lrrrrr}
    \toprule
    & \textbf{$p_{5}$} & \textbf{mediana} & \textbf{povprečje} & \textbf{$p_{95}$} & \textbf{maks.} \\
    \midrule
    cena [EUR]        & 1.550  & 14.890  & 18.137  & 49.014  & 374.000 \\
    kilometrina [km]  & 26.000 & 159.950 & 169.464 & 334.000 & 1.974.001 \\
    letnik            & 2005   & 2017    & 2016    & 2024    & 2029 \\
    moč [kW]          & 56     & 110     & 117     & 221     & 920 \\
    \bottomrule
  \end{tabular}
\end{table}

\begin{figure}[htbp]
  \centering
  \begin{subfigure}{0.48\textwidth}
    \includegraphics[width=\textwidth]{dataset_price_hist.pdf}
    \caption{oglaševana cena (log)}
  \end{subfigure}\hfill
  \begin{subfigure}{0.48\textwidth}
    \includegraphics[width=\textwidth]{dataset_km_hist.pdf}
    \caption{prevoženi kilometri}
  \end{subfigure}
  \\[1ex]
  \begin{subfigure}{0.48\textwidth}
    \includegraphics[width=\textwidth]{dataset_year_hist.pdf}
    \caption{letnik prve registracije}
  \end{subfigure}\hfill
  \begin{subfigure}{0.48\textwidth}
    \includegraphics[width=\textwidth]{dataset_brand_bar.pdf}
    \caption{oglasi po znamki}
  \end{subfigure}
  \caption{Porazdelitve ključnih atributov v normalizirani bazi.}
  \label{fig:porazdelitve}
  % vir: analysis/out/dataset_*_hist.pdf, dataset_brand_bar.pdf
\end{figure}

\paragraph{Pokritost po segmentih.}
Za pošteno oceno mora imeti algoritem dovolj primerljivih vozil v ožjem
segmentu (ista znamka in model, podoben letnik in kilometrina).
Slika~\ref{fig:pokritost} prikazuje toplotno karto za znamki Renault in
Volkswagen: vrstice so modeli, stolpci letniki, barva pa
število oglasov v tej kombinaciji. Vidno je, da so pogosti modeli srednjih
letnikov (npr.\ Volkswagen Golf ali Renault Clio) gosto pokriti, redki modeli
in skrajni letniki pa le redko, kar neposredno pojasni razlike v zanesljivosti
ocen (razdelek~\ref{sec:natancnost}).

\begin{figure}[htbp]
  \centering
  \includegraphics[width=\textwidth]{dataset_coverage_heatmap.pdf}
  \caption{Pokritost: število oglasov po modelu in letniku za znamki
  Renault in Volkswagen (po ena znamka na vrstico).}
  \label{fig:pokritost}
  % vir: analysis/out/dataset_coverage_heatmap.pdf
\end{figure}

\paragraph{Osamelci in kakovost podatkov.}
Postopek normalizacije trenutno ne izvaja filtriranja osamelcev, zato baza
vsebuje tudi vozila z očitno napačnimi vrednostmi.
Tabela~\ref{tab:osamelci} prikazuje njihov delež. Kar $10{,}7\,\%$ oglasov ima
vsaj eno vrednost zunaj smiselnih mej, največ pri kilometrini ($9{,}9\,\%$;
najvišja zabeležena vrednost je skoraj dva milijona kilometrov) in letniku
($8{,}1\,\%$; nekaj oglasov ima celo letnik v prihodnosti). Preprost filter po
interkvartilnem razponu (IQR) na ceni znotraj skupine
$(\text{znamka},\text{model})$ bi odstranil dodatnih $2{,}2\,\%$ oglasov.
Vpliv doslednega filtriranja osamelcev na natančnost obravnavamo v
razdelku~\ref{sec:obcutljivost}, uvedbo takega filtra pa v
poglavju~\ref{ch:sklep}.

\begin{table}[H]
  \centering
  \caption{Osamelci v normalizirani bazi ($85.846$ oglasov).}
  \label{tab:osamelci}
  % vir: analysis/out/dataset_outliers.csv
  \begin{tabularx}{\textwidth}{@{}X r r@{}}
    \toprule
    \textbf{Merilo} & \textbf{Št.\ vrstic} & \textbf{Delež} \\
    \midrule
    cena zunaj $[500,\ 200.000]$~EUR         & 427   & $0{,}5\,\%$ \\
    kilometrina zunaj $[1000,\ 500.000]$~km  & 8.495 & $9{,}9\,\%$ \\
    letnik zunaj $[1990,\ 2027]$             & 6.979 & $8{,}1\,\%$ \\
    vsaj ena vrednost zunaj mej              & 9.169 & $10{,}7\,\%$ \\
    dodatno po filtru IQR na ceni            & 1.878 & $2{,}2\,\%$ \\
    \bottomrule
  \end{tabularx}
\end{table}

% =====================================================================
\section{Natančnost ocene vrednosti}
\label{sec:natancnost}

To je osrednja meritev diplomskega dela: neposredno preverja zastavljeno
izhodišče o \emph{dovolj natančni} oceni (RV1). Protokol je opisan v
razdelku~\ref{sec:metodologija}, rezultate pa izdela skripta
\texttt{analysis/accuracy\_cv.py} nad vzorcem $3000$ oglasov (od tega jih je
$2986$ dobilo oceno).

\paragraph{Merilo resnice in namerni odmik ocene.}
Kot merilo resnice uporabljamo \emph{oglaševano} ceno, saj dejanske prodajne
cene niso javno dostopne. Sistem pa namenoma ne ocenjuje oglaševane cene,
temveč \emph{pričakovano prodajno ceno}: ker se rabljena vozila v pogajanjih
praviloma prodajo pod oglaševano ceno, oceno zniža s stopenjskim popustom
($9$--$15\,\%$ glede na cenovni razred). Zato je pri primerjavi z oglaševano
ceno pričakovan rahel \emph{negativni} odmik ocene, kar je zaželeno in ne
pomeni slabše ocene. To lastnost upoštevamo pri branju vseh rezultatov v
nadaljevanju: napako merimo dvakrat --- proti oglaševani ceni z vključenim
popustom (kar dejansko vrne sistem) in brez popusta (čista primerjalna
ocena).

\paragraph{Skupni rezultat.}
Brez popusta znaša MdAPE predlaganega pristopa $13{,}5\,\%$, MAPE $23{,}8\,\%$,
MAE $2798$~EUR in RMSE $5151$~EUR; z vključenim namernim popustom je MdAPE
proti oglaševani ceni $14{,}5\,\%$ (mediana predznačene napake $-6{,}8\,\%$).
Pokritost je $99{,}5\,\%$ pri $n \ge 1$ in $95{,}7\,\%$ pri $n \ge 10$.
Slika~\ref{fig:pred-vs-actual} prikazuje napovedano ceno proti dejanski
($R^2 = 0{,}87$), slika~\ref{fig:ape-hist} pa porazdelitev odstotne napake.

\begin{figure}[htbp]
  \centering
  \begin{subfigure}{0.48\textwidth}
    \includegraphics[width=\textwidth]{accuracy_pred_vs_actual.pdf}
    \caption{napovedano proti dejanskemu}
    \label{fig:pred-vs-actual}
  \end{subfigure}\hfill
  \begin{subfigure}{0.48\textwidth}
    \includegraphics[width=\textwidth]{accuracy_ape_hist.pdf}
    \caption{porazdelitev napake}
    \label{fig:ape-hist}
  \end{subfigure}
  \caption{Natančnost ocene vrednosti (leave-one-out, brez popusta).}
  % vir: analysis/out/accuracy_pred_vs_actual.pdf, accuracy_ape_hist.pdf
\end{figure}

\paragraph{Primerjava z osnovnimi modeli.}
Da presodimo, ali je predlagani pristop boljši od trivialnega, ga primerjamo s
štirimi referenčnimi modeli na istem protokolu in isti množici kandidatov
(tabela~\ref{tab:baselines}, slika~\ref{fig:baselines}): B0~globalna mediana
cene za dani model, B1~navadno povprečje cen v filtrirnem oknu, B1m~mediana cen
v oknu, B2~linearna regresija $\text{cena} \sim \text{letnik} + \text{km}$,
B3~metoda $k$~najbližjih sosedov ($k = 10$) po standardiziranih letniku in
kilometrini, ter B4~naš pristop (uteženo povprečje s podobnostnim jedrom, brez
popusta).

\begin{table}[htbp]
  \centering
  \caption{Primerjava natančnosti z osnovnimi modeli (leave-one-out).
  MdAPE in MAPE sta v odstotkih, MAE in RMSE v evrih.}
  \label{tab:baselines}
  % vir: analysis/out/accuracy_overall.csv
  \setlength{\tabcolsep}{5pt}
  \begin{tabular}{l r r r r}
    \toprule
    \textbf{Model} & \textbf{MdAPE} & \textbf{MAPE} & \textbf{MAE} & \textbf{RMSE} \\
    \midrule
    B0 globalna mediana       & 37{,}8 & 95{,}7 & 7395 & 10883 \\
    B1 povprečje v oknu       & 14{,}2 & 24{,}9 & 2930 & 5333 \\
    B1m mediana v oknu        & 13{,}1 & 22{,}5 & 2877 & 5561 \\
    B2 linearna regresija     & 12{,}9 & 22{,}4 & 2677 & 5032 \\
    B3 $k$-najbližjih sosedov & \textbf{11{,}2} & 21{,}5 & 2497 & 4883 \\
    B4 naš pristop            & 13{,}5 & 23{,}8 & 2798 & 5151 \\
    \bottomrule
  \end{tabular}
\end{table}

\begin{figure}[htbp]
  \centering
  \includegraphics[width=0.7\textwidth]{accuracy_baselines.pdf}
  \caption{MdAPE osnovnih modelov in predlaganega pristopa.}
  \label{fig:baselines}
  % vir: analysis/out/accuracy_baselines.pdf
\end{figure}

Predlagani pristop je \emph{bistveno} boljši od naivne globalne mediane
($13{,}5\,\%$ proti $37{,}8\,\%$) in primerljiv s preostalimi metodami, ki se
gibljejo med $11$ in $14\,\%$. Ni pa najboljši: metoda $k$~najbližjih sosedov
doseže $11{,}2\,\%$, tudi navadna mediana istega okna ($13{,}1\,\%$) je
malenkost natančnejša. To je pričakovano za zasnovo, katere cilj je
\emph{razložljivost}: za razliko od naštetih metod predlagani pristop poleg
točkovne ocene vrne tudi cenovni razpon in indeks zaupanja, uteži pa so
neposredno interpretabilne. Rezultat hkrati potrjuje smiselnost prehoda na
metode strojnega učenja, ki jo predlagamo v poglavju~\ref{ch:sklep}: če že
preprost $k$-NN preseže sedanji pristop, je od zmogljivejših modelov mogoče
pričakovati še boljše rezultate.

\paragraph{Odvisnost od velikosti vzorca in segmenta (RV2).}
Slika~\ref{fig:mdape-n} prikazuje MdAPE glede na število primerljivih vozil
$n$: napaka je pri $n < 10$ okoli $20\,\%$, pri $n \ge 20$ pa se ustali pri
približno $13\,\%$. To empirično podpira izbiro konstante $20$ v indeksu
zaupanja (enačba~\ref{eq:confidence}) --- pri toliko primerljivih vozilih je
ocena že blizu svoje najboljše natančnosti. Slika~\ref{fig:mdape-brand} prikazuje MdAPE po znamki: najbolje se
odrežejo pogosti, gosto pokriti modeli (Audi $9{,}8\,\%$, Toyota $10{,}8\,\%$,
Volkswagen $12{,}1\,\%$), najslabše pa Mercedes-Benz ($24{,}9\,\%$), pri
katerem se pri normalizaciji več različnih modelov zlije v isto kanonično
ime, kar poslabša izbor primerljivih vozil.
Tabela~\ref{tab:natancnost-segmenti} kaže še razčlenitev po starosti,
kilometrini in raznolikosti ponudbe.

\begin{figure}[htbp]
  \centering
  \begin{subfigure}{0.48\textwidth}
    \includegraphics[width=\textwidth]{accuracy_mdape_vs_n.pdf}
    \caption{glede na velikost vzorca $n$}
    \label{fig:mdape-n}
  \end{subfigure}\hfill
  \begin{subfigure}{0.48\textwidth}
    \includegraphics[width=\textwidth]{accuracy_mdape_by_brand.pdf}
    \caption{po znamki}
    \label{fig:mdape-brand}
  \end{subfigure}
  \caption{Odvisnost natančnosti ocene od podatkov.}
  % vir: analysis/out/accuracy_mdape_vs_n.pdf, accuracy_mdape_by_brand.pdf
\end{figure}

\begin{table}[htbp]
  \centering
  \caption{Natančnost predlaganega pristopa po segmentih (MdAPE).}
  \label{tab:natancnost-segmenti}
  % vir: analysis/out/accuracy_by_age.csv, accuracy_by_km.csv, accuracy_by_ratio.csv
  \begin{tabularx}{\textwidth}{@{}l X r r@{}}
    \toprule
    \textbf{Razčlenitev} & \textbf{Koš} & \textbf{Primeri} & \textbf{MdAPE [\%]} \\
    \midrule
    starost [leta] & 0--3   & 367 & 9{,}0 \\
                   & 4--7   & 921 & 10{,}6 \\
                   & 8--12  & 932 & 14{,}5 \\
                   & $13+$  & 766 & 22{,}9 \\
    \midrule
    kilometrina    & $<50.000$        & 280  & 9{,}6 \\
                   & $50.000$--$100.000$  & 449  & 11{,}9 \\
                   & $100.000$--$150.000$ & 634  & 12{,}4 \\
                   & $150.000$--$200.000$ & 629  & 13{,}7 \\
                   & $>200.000$       & 1021 & 16{,}7 \\
    \midrule
    raznolikost    & $0$--$0{,}10$    & 88   & 7{,}5 \\
                   & $0{,}10$--$0{,}25$ & 1002 & 8{,}4 \\
                   & $0{,}25$--$0{,}50$ & 1613 & 15{,}5 \\
                   & $>0{,}50$        & 447  & 29{,}7 \\
    \bottomrule
  \end{tabularx}
\end{table}

Vzorec pokaže jasne zakonitosti: napaka narašča s starostjo vozila ($9\,\%$ pri
najnovejših, $23\,\%$ pri vozilih, starejših od $13$~let), s prevoženo
kilometrino in --- najmočneje --- z \emph{raznolikostjo ponudbe}. Pri modelih z
enotno ponudbo (raznolikost pod $0{,}1$) je MdAPE le $7{,}5\,\%$, pri razpršeni
ponudbi (nad $0{,}5$) pa $29{,}7\,\%$. To potrjuje, da je raznolikost, ki jo
sistem uporabniku tudi prikaže, dober vgrajen pokazatelj zanesljivosti
posamezne ocene.

\paragraph{Vpliv popusta.}
Slika~\ref{fig:popust} prikazuje predznačeno napako z in brez popusta. Brez
popusta je ocena rahlo optimistična (mediana $+3{,}8\,\%$ nad oglaševano ceno),
z namernim popustom pa se premakne na $-6{,}8\,\%$ pod oglaševano ceno. Ta
premik je skladen z zasnovo sistema, ki cilja pričakovano prodajno in ne
oglaševane cene; dokončno umerjanje popusta bi bilo mogoče šele s podatki o
dejanskih prodajah (razdelek~\ref{sec:groznje}).

\begin{figure}[htbp]
  \centering
  \includegraphics[width=0.7\textwidth]{accuracy_discount_bias.pdf}
  \caption{Predznačena napaka proti oglaševani ceni, z in brez namernega popusta.}
  \label{fig:popust}
  % vir: analysis/out/accuracy_discount_bias.pdf
\end{figure}

\paragraph{Časovni razrez.}
Pri napovedovanju najnovejših $10\,\%$ oglasov zgolj iz starejših $90\,\%$
znaša MdAPE predlaganega pristopa $13{,}0\,\%$ (tabela~\ref{tab:temporal}), kar
je celo malenkost bolje od rezultata leave-one-out in potrjuje, da ocena ne
izhaja iz uhajanja informacije.

\begin{table}[htbp]
  \centering
  \caption{Natančnost pri časovnem razrezu (napoved najnovejših $10\,\%$).
  MdAPE in MAPE sta v odstotkih, MAE in RMSE v evrih.}
  \label{tab:temporal}
  % vir: analysis/out/accuracy_temporal.csv
  \setlength{\tabcolsep}{5pt}
  \begin{tabular}{l r r r r}
    \toprule
    \textbf{Model} & \textbf{MdAPE} & \textbf{MAPE} & \textbf{MAE} & \textbf{RMSE} \\
    \midrule
    B0 globalna mediana       & 35{,}9 & 74{,}5 & 7596 & 11222 \\
    B1m mediana v oknu        & 12{,}8 & 19{,}7 & 3088 & 5805 \\
    B3 $k$-najbližjih sosedov & 10{,}5 & 18{,}2 & 2667 & 4890 \\
    B4 naš pristop            & 13{,}0 & 20{,}6 & 2994 & 5347 \\
    \bottomrule
  \end{tabular}
\end{table}

% =====================================================================
\section{Občutljivost algoritma na parametre}
\label{sec:obcutljivost}

Ocenjevalni algoritem vsebuje več hevristično določenih konstant: dovoljeno
odstopanje kilometrine ($\pm 40.000$~km) in letnika ($\pm 3$~leta), uteži
delnih ocen ($0{,}6$ in $0{,}4$), oblika padajoče funkcije in stopenjski
popust. V tem razdelku vsako spreminjamo ob nespremenjenih ostalih in merimo
MdAPE ter pokritost (vzorec $2000$ oglasov, leave-one-out). Rezultate izdela
skripta \texttt{analysis/sensitivity.py}.

\paragraph{Okno primerljivosti.}
Tabela~\ref{tab:sens-okno} in slika~\ref{fig:sens-km} kažeta, da je algoritem
na širino \emph{kilometrinskega} okna skoraj neobčutljiv: MdAPE se med
$\pm 10.000$~km in $\pm 80.000$~km giblje le med $13{,}2$ in $13{,}6\,\%$.
Nasprotno je odvisen od \emph{letnikovega} okna: pri $\pm 1$~letu je MdAPE
$12{,}6\,\%$, pri $\pm 5$~letih pa $15{,}2\,\%$. Ožje okno da bolj primerljiva
vozila, a manjši vzorec (nižja pokritost); gre za znano razmerje med
pristranskostjo in varianco~\cite{bias-variance}. Toplotna karta
(slika~\ref{fig:sens-heatmap}) potrjuje, da natančnost določa skoraj izključno
letnikovo okno. Sedanja izbira $\pm 3$~leta je razumen kompromis med
natančnostjo in pokritostjo; morebitna zožitev na $\pm 2$~leti bi MdAPE
znižala za približno $0{,}3$~odstotne točke ob še vedno visoki pokritosti.

\begin{table}[htbp]
  \centering
  \caption{Občutljivost na širino okna primerljivosti. Pokritost je delež
  poizvedb z vsaj $5$ primerljivimi vozili.}
  \label{tab:sens-okno}
  % vir: analysis/out/sensitivity_km.csv, sensitivity_year.csv
  \setlength{\tabcolsep}{5pt}
  \begin{tabular}{l r r r}
    \toprule
    \textbf{Parameter} & \textbf{MdAPE [\%]} & \textbf{Pokritost [\%]} & \textbf{Mediana $n$} \\
    \midrule
    kilometrina $\pm 10.000$~km & 13{,}2 & 92{,}3 & 53 \\
    kilometrina $\pm 40.000$~km & 13{,}4 & 98{,}1 & 196 \\
    kilometrina $\pm 80.000$~km & 13{,}6 & 99{,}0 & 362 \\
    \midrule
    letnik $\pm 1$~leto  & 12{,}6 & 95{,}6 & 97 \\
    letnik $\pm 3$~leta  & 13{,}4 & 98{,}1 & 196 \\
    letnik $\pm 5$~let   & 15{,}2 & 98{,}5 & 262 \\
    \bottomrule
  \end{tabular}
\end{table}

\begin{figure}[htbp]
  \centering
  \begin{subfigure}{0.48\textwidth}
    \includegraphics[width=\textwidth]{sensitivity_km.pdf}
    \caption{okno kilometrine}
  \end{subfigure}\hfill
  \begin{subfigure}{0.48\textwidth}
    \includegraphics[width=\textwidth]{sensitivity_year.pdf}
    \caption{okno letnika}
  \end{subfigure}
  \caption{Občutljivost na širino okna primerljivosti (MdAPE in pokritost).}
  \label{fig:sens-km}
  % vir: analysis/out/sensitivity_km.pdf, sensitivity_year.pdf
\end{figure}

\begin{figure}[htbp]
  \centering
  \includegraphics[width=0.7\textwidth]{sensitivity_km_year_heatmap.pdf}
  \caption{MdAPE [\%] po mreži dovoljenih odstopanj kilometrine in letnika.}
  \label{fig:sens-heatmap}
  % vir: analysis/out/sensitivity_km_year_heatmap.pdf
\end{figure}

\paragraph{Uteži, oblika jedra in popust.}
Tabela~\ref{tab:sens-ostalo} povzema preostale parametre. Razmerje uteži med
kilometrino in letnikom na natančnost praktično ne vpliva ($13{,}4$ do
$13{,}6\,\%$ za vsa preizkušena razmerja), prav tako ne zamenjava linearne
padajoče funkcije z eksponentno ($13{,}4$ proti $13{,}3\,\%$). Edina parametra,
ki opazno premakneta rezultat, sta torej širina letnikovega okna in namerni
popust: ta predznačeno napako proti oglaševani ceni premakne iz $+3{,}8\,\%$ v
$-6{,}8\,\%$, MdAPE pa zviša z $13{,}4$ na $14{,}5\,\%$. Ta premik je posledica
zavestne odločitve, da sistem ocenjuje pričakovano prodajno in ne oglaševane
cene, in ne slabše ocene.

\begin{table}[htbp]
  \centering
  \caption{Občutljivost na uteži, obliko jedra in popust.}
  \label{tab:sens-ostalo}
  % vir: analysis/out/sensitivity_weights.csv, sensitivity_kernel.csv, sensitivity_discount.csv
  \begin{tabularx}{\textwidth}{@{}l X r >{\raggedleft\arraybackslash}p{2.3cm}@{}}
    \toprule
    \textbf{Parameter} & \textbf{Nastavitev} & \textbf{MdAPE [\%]} & \textbf{Predznačena napaka [\%]} \\
    \midrule
    uteži (km / letnik) & $0{,}5 / 0{,}5$              & 13{,}4 & $+3{,}7$ \\
    uteži (km / letnik) & $0{,}6 / 0{,}4$ (privzeto)   & 13{,}4 & $+3{,}8$ \\
    uteži (km / letnik) & $0{,}8 / 0{,}2$              & 13{,}6 & $+3{,}8$ \\
    \midrule
    jedro               & linearno (privzeto)         & 13{,}4 & $+3{,}8$ \\
    jedro               & eksponentno                 & 13{,}3 & $+3{,}8$ \\
    \midrule
    popust              & brez                        & 13{,}4 & $+3{,}8$ \\
    popust              & stopenjski (privzeto)        & 14{,}5 & $-6{,}8$ \\
    \bottomrule
  \end{tabularx}
\end{table}

\paragraph{Sklep o robustnosti.}
Algoritem je robusten: na večino svojih konstant se skoraj ne odziva, kar
pomeni, da njihova hevristična izbira ni sporna. Njegova natančnost je odvisna
predvsem od kakovosti in gostote podatkov (razdelek~\ref{sec:natancnost}), ne
od finega nastavljanja parametrov.

% =====================================================================
\section{Latenca posameznega vrednotenja}
\label{sec:latenca}

Merimo latenco vrste (a) in (b) iz razdelka~\ref{sec:mere}. Meritev izvede
\texttt{benchmarks/bench\_algo.py}: $2000$ neposrednih klicev funkcije nad
naključnimi realnimi vhodi (brez HTTP), po $100$ ogrevalnih klicih, ki se ne
merijo. Ločeno merimo čas poizvedbe SQL, čas točkovanja v Pythonu in čas
serializacije odgovora.

\paragraph{Ozko grlo: pregled celotne tabele.}
Prvotna izvedba poizvedbe ni imela ustreznega indeksa.
Izsek~\ref{lst:explain-before} prikazuje izpis \texttt{EXPLAIN (ANALYZE,
BUFFERS)}: načrt je \texttt{Seq Scan}, baza pri vsakem vrednotenju prebere
celotno tabelo in zavrže $85.273$ vrstic, ki ne ustrezajo pogojem
(\texttt{Rows Removed by Filter}). Ker so podatki v celoti v predpomnilniku
(\texttt{shared hit}, brez \texttt{read}), je čas stabilen, a raste linearno z
velikostjo tabele.

\begin{center}
\captionof{listing}{\texttt{EXPLAIN (ANALYZE, BUFFERS)} pred dodajanjem indeksa.}
\label{lst:explain-before}
\begin{minted}{text}
Seq Scan on normalized_market_listing
    (cost=0.00..3314.15 rows=12 width=97)
    (actual time=0.555..6.847 rows=573 loops=1)
  Filter: ((kilometers >= 18500) AND (kilometers <= 98500)
       AND (registration_year >= 2019) AND (registration_year <= 2025)
       AND (brand_name = 'Citroën') AND (model_name = 'C3'))
  Rows Removed by Filter: 85273
  Buffers: shared hit=1168
Execution Time: 6.880 ms
\end{minted}
\end{center}

\paragraph{Učinek sestavljenega indeksa.}
Po dodajanju sestavljenega indeksa na
$(\text{brand\_name},\ \text{model\_name},\ \text{registration\_year},\
\text{kilometers})$ se načrt spremeni v \texttt{Bitmap Index Scan}
(izsek~\ref{lst:explain-after}): baza prebere le nekaj strani indeksa in
pripadajoče vrstice, čas izvajanja pa pade z $6{,}88$~ms na $0{,}69$~ms.
Latenca poizvedbe SQL se v mediani zmanjša z $9{,}7$~ms na $2{,}0$~ms
(tabela~\ref{tab:indeks-pred-po}, slika~\ref{fig:indeks-pred-po}). Pomembneje
od trenutnega pohitrenja je, da postane čas poizvedbe neodvisen od velikosti
tabele: namesto $O(\text{tabela})$ postane $O(\text{zadetki})$.

\begin{table}[htbp]
  \centering
  \caption{Latenca poizvedbe SQL pred in po dodajanju indeksa, v milisekundah.}
  \label{tab:indeks-pred-po}
  % vir: benchmarks/report.md (results_algo_before/after.json)
  \begin{tabular}{lrrrr}
    \toprule
    & \textbf{p50} & \textbf{p95} & \textbf{p99} & \textbf{Povprečje} \\
    \midrule
    pred (Seq Scan)  & 9{,}7 & 21{,}0 & 24{,}8 & 11{,}3 \\
    po (Index Scan)  & 2{,}0 & 8{,}6  & 12{,}5 & 3{,}0 \\
    \bottomrule
  \end{tabular}
\end{table}

\begin{figure}[htbp]
  \centering
  \includegraphics[width=0.6\textwidth]{latency_index_before_after.pdf}
  \caption{Latenca poizvedbe SQL pred in po dodajanju sestavljenega indeksa.}
  \label{fig:indeks-pred-po}
  % vir: analysis/out/latency_index_before_after.pdf
\end{figure}

\begin{center}
\captionof{listing}{\texttt{EXPLAIN (ANALYZE, BUFFERS)} po dodajanju indeksa.}
\label{lst:explain-after}
\begin{minted}{text}
Bitmap Heap Scan on normalized_market_listing
    (cost=5.84..50.49 rows=12 width=97)
    (actual time=0.202..0.615 rows=573 loops=1)
  Recheck Cond: (brand_name = 'Citroën' AND model_name = 'C3'
       AND registration_year BETWEEN 2019 AND 2025
       AND kilometers BETWEEN 18500 AND 98500)
  Buffers: shared hit=195
  ->  Bitmap Index Scan on idx_nml_valuation
        (actual time=0.166..0.166 rows=573 loops=1)
        Buffers: shared hit=6
Execution Time: 0.688 ms
\end{minted}
\end{center}

\paragraph{Razčlenitev latence po komponentah.}
Tabela~\ref{tab:latenca-vrednotenja} in slika~\ref{fig:latency-breakdown}
prikazujeta latenco celotnega vrednotenja po dodajanju indeksa. Tipično
vrednotenje traja $2{,}0$~ms, v najslabšem odstotku primerov (p99) pa
$13{,}5$~ms. Skoraj celotna latenca ($\approx 88\,\%$ povprečja) odpade na
poizvedbo SQL; točkovanje v Pythonu ($0{,}4$~ms povprečno) in serializacija
odgovora ($\approx 0{,}01$~ms) sta zanemarljiva. Poizvedba v povprečju vrne
$287$ primerljivih vozil (mediana $175$, največ $1495$).

\begin{table}[htbp]
  \centering
  \caption{Latenca posameznega vrednotenja po dodajanju indeksa
  ($N = 2000$), v milisekundah.}
  \label{tab:latenca-vrednotenja}
  % vir: benchmarks/report.md (results_algo.json)
  \begin{tabular}{lrrrr}
    \toprule
    \textbf{Komponenta} & \textbf{p50} & \textbf{p95} & \textbf{p99} & \textbf{Povprečje} \\
    \midrule
    Skupaj                & 2{,}0  & 9{,}7 & 13{,}5 & 3{,}1 \\
    Poizvedba SQL         & 1{,}8  & 8{,}4 & 11{,}5 & 2{,}7 \\
    Točkovanje (Python)   & 0{,}2  & 1{,}2 & 1{,}8  & 0{,}4 \\
    Serializacija (JSON)  & 0{,}01 & 0{,}02 & 0{,}03 & 0{,}01 \\
    \bottomrule
  \end{tabular}
\end{table}

\begin{figure}[htbp]
  \centering
  \includegraphics[width=0.6\textwidth]{latency_breakdown.pdf}
  \caption{Razčlenitev latence posameznega vrednotenja po komponentah.}
  \label{fig:latency-breakdown}
  % vir: analysis/out/latency_breakdown.pdf
\end{figure}

Vrednotenje se torej izvede daleč pod nefunkcionalno zahtevo $2$~s
(razdelek~\ref{sec:zahteve}); prostora za rast podatkovne baze je ob
prisotnosti indeksa dovolj.

% =====================================================================
\section{Obnašanje pod obremenitvijo}
\label{sec:obremenitev}

Merimo end-to-end latenco (vrsta (c)) in prepustnost končne točke
\texttt{POST /valuate} pri naraščajoči sočasnosti. Obremenitev ustvarja
skripta \texttt{benchmarks/loadtest.py} z več vzporednimi odjemalci, ki v
fiksnem $20$-sekundnem oknu neprekinjeno pošiljajo zahtevke z naključnimi
realnimi vhodi. Test se samodejno prekine ob deležu napak nad $5\,\%$ ali
$p_{95}$ nad $5$~s. Med testom vzorčimo tudi zasedenost baze
(\texttt{pg\_stat\_activity}).

\begin{table}[htbp]
  \centering
  \caption{Prepustnost in latenca pri naraščajoči sočasnosti.
  Latence so v milisekundah, $L$ je ocena po Littlovem zakonu.}
  \label{tab:obremenitev}
  % vir: benchmarks/results_load.json (loadtest.py --db-sample)
  \begin{tabular}{r r r r r r r}
    \toprule
    \textbf{Sočasnost} & \textbf{Prep.\ [z/s]} & \textbf{p50} & \textbf{p95} & \textbf{p99}
      & \textbf{Napake} & \textbf{$L$} \\
    \midrule
    1  & \res{x} & \res{x} & \res{x} & \res{x} & \res{x} & \res{x} \\
    5  & \res{x} & \res{x} & \res{x} & \res{x} & \res{x} & \res{x} \\
    10 & \res{x} & \res{x} & \res{x} & \res{x} & \res{x} & \res{x} \\
    20 & \res{x} & \res{x} & \res{x} & \res{x} & \res{x} & \res{x} \\
    40 & \res{x} & \res{x} & \res{x} & \res{x} & \res{x} & \res{x} \\
    \bottomrule
  \end{tabular}
\end{table}

\paragraph{Rezultat.}
\res{Prepustnost narašča približno linearno do sočasnosti X, nato se ustali pri
približno X zahtevkih na sekundo. Vse latence ostanejo pod nefunkcionalno
zahtevo 2~s.}

\paragraph{Skladnost meritve (Littlov zakon).}
\res{Stolpec $L$ (prepustnost krat povprečna latenca) ostaja blizu nastavljene
sočasnosti, kar potrjuje notranjo skladnost meritve.}

\paragraph{Dokaz ozkega grla.}
Teoretični strop postavitve je $1000 / S \approx 324$ vrednotenj na sekundo,
kjer je $S = 3{,}1$~ms povprečni čas ene evalvacije. Dejanska meja je nižja,
ker ves promet serializira \emph{ena sama deljena povezava do baze}
(\texttt{app/database.py}) in en sam strežniški proces.
\res{To potrjuje vzorčenje baze med testom: število aktivnih poizvedb kljub
naraščajoči sočasnosti odjemalcev ne preseže X.}

% =====================================================================
\section{Preverjanje funkcionalnih zahtev}
\label{sec:funkcionalne}

Tabela~\ref{tab:sledljivost} povezuje vsako funkcionalno zahtevo iz
razdelka~\ref{sec:zahteve} z načinom preverjanja in izidom.

\begin{table}[htbp]
  \centering
  \caption{Sledljivost funkcionalnih zahtev do preverjanja.}
  \label{tab:sledljivost}
  % vir: ročni scenariji + imena Jest testov v car-stock-market-be
  \begin{tabularx}{\textwidth}{@{}X X c@{}}
    \toprule
    \textbf{Zahteva} & \textbf{Način preverjanja} & \textbf{Izid} \\
    \midrule
    Registracija uporabnika              & \res{test / scenarij} & \res{OK} \\
    Prijava in odjava                    & \res{test / scenarij} & \res{OK} \\
    Objava oglasa (le registrirani)      & \res{test / scenarij} & \res{OK} \\
    Vrednotenje vozila po atributih      & \res{test / scenarij} & \res{OK} \\
    Všečkanje oglasov                    & \res{test / scenarij} & \res{OK} \\
    Ažuriranje lastnih/všečkanih oglasov & \res{test / scenarij} & \res{OK} \\
    Pregled oglasov (anonimni)           & \res{test / scenarij} & \res{OK} \\
    \bottomrule
  \end{tabularx}
\end{table}

\paragraph{Pokritost testov.}
Avtomatizirani testi zalednega sistema (Jest + Supertest,
razdelek~\ref{sec:testiranje-be}) obsegajo \res{X} testov v \res{X} zbirkah in
pokrivajo \res{X}\,\% stavkov ter \res{X}\,\% vej. Testi se izvedejo ob vsakem
zahtevku za združitev na cevovodu neprekinjene integracije
(razdelek~\ref{subsec:ci}).

% =====================================================================
\section{Preverjanje nefunkcionalnih zahtev}
\label{sec:nefunkcionalne}

Tabela~\ref{tab:nfr} za vsako nefunkcionalno zahtevo iz
razdelka~\ref{sec:zahteve} navaja merilo, izmerjeno vrednost in sodbo.

\begin{table}[htbp]
  \centering
  \caption{Preverjanje nefunkcionalnih zahtev.}
  \label{tab:nfr}
  \setlength{\tabcolsep}{4pt}
  \begin{tabularx}{\textwidth}{@{}>{\raggedright\arraybackslash}X
                                 >{\raggedright\arraybackslash}X
                                 >{\raggedright\arraybackslash}X c@{}}
    \toprule
    \textbf{Zahteva} & \textbf{Merilo} & \textbf{Izmerjeno} & \textbf{Sodba} \\
    \midrule
    Odziv vmesnika na dejanje uporabnika & $< 100$~ms & \res{x}~ms  & \res{OK} \\
    Zahtevki HTTP (razen vrednotenja)    & $< 300$~ms & \res{x}~ms ($p_{95}$) & \res{OK} \\
    Vrednotenje vozila                   & $< 2$~s    & $13{,}5$~ms ($p_{99}$, brez omrežja) & OK \\
    Vizualna privlačnost / dostopnost    & Lighthouse & \res{x} / \res{x} & \res{OK} \\
    Prilagodljiva postavitev             & preizkus na več velikostih zaslona & \res{da} & \res{OK} \\
    Varovanje uporabniških podatkov      & bcrypt, JWT z rokom & \res{opis} & \res{OK} \\
    Odzivnost pri zahtevnejših opravilih & asinhrona obdelava & \res{opis} & \res{OK} \\
    \bottomrule
  \end{tabularx}
\end{table}

\paragraph{Meritev odzivnosti vmesnika in endpointov.}
\res{Odzivnost uporabniškega vmesnika smo izmerili z orodjem Lighthouse na treh
ključnih straneh; ocene za performanco so X, za dostopnost pa X. Latenco
glavnih endpointov zalednega sistema smo izmerili z orodjem autocannon:
$p_{95}$ je v vseh primerih pod X~ms, torej pod zahtevo 300~ms.}

\paragraph{Varnost.}
Gesla so pred shranjevanjem zgoščena z algoritmom bcrypt
(razdelek~\ref{sec:avtentikacija}), žetoni JWT imajo rok veljavnosti $7$~dni,
poti vmesnika API pa ščiti sredinski preverjevalnik žetonov. \res{Orodje
\texttt{npm audit} ob zadnjem pregledu ni javilo kritičnih ranljivosti v
odvisnostih.}

% =====================================================================
\section{Uporabniški preizkus}
\label{sec:uporabniki}

Izvedli smo formativni uporabniški preizkus~\cite{nielsen1994} z namenom
odkriti težave uporabnosti, ne pa statistično ovrednotiti izdelka.

\paragraph{Udeleženci in postopek.}
V preizkus smo vključili \res{N} udeležencev iz kroga družine in bližnjih
prijateljev, katerih starost in računalniška pismenost sta se med seboj
razlikovali (tabela~\ref{tab:udelezenci}). Vsak udeleženec je prejel enak
seznam nalog: registracija, objava oglasa, všečkanje oglasa, izbira
slovenskega jezika v nastavitvah in vrednotenje poljubnega vozila. Med
izvajanjem je udeleženec razmišljal na glas, opazovalec pa je beležil, ali je
nalogo dokončal samostojno, in vsako opaženo težavo.

\begin{table}[htbp]
  \centering
  \caption{Udeleženci uporabniškega preizkusa.}
  \label{tab:udelezenci}
  \begin{tabular}{c c c l}
    \toprule
    \textbf{\#} & \textbf{Starost (razpon)} & \textbf{Rač.\ pismenost (1--5)} & \textbf{Naprava} \\
    \midrule
    1 & \res{x} & \res{x} & \res{x} \\
    2 & \res{x} & \res{x} & \res{x} \\
    3 & \res{x} & \res{x} & \res{x} \\
    \res{\dots} & & & \\
    \bottomrule
  \end{tabular}
\end{table}

\paragraph{Rezultati.}
Vsi udeleženci so vse naloge dokončali brez pomoči, kar potrjuje, da so glavne
poti v vmesniku dovolj jasne (RV4). Kljub temu so odkrili nekaj
pomanjkljivosti; tabela~\ref{tab:tezave} jih razvršča po resnosti (lestvica
$0$--$4$ po Nielsenu~\cite{nielsen1994}) in navaja, ali so bile odpravljene.

\begin{table}[htbp]
  \centering
  \caption{Težave, odkrite med uporabniškim preizkusom
  (resnost po lestvici $0$--$4$).}
  \label{tab:tezave}
  \setlength{\tabcolsep}{5pt}
  \begin{tabularx}{\textwidth}{@{}>{\raggedright\arraybackslash}X c c@{}}
    \toprule
    \textbf{Težava} & \textbf{Resnost} & \textbf{Status} \\
    \midrule
    Premajhni naslovi podstrani v glavi glavne strani       & \res{x} & odpravljeno \\
    Manjkajoče polje za ponovni vnos gesla pri registraciji & \res{x} & odpravljeno \\
    Velika števila brez ločil za tisočice                   & \res{x} & odpravljeno \\
    Manjkajoč opis pod slikami oblik vozil v obrazcu        & \res{x} & odpravljeno \\
    Premalo vozil, dostopnih za vrednotenje                 & \res{x} & \res{v načrtu} \\
    \bottomrule
  \end{tabularx}
\end{table}

Omejitve preizkusa (majhen in priložnostni vzorec, brez standardiziranega
vprašalnika, brez merjenja časa po nalogah) obravnavamo v
razdelku~\ref{sec:groznje}.

% =====================================================================
\section{Simulacija uporabe}
\label{sec:simulacija}

Za ponazoritev delovanja rešitve povzemamo tipičen potek uporabe. Uporabnik je
ob vstopu usmerjen na glavno stran (slika~\ref{fig:glavna-stran}), kjer sta
izpostavljeni prijava in registracija ter pregled novih in vseh objavljenih
oglasov. Registriran uporabnik lahko objavi oglas prek vodene forme
(slika~\ref{fig:objava-oglasa}), ovrednoti poljubno vozilo, všečka oglase in si
ogleda svoj profil z všečkanimi vozili (slika~\ref{fig:liked-cars}).

\begin{figure}[H]
  \centering
  
\includegraphics[width=0.9\textwidth]{glavna-stran.png}
  \caption{Glavna stran prototipa}
  \label{fig:glavna-stran}
\end{figure}
\begin{figure}[H]
  \centering
  \includegraphics[width=0.9\textwidth]{objava-oglasa.png}
  \caption{Vodena forma za objavo oglasa}
  \label{fig:objava-oglasa}
\end{figure}
\begin{figure}[H]
  \centering
  \includegraphics[width=0.9\textwidth]{liked-cars.png}
  \caption{Stran uporabnika s prikazanimi všečkanimi vozili}
  \label{fig:liked-cars}
\end{figure}

\paragraph{Primerjava ocen z obstoječimi portali.}
Za ilustracijo smo delovanje sistema primerjali s ponudbo na dveh spletnih
portalih na treh vozilih (tabela~\ref{tab:vrednotenje}). Sistematično in
kvantitativno oceno natančnosti podaja razdelek~\ref{sec:natancnost}; ta
primerjava je zgolj kvalitativna ponazoritev.

\begin{table}[htbp]
  \centering
  \caption{Primerjava ocenjenih cen s ponudbo na spletnih portalih (v evrih).}
  \label{tab:vrednotenje}
  \setlength{\tabcolsep}{5pt}
  \begin{tabular}{lccc}
    \toprule
    & \textbf{Audi A3} & \textbf{BMW ser.\ 2} & \textbf{Citroën Berlingo} \\
    \midrule
    slov.\ portal (razpon)     & 4.000--6.000  & 9.000--12.000  & 7.500--10.000 \\
    nem.\ portal (razpon)      & 4.500--9.000  & pribl.\ 11.100 & 8.000--12.000 \\
    \midrule
    Sistem --- dobra cena      & 4.500  & 7.600   & 7.600  \\
    Sistem --- sredinska cena  & 6.350  & 9.300   & 9.941  \\
    Sistem --- slaba cena      & 7.600  & 10.500  & 12.100 \\
    Raznolikost                & $49\,\%$ & $31\,\%$ & $45\,\%$ \\
    \bottomrule
  \end{tabular}
\end{table}

Pri vozilih z nizko raznolikostjo ponudbe (Audi~A3, BMW serije~2) se sredinska
ocena dobro ujema s cenami na obeh portalih. Pri modelu Citroën Berlingo je
raznolikost bistveno višja, zato je razpon med dobro in slabo oceno širši,
sredinska ocena pa manj zanesljiva --- kar je skladno z ugotovitvijo iz
razdelka~\ref{sec:natancnost}, da napaka strmo narašča z raznolikostjo ponudbe.
Pri interpretaciji je treba upoštevati, da sistem ocenjuje pričakovano prodajno
in ne oglaševano ceno.

% =====================================================================
\section{Grožnje veljavnosti}
\label{sec:groznje}

\paragraph{Notranja veljavnost.}
Kot merilo resnice pri merjenju natančnosti uporabljamo \emph{oglaševano} ceno,
ne dejanske prodajne cene, ki ni javno dostopna. Namerni popust, s katerim
sistem cilja pričakovano prodajno ceno, je zato ocenjen in ne umerjen na
dejanskih prodajah. Isti fizični oglas se lahko pojavi pri več virih; prečno
preverjanje izloči zapise z istim izvornim oglasom, ne more pa zaznati
podvojitev med viri.

\paragraph{Zunanja veljavnost.}
Podatki pokrivajo $10$ znamk in pretežno slovenski trg (z manjšim nemškim
vzorcem); znamka z največ oglasi na trgu (okoli $8000$ oglasov) v šifrantu ni
zajeta in je iz analize izpuščena. Uporabniški preizkus je temeljil na majhnem
in priložnostnem vzorcu.

\paragraph{Konstruktna veljavnost.}
Pojem \enquote{tržna vrednost} operacionaliziramo kot uteženo povprečje
oglaševanih cen primerljivih vozil, znižano za pričakovani pogajalski popust;
\enquote{raznolikost ponudbe} pa kot relativno širino interkvartilnega razpona
teh cen. Drugačna operacionalizacija bi lahko dala drugačne rezultate.

\paragraph{Zanesljivost.}
Meritve latence in obremenitve so bile izvedene na deljenem produkcijskem
strežniku, kjer lahko na rezultate vpliva ozadje. Meritve natančnosti so
ponovljive nad fiksnim posnetkom baze.

% =====================================================================
\section{Povzetek ugotovitev}
\label{sec:povzetek-analize}

Tabela~\ref{tab:rv-povzetek} povzema odgovore na raziskovalna vprašanja glede
na vnaprej postavljena merila.

\begin{table}[htbp]
  \centering
  \caption{Povzetek odgovorov na raziskovalna vprašanja.}
  \label{tab:rv-povzetek}
  \begin{tabularx}{\textwidth}{@{}l X c@{}}
    \toprule
    \textbf{RV} & \textbf{Glavni rezultat} & \textbf{Sodba} \\
    \midrule
    RV1 natančnost & MdAPE $13{,}5\,\%$ ($9$--$10\,\%$ za dobro pokrita vozila),
                     bistveno bolje od naivne ocene ($37{,}8\,\%$), primerljivo z
                     drugimi metodami & izpolnjeno \\
    RV2 odvisnost  & napaka pada z velikostjo vzorca (ustali se pri $n\approx20$),
                     narašča s starostjo, kilometrino in raznolikostjo ponudbe
                     ($7{,}5\,\%$ do $29{,}7\,\%$) & pojasnjeno \\
    RV3 hitrost    & vrednotenje $2$~ms (p50), $13{,}5$~ms (p99); indeks odpravi
                     odvisnost od velikosti baze; \res{meja X~z/s} & \res{izpolnjeno} \\
    RV4 uporabnost & vse naloge dokončane brez pomoči; \res{X} težav odpravljenih & izpolnjeno \\
    \bottomrule
  \end{tabularx}
\end{table}

\paragraph{Sodba o zastavljenih izhodiščih.}
Rezultati \emph{potrjujejo} izhodišča, ki smo si jih zastavili na začetku. Sistem na podlagi štirih
atributov vozila vrne oceno vrednosti z mediano odstotne napake $13{,}5\,\%$
(brez namernega popusta), pri dobro pokritih vozilih pa $9$--$10\,\%$, kar je
znotraj zastavljenega praga $15\,\%$ in bistveno bolje od naivne metode.
Predlagani pristop se med primerljivo natančnimi metodami ne uvršča na prvo
mesto --- preprost $k$-najbližjih sosedov doseže $11{,}2\,\%$ --- je pa edini,
ki poleg točkovne ocene vrne tudi razložljiv cenovni razpon in indeks
zaupanja, ta pa se izkaže za dober pokazatelj dejanske zanesljivosti
posamezne ocene. Glavne omejitve so odvisnost od enega prevladujočega
podatkovnega vira, nižja natančnost pri starejših vozilih in modelih z
razpršeno ponudbo ter merilo resnice, ki temelji na oglaševanih in ne
prodajnih cenah. Konkretne smeri za izboljšave --- vključno s prehodom na
metode strojnega učenja, ki jih ta analiza dodatno utemeljuje --- navajamo v
poglavju~\ref{ch:sklep}.


\chapter{Sklep}
\label{ch:sklep}
Cilj diplomskega dela je bil na podlagi jassno opredeljenih zahtev razviti prototip, ki uporabnikom omogoča objavo in
pregled oglasov rabljenih vozil, ter vanjo vgraditi sistem za samodejno ocenjevanje
tržne vrednosti vozila na podlagi primerljivih vozil na trgu. Oba cilja sta bila dosežena:
delujoč prototip z ločeno odjemalsko in zaledno komponento ter komponento za
vrednotenje vozil (poglavje~\ref{ch:struktura}), zasnovali pa smo tudi "cevovod" za zajem,
normalizacijo in primerjavo tržnih podatkov po
pristopu primerljivih prodaj z uteženim povprečjem (razdelka~\ref{subsec:filtriranje}
in~\ref{subsec:uteceno-povprecje}). Prototip je nameščen v produkcijskem okolju in
podpira vse zastavljene funkcionalne zahteve, od registracije, prijave ter objave in
pregleda oglasov do všečkanja in vrednotenja poljubnega vozila; izpolnjene so tudi
nefunkcionalne zahteve glede odzivnosti vmesnika, varovanja uporabniških podatkov s
zgoščevanjem gesel in časovno omejenimi žetoni ter razširljivosti arhitekture z ločenimi
komponentami in samodejnimi testi. Podatkovna baza trenutno obsega približno 86.000
normaliziranih oglasov za deset znamk vozil.\\
Na začetku smo si zastavili, da je mogoče uporabniku na podlagi zbranih tržnih
podatkov in primerjave s podobnimi vozili ponuditi uporabno in dovolj natančno oceno tržne
vrednosti. Sistem na podlagi znamke, modela, letnika in prevoženih
kilometrov vrne oceno cene skupaj z indeksom zaupanja, ki uporabnika pregledno opozori na
zanesljivost posamezne ocene glede na velikost vzorca primerljivih vozil.
Sistematično prečno preverjanje (razdelek~\ref{sec:natancnost}) je pokazalo mediano
absolutne odstotne napake $13{,}5\,\%$, pri dobro pokritih vozilih pa $9$--$10\,\%$;
ocena je bistveno natančnejša od naivne referenčne metode in znotraj zastavljenega
praga $15\,\%$. Napaka strmo narašča z raznolikostjo ponudbe (od $7{,}5\,\%$ pri enotni
do $29{,}7\,\%$ pri razpršeni ponudbi), zato indeks zaupanja in prikazani cenovni razpon
uporabnika pregledno opozorita na manj zanesljive ocene. Preizkus z uporabniki je potrdil
razumljivost in enostavnost uporabniškega vmesnika. Glavne omejitve rešitve so odvisnost
od enega prevladujočega podatkovnega vira, nižja zanesljivost ocene pri starejših vozilih
in modelih z razpršeno ponudbo ter omejen obseg preizkusa z uporabniki.\\
\begin{itemize}
    \item \textbf{Razširitev nabora vhodnih spremenljivk.} Ko bo na voljo dovolj
    normaliziranih podatkov, bi bilo smiselno v algoritem vključiti dodatne parametre, (kot
    so oprema, vrsta goriva, tip menjalnika in stanje vozila, kar bi izboljšalo natančnost
    ocene). V prihodnjem se bo z velikostjo uporabnikov, večal tudi moj dostop do podatkov in ne
    bomo rabili biti odvisni od ostalih virov podatkov.
    \item \textbf{Dinamično prilagajanje intervalov primerljivosti.} Namesto fiksnih mej bi
    lahko intervale kilometrine in letnika prilagajali sproti glede na gostoto podatkov v
    posameznem segmentu (na primer širši interval pri redkih modelih in ožji pri pogostih),
    s čimer bi omilili razmerje med pristranskostjo in varianco.
    \item \textbf{Izboljšanje kakovosti podatkov.} Analiza (razdelek~\ref{sec:podatki}) je
    pokazala, da približno $11\,\%$ normaliziranih oglasov vsebuje očitno napačne vrednosti,
    zato bi bilo smiselno v cevovod dodati samodejno filtriranje osamelcev (npr.\ po
    interkvartilnem razponu). Dopolniti je treba tudi šifrant znamk in modelov (manjkajoča
    prevladujoča znamka, združevanje modelov pri nekaterih znamkah).
    \item \textbf{Uvedba metod strojnega učenja.} Analiza natančnosti
    (razdelek~\ref{sec:natancnost}) je pokazala, da že preprosta metoda $k$-najbližjih
    sosedov preseže sedanji pristop primerljivih prodaj, kar dodatno utemeljuje ta korak.
    Ko bo obseg lastnih, normaliziranih podatkov zadosten, bi lahko sedanji, razložljivi
    pristop dopolnili ali nadomestili z zmogljivejšimi modeli (na primer regresijskimi
    modeli, naključnimi gozdovi ali nevronskimi mrežami), ki bi lahko predstavili bolj
    zahtevne, nelinearne odvisnosti med lastnostmi vozila in ceno~\cite{gegic2019car}.
    \item \textbf{Zmanjšanje odvisnosti od zunanjih virov.} Z rastjo števila uporabnikov in
    oglasov, na platformi, bomo lahko delež zunanjih podatkov
    postopno zmanjševali in ocene vse bolj gradili na podatkih, ki so zbrani znotraj same
    aplikacije. To bo izboljšalo tako natančnost kot vzdržnost sistema, saj se s tem
    zmanjša odvisnost od krhkih postopkov zajema s spletnih strani.
    \item \textbf{Postopna migracija zalednega sistema na TypeScript.} Uvedba statičnega
    preverjanja tipov bi dolgoročno izboljšala zanesljivost in vzdržljivost zalednega dela
    aplikacije~\cite{gao2017typescript}.
    \item \textbf{Nadgradnja infrastrukture in spremljanja delovanja.} Smiselna nadaljnja
    koraka sta uvedba centraliziranega beleženja dogodkov in spremljanja delovanja sistema
    (angl.\ \emph{monitoring}) v produkcijskem okolju. To bi olajšalo odkrivanje napak in
    analizo obremenitve sistema, ter razširitev nabora avtomatiziranih testov na
    del aplikacije za vrednotenje vozil in odjemalski del, ki v tem trenutku nista pokrita enako
    temeljito kot zaledni sistem, za katerega teste ob vsakem zahtevku za združitev že
    izvaja orodje GitHub Actions.
\end{itemize}


%%%%%%%%%%%%%%%%%%%%%%%%%%%%%%%%%%%%%%%%
%  LITERATURA
%%%%%%%%%%%%%%%%%%%%%%%%%%%%%%%%%%%%%%%%

%\cleardoublepage
\chapter*{Literatura}
\addcontentsline{toc}{chapter}{Literatura}

\defbibfilter{ostalo}{
  not type=article
  and not type=inproceedings
}
\printbibliography[
  type=article,
  title={Članki v revijah},
  heading=subbibintoc
]
\newpage
\printbibliography[
  type=inproceedings,
  title={Članki v zbornikih konferenc},
  heading=subbibintoc
]
\newpage
\printbibliography[
  filter=ostalo,
  title={Ostala literatura},
  heading=subbibintoc
]

\chapter*{Dodatek}
\addcontentsline{toc}{chapter}{Dodatek}

\section*{Koda: Definicija modela User}

\begin{center}
\captionof{listing}{Definicija modela \texttt{User} z ORM knjižnico Sequelize.}
\label{code:user-model}
\begin{minted}{javascript}
import { DataTypes } from "sequelize";
import sequelize from "#/config/db.js";

const User = sequelize.define(
  "User",
  {
    id: {
      type: DataTypes.INTEGER,
      autoIncrement: true,
      primaryKey: true,
    },
    firstName: {
      type: DataTypes.STRING,
      allowNull: false,
    },
    lastName: {
      type: DataTypes.STRING,
    },
    email: {
      type: DataTypes.STRING,
      allowNull: false,
      unique: true,
      validate: {
        isEmail: true,
      },
    },
    telephone: {
      type: DataTypes.STRING,
      allowNull: false,
    },
    password: {
      type: DataTypes.STRING,
      allowNull: false,
    },
    soldCars: {
      type: DataTypes.INTEGER,
      defaultValue: 0,
    },
  },
  { timestamps: true },
);

export default User;
\end{minted}
\end{center}

\begin{center}
\captionof{listing}{Primer normalizacije podatkov}
\label{code:normalization-python}
\begin{minted}{python}
def extract_data_from_raw_payload(raw
    valid = True
    data = {}

    raw_brand = raw_payload.get("brand", None)
    raw_model = raw_payload.get("model", None)

    if not raw_brand:
        valid = False
    if not raw_model:
        valid = False

    if valid:
        brand_match = match_brand(_preprocess(raw_brand), SOURCE)
        if brand_match is None:
            print(f"Could not match brand '{raw_brand}' for raw_id: {id}")
            valid = False
        else:
            brand_id, brand_canonical = brand_match
            data["brand_name"] = bran

            model_match = match_model
                _preprocess(raw_model), brand_id, brand_canonical, SOURCE
            )
            if model_match is None:
                print(
                    f"Could not match model '{raw_model}' (brand: '{brand_canonical}') for raw_id: {id}"
                )
                valid = False
            else:
                _, model_canonical = model_match
                data["model_name"] =
\end{minted}
\end{center}

\end{document}
